\documentclass[a4paper,12pt]{article}

\usepackage{preamble}



\pgfmathsetmacro\sprayRadius{.4pt}
\pgfmathsetmacro\sprayPeriod{.6cm}

\pgfdeclarepatternformonly{spray}{\pgfpoint{-\sprayRadius}{-\sprayRadius}}{\pgfpoint{1cm + \sprayRadius}{1cm + \sprayRadius}}{\pgfpoint{\sprayPeriod}{\sprayPeriod}}{
    \foreach \x/\y in {2/53,6/52,11/48,23/49,20/47,32/46,41/47,47/51,56/52,46/44,4/43,16/42,33/41,41/37,49/35,55/31,37/35,44/30,28/37,24/36,17/37,7/38,0/31,8/29,18/31,28/30,37/28,30/27,46/24,51/21,24/23,12/24,4/21,18/19,12/16,31/21,38/18,26/16,46/16,56/12,52/10,45/8,51/4,37/12,35/7,24/9,14/9,2/12,8/6,15/4,27/0,34/1,40/1} {
    \pgfpathcircle{\pgfpoint{(\x + random()) / 57 * \sprayPeriod}{\sprayPeriod - (\y + random()) / 55 * \sprayPeriod}}{\sprayRadius}
    }
\pgfusepath{fill}
}

\begin{titlepage}
\title{Treatise on Differential Forms}
\author{Abdul Musthakin}
\date{February 2026}
\end{titlepage}

\begin{document}
\maketitle

\pagebreak

\tableofcontents

\pagebreak

\section{Introduction}\label{sec:introduction}

\subsection{Aims and Prerequisites}\label{subsec:introduction:aims_and_prerequisites}

This text is an account of differential forms and de Rham cohomology.
Its aim, in this version, is to cover Chapters I and II of Bott and Tu \cite{bott1982differential} in a self-contained way, starting from multilinear algebra and smooth manifolds.
Later versions will follow the later chapters of the same book.

The reader is assumed to be familiar with linear algebra and multivariable calculus, and to be comfortable reading proofs.
No knowledge of tensors or manifolds is assumed.
The point-set topology that we use is collected in \cref{sec:topological_spaces}.
Results from point-set topology are cited rather than proved, as are some results from the theory of smooth manifolds that are not about differential forms.

\subsection{Notation and Conventions}\label{subsec:introduction:notation_and_conventions}

In general, vectors are not bolded, though they may be when they are Euclidean.

$0 \in \N$.

$\R^n$ refers to $n$-dimensional Euclidean space.
We write $\R$ for the space $\R^1$ as well as for the set of real numbers, though we do not do this for other spaces such as $\mathbb{S}^1$ or $\mathbb{T}^1$.

In differential geometry, the term \enquote*{function} is often used to specifically refer to functions with codomain $\R$ (or perhaps $\C$).
Then, other maps, which may still be actual functions, are referred to as \enquote*{maps}.
This convention will not be used in this text.

Standard tensor notation is used.

Einstein summation convention is used.

Typically:

\begin{table}[!htbp]
\begin{tabular}{@{}llc@{}}
\toprule
Object                          & Style                   & Examples                                    \\ \midrule
Variables                       & Lowercase, Latin        & $a, b, c$                                   \\
Standard Functions              & Lowercase, Latin        & $f, g, h$                                   \\
Scalars                         & Lowercase, Greek        & $\lambda, \mu, \alpha$                      \\
Sets                            & Uppercase, Latin        & $A, S, X$                                   \\
Special sets                    & Blackboard Bold         & $\mathbb{R}, \mathbb{S}^2, \mathbb{C}$      \\
Matrices                        & Uppercase, Latin, Bold  & $\mathbf{M}, \mathbf{A}, \mathbf{B}$        \\
Spaces \& Manifolds             & Uppercase, Latin        & $X, M, V$                                   \\
Categories                      & Uppercase, Script       & $\mathscr{C}, \mathscr{A}, \mathscr{B}$     \\
Lie Algebras                    & Fraktur                 & $\mathfrak{g}, \mathfrak{h}, \mathfrak{so}$ \\
Groups \& Algebras              & Uppercase, Latin        & $G, A, B$                                   \\
Groupoids \& Lie Algebroids    & Uppercase, Calligraphic & $\mathcal{G}, \mathcal{A}, \mathcal{B}$     \\ \bottomrule
\end{tabular}
\end{table}

Throughout, $K$ denotes a field.
From \cref{subsec:multilinear_algebra:tensor_algebra} onwards, $K$ has characteristic zero, and from \cref{sec:manifolds} onwards the ground field is $\R$.
General commutative rings $R$ are used only for multilinear maps and tensor products in \cref{subsubsec:multilinear_algebra:tensors:tensor_products} and for (co)chain complexes in \cref{subsubsec:de_rham_theory:de_rham_cohomology:de_rham_complex}.

Wedge products of alternating tensors and of differential forms use the determinant convention
\[
    S \wedge T = \frac{(p+q)!}{p!\,q!} \alt(S \otimes T)
\]
for tensors of rank $p$ and $q$, so that $(\omega^1 \wedge \cdots \wedge \omega^k)(v_1, \ldots, v_k) = \det(\omega^i(v_j))$.
The wedge product on the quotient $\extp(V)$ of $T(V)$ is induced by $\otimes$ and carries no such factor, so the identification of $\extp^k(V)$ with $\Lambda^k(V)$ in \cref{prop:exterior_algebra_quotient_subspace_equivalence} rescales by $k!$ in degree $k$.

Each of the following objects has one symbol in this text, except that the pushforward may also be written as $f_*$.
Capital letters denote spaces and algebras, and lower-case words denote operators.
\begin{table}[!htbp]
\begin{tabular}{@{}ll@{}}
\toprule
Object                                                            & Symbol                                \\ \midrule
Tensor powers $V^{\otimes k}$ and $(V^*)^{\otimes k}$             & $T^k(V)$, $T_k(V)$                    \\
Type $(p, q)$ tensors                                             & $T^p_q(V)$                            \\
Staggered tensors of rank $k$                                     & $\mathcal{T}^k(V)$                    \\
Tensor algebra                                                    & $T(V)$                                \\
Symmetric and antisymmetric algebras                              & $\Sym(V)$, $\Asym(V)$                 \\
Symmetrization and antisymmetrization of a tensor                 & $\sym$, $\asym$                       \\
Alternating tensors (a subspace of $T^k(V)$)                      & $\Lambda^k(V)$, $\Lambda(V)$          \\
Alternation of a tensor                                           & $\alt$                                \\
Alternating $k$-multilinear maps $V^{\times k} \to W$             & $\Alt^k(V, W)$                        \\
Exterior algebra (a quotient of $T(V)$)                           & $\extp(V)$, $\extp^k(V)$              \\
Pushforward of a smooth map $f$, at a point and globally           & $Df_p$, $Df$                          \\
Exterior derivative, including the differential of a function     & $\d$                                  \\
Differential forms                                                & $\Omega^k(M)$                         \\ \bottomrule
\end{tabular}
\end{table}

\pagebreak

\section{Multilinear Algebra}\label{sec:multilinear_algebra}

\subsection{Tensors}\label{subsec:multilinear_algebra:tensors}

There are many different ways to motivate the discussion of tensors, and there are then many different ways to actually discuss them.
Tensors may be thought of as multidimensional arrays just as vectors and matrices may be thought of as one- and two-dimensional arrays, respectively.
They may also be thought of as multilinear maps.
Whilst we will view tensors in an even more abstract sense, we still need to define multilinear maps.
\begin{definition}[Multilinear Map]\label{def:multilinear_map}
Let $R$ be a commutative ring, and let $V_1, \dots, V_n$, and $W$ be $R$-modules.
A map $f \colon V_1 \times V_2 \times \dots \times V_n \to W$ is $n$-multilinear if it is linear in each of its arguments.
That is, for each $i \in \{1, \dots, n\}$, any elements $v_{(j)} \in V_j$ where $j \neq i$, we have
\begin{equation}
    f(v_{(1)}, \ldots, \lambda u + \mu w, \ldots, v_{(n)}) = \lambda f(v_{(1)}, \ldots, u, \ldots, v_{(n)}) + \mu f(v_{(1)}, \ldots, w, \ldots, v_{(n)})
\end{equation}
for any scalars $\lambda, \mu \in R$.
A two-multilinear map is called bilinear and a three-multilinear map is called trilinear.
\end{definition}

\subsubsection{Tensor Products}\label{subsubsec:multilinear_algebra:tensors:tensor_products}

The tensor product between $R$-modules can be interpreted as capturing the property of bilinearity, as it is used to uniquely factor bilinear maps into linear maps.
When we later see that tensors can really be considered bilinear (and more generally multilinear) maps themselves, we may consider them to be the most fundamental versions of such maps due to the aforementioned property, which is made precise via its universal property.
\begin{definition}[Tensor Product]\label{def:tensor_product}
    Let $R$ be a commutative ring, and let $V$ and $W$ be $R$-modules.
    The tensor product of $V$ and $W$ is the $R$-module $V \otimes_R W$ equipped with a bilinear map $\varphi \colon V \times W \to V \otimes_R W$ such that for any bilinear map $h \colon V \times W \to Z$, where $Z$ is an $R$-module, there is a unique linear map $\tilde{h} \colon V \otimes_R W \to Z$ making the following diagram commute.
    \[\begin{tikzcd}
        {V \times W} && {V \otimes W} \\
        \\
        && Z
        \arrow["\varphi", from=1-1, to=1-3]
        \arrow["h"', from=1-1, to=3-3]
        \arrow["{\tilde{h}}", dashed, from=1-3, to=3-3]
    \end{tikzcd}\]
    We call the map $\varphi$ the tensor product of the elements of $V$ and $W$, writing $\varphi(v, w) = v \otimes w$.
    A tensor is an element of the tensor product of $R$-modules.
    We may call the product $V \otimes_R W$ a tensor product space, and we often denote $\otimes_R$ as just $\otimes$.
\end{definition}
As the tensor product is defined using a universal property, it is unique up to a unique isomorphism.
Whilst we have defined the tensor product in the generality of $R$-modules, we will mainly focus on the case where we have $K$-vector spaces.
Vector spaces should be assumed to be over $K$ unless otherwise stated.
Furthermore, we will consider vector spaces to be finite-dimensional unless otherwise stated.

We will provide a construction of the tensor product of two vector spaces, so that we have a proof that they exist.
For this, consider two $K$-vector spaces $V$ and $W$, and consider the $K$-vector space $L$ whose basis is $V \times W$.
This space may be obtained by considering vector space of functions $V \times W \to K$ with finitely many nonzero values and identifying $(v, w)$ with the function that takes the value $1$ on $(v, w)$ and $0$ otherwise.
Define $S$ to be the subspace of $L$ generated by the vectors of the following forms, with $v, v' \in V$, $w, w' \in W$, and $\lambda \in K$.
\begin{align}
    \begin{gathered}
        (v, w + w') - (v, w) - (v, w') \\
        (v + v', w) - (v, w) - (v', w) \\
        (\lambda v, w) - \lambda (v, w) \\
        (v, \lambda w) - \lambda (v, w)
    \end{gathered}
\end{align}
Then, we define $V \otimes W = L/S$.
The image of $(v, w)$ is the equivalence class $[(v, w)]$ which is denoted $v \otimes w$.
With these conditions, we effectively set some elements of the tensor product space to be the same so that the tensor product is a bilinear map.
Indeed, this construction of the tensor product highlights the nature of tensors as multilinear maps.
It remains to be shown that the space defined here actually satisfies the universal property.
\begin{proposition}\label{prop:tensor_product_space_construction_satisfies_universal_property}
    Given two $K$-vector spaces $V$ and $W$, the tensor product $V \otimes W$ constructed as the aforementioned quotient space $L/S$ satisfies the universal property of the tensor product.
\end{proposition}
\begin{proof}
    The map $\varphi \colon V \times W \to V \otimes W$ is given by $\varphi(v, w) = [(v, w)] = v \otimes w$.
    We can verify that $\varphi$ is a bilinear map.
    Let, $v, v' \in V$, $w, w' \in W$, and let $\lambda \in K$.
    The subspace $S$ contains the element
    \begin{equation*}
        (v + v', w) - (v, w) - (v', w).
    \end{equation*}
    In the quotient space, this is equivalent to the zero vector, so
    \begin{equation*}
        [(v + v', w)] - [(v, w)] - [(v', w)] = 0.
    \end{equation*}
    Rearranging the equation and rewriting the terms gives us
    \begin{equation*}
        (v + v') \otimes w = v \otimes w + v' \otimes w.
    \end{equation*}
    Similarly, we get
    \begin{equation*}
        (\lambda v) \otimes w = \lambda (v \otimes w).
    \end{equation*}
    This shows that $\varphi$ is linear in its first argument, and it can be similarly shown that it is linear in its second argument, so it is a bilinear map.

    Let $Z$ be a $K$-vector space.
    Given a bilinear map $h \colon V \times W \to Z$, there exists a unique linear map $f \colon L \to Z$ such that $f(v, w) = h(v, w)$, extending $h$.
    We wish to define $\tilde{h}(v \otimes w) = f(v, w)$, but this must be well-defined.
    That is, given that $[x] = [y]$ for $x, y \in L$, we need $\tilde{h}([x]) = \tilde{h}([y])$.
    Equivalently, we need $f(x) = f(y)$, which is the same as $f(x - y) = 0$.
    If $[x] = [y]$, then $x - y \in S$, so it is sufficient for $S$ to be in the kernel of $f$.
    Applying $f$ to any of the four generators of $S$ yields zero, so $S \subseteq \ker f$.
    Hence, there exists a linear map $\tilde{h} \colon L/S \to Z$ such that $h = \tilde{h} \circ \varphi$.

    To show that $\tilde{h}$ is unique, suppose that there exists a linear map $\hat{h} \colon L/S \to Z$ such that $h = \hat{h} \circ \varphi$.
    We have $\hat{h}(v \otimes w) = h(v, w) = \tilde{h}(v \otimes w)$, so $\hat{h}$ and $\tilde{h}$ agree on the elements $v \otimes w$, which span $V \otimes W$.
    All other elements are linear combinations of the elements $v \otimes w$, and the two maps are linear.
    Therefore, $\hat{h} = \tilde{h}$.
\end{proof}

A tensor of the form $v \otimes w \in V \otimes W$ is called an \emph{elementary tensor} or a \emph{pure tensor} or a \emph{decomposable tensor}.
Not all elements of $V \otimes W$ may be expressed as such tensors, but they can all be expressed using them.
Particularly, elementary tensors formed by taking the tensor product of the basis vectors of $V$ and $W$ form a basis for $V \otimes W$.
Given a basis $\{e_i\}_{i \in I}$ for $V$ and a basis $\{f_j\}_{j \in J}$ for $W$, we have
\begin{equation}
    V \otimes W = \operatorname{span}\{e_i \otimes f_j : (i, j) \in I \times J\}.
\end{equation}
These elements are also linearly independent.
Let $e^k$ and $f^l$ be the coordinate functionals of the two bases.
The bilinear map $(v, w) \mapsto e^k(v) f^l(w)$ induces, by the universal property, a linear map $V \otimes W \to K$ which takes the value $\delta_{ik}\delta_{jl}$ on $e_i \otimes f_j$.
If $V$ is $n$ dimensional and $W$ is $m$ dimensional, it follows that $V \otimes W$ is $mn$ dimensional.
Note that, in general, an element of $V \otimes W$ is not an elementary tensor.
We can only combine and simplify tensors when we are allowed to via their properties of being a bilinear map.
For instance, given the tensor product space $\R^2 \otimes \R^2$, the tensor $e_1 \otimes e_1 + e_2 \otimes e_2$ cannot be further simplified (it is not elementary, as its matrix has rank two), but we have $e_1 \otimes e_1 + e_1 \otimes e_2 = e_1 \otimes (e_1 + e_2)$.

The tensor product of vector spaces is essentially associative and commutative.
That is, although the product does not strictly have these properties in the sense of equality, we have canonical isomorphisms instead.
By a slight abuse of notation, we can consider the tensor product to literally be associative and commutative, so we can make sense of expressions like $U \otimes V \otimes W$.
\begin{proposition}\label{prop:tensor_product_associative_commutative}
    The tensor product is associative and commutative up to canonical isomorphisms.
    Given $K$-vector spaces $U$, $V$, and $W$, we have
    \begin{equation}
        (U \otimes V) \otimes W \cong U \otimes (V \otimes W)
    \end{equation}
    via the canonical isomorphism that maps $(u \otimes v) \otimes w$ to $u \otimes (v \otimes w)$.
    We also have
    \begin{equation}
        V \otimes W \cong W \otimes V
    \end{equation}
    via the canonical isomorphism that maps $v \otimes w$ to $w \otimes v$.
\end{proposition}
\begin{proof}
    For some fixed $w \in W$, consider the map $f_w \colon U \times V \to U \otimes (V \otimes W)$ defined by $f_w(u, v) = u \otimes (v \otimes w)$.
    By the universal property, this induces a unique linear map $\varphi_w \colon U \otimes V \to U \otimes (V \otimes W)$.
    Now, consider the map $h \colon (U \otimes V) \times W \to U \otimes (V \otimes W)$ defined by $h(x, w) = \varphi_w(x)$.
    This is linear in $x$ as $\varphi_w$ is linear, and it is linear in $w$ by the uniqueness of the induced maps, as $\varphi_{\lambda w + \mu w'}$ and $\lambda \varphi_w + \mu \varphi_{w'}$ are both linear and agree with $f_{\lambda w + \mu w'}$ after composing with $\otimes$.
    Hence, by the universal property, there exists a unique linear map $\alpha \colon (U \otimes V) \otimes W \to U \otimes (V \otimes W)$ given by $\alpha((u \otimes v) \otimes w) = u \otimes (v \otimes w)$.
    \[\begin{tikzcd}
        {U \times V} && {U \otimes V} && {(U \otimes V) \otimes W} \\
        \\
        && {U \otimes (V \otimes W)} && {(U \otimes V) \times W}
        \arrow["\otimes", from=1-1, to=1-3]
        \arrow["{f_w}"', from=1-1, to=3-3]
        \arrow["{\varphi_w}", dashed, from=1-3, to=3-3]
        \arrow["\alpha"', dashed, from=1-5, to=3-3]
        \arrow["\otimes"', from=3-5, to=1-5]
        \arrow["h", from=3-5, to=3-3]
    \end{tikzcd}\]
    We can similarly construct its inverse, the linear map $\beta \colon U \otimes (V \otimes W) \to (U \otimes V) \otimes W$ given by $\beta(u \otimes (v \otimes w)) = (u \otimes v) \otimes w$.
    The compositions $\beta \circ \alpha$ and $\alpha \circ \beta$ are the identity maps on the generators of their respective spaces, they are the identity maps for the entire space.
    Additionally, they are independent of a choice of basis for the tensor product.
    Therefore, there exists a canonical isomorphism giving $(U \otimes V) \otimes W \cong U \otimes (V \otimes W)$.

    We will now show commutativity.
    Consider the function $f \colon V \times W \to W \otimes V$ given by $f(v, w) = w \otimes v$.
    It is bilinear in $(v, w)$ because $\otimes$ is bilinear in each argument.
    By the universal property, there exists a unique linear map $\Phi \colon V \otimes W \to W \otimes V$ given by $\Phi(v \otimes w) = w \otimes v$.
    \[\begin{tikzcd}
        {V \times W} && {V \otimes W} \\
        \\
        && {W \otimes V}
        \arrow["\otimes", from=1-1, to=1-3]
        \arrow["f"', from=1-1, to=3-3]
        \arrow["\Phi", dashed, from=1-3, to=3-3]
    \end{tikzcd}\]
    Again, we can find an inverse map $\Phi^{-1} \colon W \otimes V \to V \otimes W$ given by $\Phi^{-1}(w \otimes v) = v \otimes w$.
    Thus, $\Phi$ is an isomorphism, and since it is independent of any choice of basis, it is a canonical isomorphism.
    Therefore, we have $V \otimes W \cong W \otimes V$.
\end{proof}
\begin{remark}
    It is important to note the tensor product of vectors themselves is not, in general, commutative.
    Talking about associativity of the tensor product of vectors is not particularly meaningful, as $(u \otimes v) \otimes w$ and $u \otimes (v \otimes w)$ live in different spaces.
    However, we can compare the tensors $v_1 \otimes v_2$ and $v_2 \otimes v_1$, where $v_1, v_2 \in V$.
    These two tensors need not be equal.
\end{remark}
We can thus consider a space consisting of many different vector spaces connected via tensor products.
By associativity, we can unambiguously omit parentheses in the expressions for such spaces.
The universal property in \Cref{def:tensor_product} is extended by considering multilinear maps instead of bilinear maps, so we obtain the following commutative diagram characterizing the tensor product of $n$ vector spaces.
\[\begin{tikzcd}
        {V_1 \times \cdots \times V_n} && {V_1 \otimes \cdots \otimes V_n} \\
        \\
        && Z
        \arrow["\varphi", from=1-1, to=1-3]
        \arrow["h"', from=1-1, to=3-3]
        \arrow["{\tilde{h}}", dashed, from=1-3, to=3-3]
\end{tikzcd}\]

\subsubsection{Tensor Spaces}\label{subsubsec:multilinear_algebra:tensors:tensor_spaces}

With a satisfactory introduction to tensor products, we are able to expand on the concept of tensors themselves.
We have defined tensors as being elements of the tensor product of two vector spaces.
This will still be the case, but we usually consider that tensor product space as being made up of many vector spaces.
We may consider taking the tensor product of a vector space $V$ with itself many times, and we denote this using tensor powers.
From that, we may define tensor types.
\begin{definition}[Tensor Power]\label{def:tensor_power}
    Given a finite-dimensional $K$-vector space $V$, the $n$-th tensor power of $V$ is given by
    \begin{equation}
    V^{\otimes n} = \underbrace{V \otimes \cdots \otimes V}_{n}.
    \end{equation}
    The zeroth tensor power of $V$ is defined to be the base field $K$.
    The space $V^{\otimes k}$ is denoted by $T^k(V)$ and $(V^*)^{\otimes k}$ is denoted by $T_k(V)$; these are the cases $T^k_0(V)$ and $T^0_k(V)$ of \cref{def:tensor_space_and_type}.
\end{definition}
\begin{definition}[Tensor Space and Type]\label{def:tensor_space_and_type}
    Given a finite-dimensional $K$-vector space $V$, the type $(p, q)$ tensor space of $V$ is defined by
    \begin{equation}
        T^p_q(V) = V^{\otimes p} \otimes (V^*)^{\otimes q} = \underbrace{V \otimes \cdots \otimes V}_{p} \otimes \underbrace{V^* \otimes \cdots \otimes V^*}_{q}.
    \end{equation}
    An element of $T^p_q(V)$ is called a type $(p, q)$ tensor; we may shorten this to a $(p, q)$ tensor.
    The rank (also called order or degree) of such a tensor is equal to $p + q$.
    We define $T^0_0(V)$ to be $K$.
\end{definition}
From the universal property of the tensor product, we have $\Hom(V \otimes W, Z) \cong \Multi(V \times W, Z)$, with the latter space being the space of multilinear maps from $V \times W$ to $Z$.
Considering $Z$ to be the base field $K$, we have $(V \otimes W)^* \cong \Multi(V \times W, K)$.
Then, consider the natural isomorphisms $(V \otimes W)^* \cong V^* \otimes W^*$ and $V \cong V^{**}$.
These may be used to arrive at
\begin{align}
    \begin{split}
        T^p_q(V) &= V^{\otimes p} \otimes (V^*)^{\otimes q} \\
    &\cong (V^{**})^{\otimes p} \otimes (V^*)^{\otimes q} \\
    &\cong ((V^*)^{\otimes p} \otimes V^{\otimes q})^{*} \\ 
    &\cong \Multi((V^*)^{\times p} \times V^{\times q}, K).
    \end{split}
\end{align}
Hence, a type $(p, q)$ tensor is essentially a multilinear map of $p$ dual vectors and $q$ vectors.
The tensor itself is composed of $p$ vectors and $q$ dual vectors (slots of the tensor), with each such object taking its dual as input.
This justifies alternative definitions of tensors as literally being multilinear maps and lets us consider the objects in a less abstract sense.
Henceforth, we shall use both interpretations of tensors as is convenient.
By using vectors and dual vectors in our definition, we allow for type $(p, q)$ tensors to encompass many different mathematical objects.
For instance, a type $(0, 2)$ tensor $B$ is essentially a bilinear form $V \times V \to K$ via $(v, w) \mapsto B(v, w)$.

A $(0, 0)$ tensor is a scalar.
A $(1, 0)$ tensor is a vector, and a $(0, 1)$ tensor is a dual vector.
As we increase the first number in the pair, we get bivectors, trivectors, and -- in general -- multivectors.
A $(1, 1)$ tensor is a map $T \colon V^* \times V \to K$, but this is essentially a linear transformation $L \colon V \to V$ via the isomorphism $\Theta \colon \End(V) \to T^1_1(V)$ with $\Theta(A)(\omega, v) = \omega(A(v))$.\footnote{$\End(V)$ is the space of endomorphisms of $V$.}
The tensor can be viewed as having two slots, and filling the vector slot gives a map $V^* \to K$ which is a double dual vector.
This extends to general tensors $T \colon (V^*)^{\times p} \times V^{\times q} \to K$; they can be viewed as linear maps $L \colon V^{\otimes q} \to V^{\otimes p}$.

Given a basis $\{e_1, \ldots, e_n\}$ for $V$ and the dual basis $\{\varepsilon^1, \ldots, \varepsilon^n\}$ for $V^*$, there is an induced basis for the tensor space $T^p_q(V)$ that consists of elements of the form $e_{i_1} \otimes \cdots \otimes e_{i_p} \otimes \varepsilon^{j_1} \otimes \cdots \otimes \varepsilon^{j_q}$, with the indices $i_k$ and $j_m$ taking values in $\{1, \ldots, n\}$.
Such a basis tensor may be thought of as acting on $q$ vectors $v_{(1)}, \ldots, v_{(q)}$ and $p$ dual vectors $\omega^{(1)}, \ldots, \omega^{(p)}$ to give $\omega^{(1)}(e_{i_1}) \cdots \omega^{(p)}(e_{i_p}) \varepsilon^{j_1} (v_{(1)}) \cdots \varepsilon^{j_q}(v_{(q)})$.\footnote{We use parentheses to distinguish a listing of different objects from a listing of the components of the same object.}
With respect to this basis, we may write an arbitrary component of a type $(p, q)$ tensor as the coefficient $T^{i_1 \ldots i_p}_{j_1 \ldots j_q}$.
The following uses Einstein summation convention, summing over all indices.
\begin{equation}
    T = T^{i_1 \ldots i_p}_{j_1 \ldots j_q} e_{i_1} \otimes \cdots \otimes e_{i_p} \otimes \varepsilon^{j_1} \otimes \cdots \otimes \varepsilon^{j_q}
\end{equation}
Viewing $T$ as a multilinear map, the component $T^{i_1 \ldots i_p}_{j_1 \ldots j_q}$ is given by evaluating $T$ on the basis vectors and dual vectors.
\begin{equation}\label{eq:tensor_component}
    T^{i_1 \ldots i_p}_{j_1 \ldots j_q} = T(\varepsilon^{i_1}, \ldots, \varepsilon^{i_p}, e_{j_1}, \ldots, e_{j_q})
\end{equation}
The upper indices of a tensor component may be called contravariant indices, and the lower indices may be called covariant indices.
This refers to the way a tensor component transforms under a change of basis.
Consider a change from a basis $\{e_\mu\}$ to a basis $\{\bar{e}_\nu\}$ under which the components of a vector transform via the matrix $A^\nu_\mu$, so that $\bar{v}^\nu = A^\nu_\mu v^\mu$; denote its inverse by $W^\mu_\nu$.
Note that matrices represent linear transformations and thus can be viewed as type $(1, 1)$ tensors -- hence we wrote their indices like we do for tensors.
The basis vectors transform via the inverse, $\bar{e}_\nu = W^\mu_\nu e_\mu$, and the basis dual vectors transform via $A^\nu_\mu$, $\bar{\varepsilon}^\nu = A^\nu_\mu \varepsilon^\mu$.
By the multilinearity of tensors and \eqref{eq:tensor_component}, the new tensor component is given by
\begin{equation}
    \overbar{T}^{\alpha_1 \ldots \alpha_p}_{\beta_1 \ldots \beta_q}
    = A^{\alpha_1}_{\mu_1} \cdots A^{\alpha_p}_{\mu_p}
      W^{\nu_1}_{\beta_1} \cdots W^{\nu_q}_{\beta_q}
      \, T^{\mu_1 \ldots \mu_p}_{\nu_1 \ldots \nu_q}.
\end{equation}
It can be seen that covariance is characterized by transforming with the basis, whilst contravariance is characterized by transforming opposite to it.
If the transformation is of coordinates, from $x^\mu$ to $\bar{x}^a$, then the tensor components, with respect to a particular index, transform by either the Jacobian (contravariantly) or the inverse Jacobian (covariantly).
For a type $(1,1)$ tensor, a component transforms as follows.
\begin{equation}
    \bar{T}^\alpha_\beta = \frac{\partial \bar{x}^\alpha}{\partial x^\mu} \frac{\partial x^\nu}{\partial \bar{x}^\beta} T^\mu_\nu
\end{equation}
For the contravariant index $\alpha$, we get an entry from the Jacobian.
For the covariant $\beta$, we get an entry from the inverse Jacobian.
For a type $(p, q)$ tensor, the transformation law is
\begin{equation}
    \bar{T}^{\alpha_1 \ldots \alpha_p}_{\beta_1 \ldots \beta_q} = \frac{\partial \bar{x}^{\alpha_1}}{\partial x^{\mu_1}} \cdots \frac{\partial \bar{x}^{\alpha_p}}{\partial x^{\mu_p}} \frac{\partial x^{\nu_1}}{\partial \bar{x}^{\beta_1}} \cdots \frac{\partial x^{\nu_q}}{\partial \bar{x}^{\beta_q}} T^{\mu_1 \ldots \mu_p}_{\nu_1 \ldots \nu_q}.
\end{equation}
The tensor product of two tensors $S$ and $T$ of type $(p, q)$ and $(r, s)$ is a tensor $S \otimes T$ of type $(p + r, q + s)$.
As a multilinear map, it is evaluated at a given set of inputs by multiplying the tensors evaluated at those inputs.
Suppose that $x_{i}$ is a generic input and the ranks of the two tensors are $m$ and $n$, respectively.
Then,
\begin{equation}
    (S \otimes T)(x_1, \ldots, x_{m+n}) = S(x_1, \ldots, x_m) T(x_{m+1}, \ldots, x_{m+n}).
\end{equation}
This means that the components of $S \otimes T$ is the product of the components of the two individual tensors.
\begin{equation}
    (S \otimes T)^{i_1 \ldots i_{p+r}}_{j_1 \ldots j_{q+s}} = S^{i_1 \ldots i_p}_{j_1 \ldots j_q} T^{i_{p+1} \ldots i_{p+r}}_{j_{q+1} \ldots j_{q+s}}
\end{equation}
There are two main conventions for the meaning of tensor indices.
As shown above, we may use indices to denote a specific component of a tensor with respect to a particular basis.
For instance, $T^{\mu \nu}$ may refer to a generic component of a type $(2, 0)$ tensor.\footnote{It is also convention in orthodox physics for Greek indices to range from $0$ to $3$ for the four dimensions of spacetime. Latin indices range from 1 to 3 for the three dimensions of Euclidean space.}
We may call this component index notation, as opposed to abstract index notation.
In the latter, we use $T^{\mu \nu}$ to refer to a type $(2, 0)$ tensor, the actual object that does not depend on any basis.
The indices refer to the type of tensor, showing what slots it has for vectors and dual vectors, allowing for us to consider operations on tensors without being limited to a particular coordinate system.

Repeated indices in abstract index notation indicate summation over them just as they do for tensor components.
Einstein summation convention applies in the exact same way for abstract indices.
Of course, the objects being summed in either case are the actual components of the tensors, though implicit summation does not require any discussion of particular components.

When we use indices at all, they will usually be abstract indices unless stated otherwise or suggested by context.
Regardless, the nature of upper and lower indices as contravariant and covariant remains the same.
The idea may be extended to the whole tensor object, calling tensors with only upper indices contravariant and tensors with only lower indices covariant.
Tensors which have a mixture of upper and lower indices may be said to be mixed variance.

Whilst we have introduced type $(p, q)$ tensors as being elements of $T^p_q(V)$, this may hide a subtlety of tensors -- especially in physics.
Given any tensor space of the form $V \otimes \cdots \otimes V \otimes V^* \otimes \cdots V^*$, there exist natural braiding maps shuffling around the copies of $V$ and $V^*$.
Thus, a tensor $T \in V \otimes V^*$ is naturally associated with a tensor $G \in V^* \otimes V$.
Given a basis on $V$ and the dual basis on $V^*$, we have $\tensor{T}{^a_b} = T(\varepsilon^a, e_b) = G(e_b, \varepsilon^a) = \tensor{G}{_b^a}$.
It is thus not a problem to write $T^a_b$ without staggering the indices.

However, the tensor in $V^* \otimes V$ is often not $G$.
Suppose we have a symmetric nondegenerate bilinear form $g \in V^* \otimes V^*$, a type $(0, 2)$ tensor; as it is nondegenerate, we have an inverse $g^{-1} \in V \otimes V$, a type $(2, 0)$ tensor.
We also have isomorphisms $g^\flat \colon V \to V^*$, $v \mapsto g(v, \cdot)$ and $g^\sharp \colon V^* \to V$, $\omega \mapsto g^{-1}(\omega, \cdot)$.
We can construct an isomorphism $\mu \colon V \times V^* \to V^* \times V$ with $\mu(v, \omega) = (g^\flat(v), g^\sharp(\omega))$.
Then, we have the tensor $F \in V^* \otimes V$ defined by $F = T \circ \mu$.
A component $\tensor{F}{_b^a}$ of this tensor is not necessarily equal to a component of $\tensor{G}{_b^a}$, despite them both being bilinear maps $V \times V^* \to K$.

When dealing with tensors, we often mean $F$ instead of $G$, and we often refer to everything as just the tensor $T$.
Hence, there is a reason to distinguish between $\tensor{T}{^a_b}$ and $\tensor{T}{_b^a}$.
The tensor space which includes different types of tensors, distinguished by staggered indices, is $(V \oplus V^*)^{\otimes k}$.
This is the space of all rank $k$ tensors, where copies of $V$ and $V^*$ may be permuted in any order.
For instance, if $k = 2$, we have $(V \oplus V^*)^{\otimes 2} \cong (V \otimes V) \oplus (V \otimes V^*) \oplus (V^* \otimes V) \oplus (V^* \otimes V^*)$ by distributivity.
In general, if we denote such a space by $\mathcal{T}^k(V)$, then we have
\begin{equation}
    \mathcal{T}^k(V) \coloneq (V \oplus V^*)^{\otimes k} \cong \bigoplus_{p+q=k} T^p_q(V).
\end{equation}
We can use this to construct a similar space specifically for type $(p, q)$ tensors by requiring $p$ copies of $V$ or equivalently $q$ copies of $V^*$.
\begin{definition}[Staggered Tensor Space]\label{def:staggered_tensor_space}
    Given a finite-dimensional $K$-vector space $V$, the staggered tensor space of rank $k$ tensors on $V$ is defined as $\mathcal{T}^k(V) = (V \oplus V^*)^{\otimes k}$.
    The type $(p, q)$ staggered tensor space of $V$ is defined as
    \begin{equation}
        \mathcal{T}^p_q(V) = \operatorname{span}\{x_1 \otimes \cdots \otimes x_{p+q} \in \mathcal{T}^{p+q}(V) : x_i \in V \text{ or } x_i \in V^*, \#\{i : x_i \in V \} = p\}.
    \end{equation}
    Elements of this space are called type $(p, q)$ tensors.
\end{definition}
We could talk about elements of tensor spaces with the copies of $V$ and $V^*$ shuffled around beforehand, but now we have given a proper definition as to what a type $(p, q)$ tensor means in any context.
We will continue to work with $T^p_q(V)$ instead of $\mathcal{T}^p_q(V)$ in our definitions for simplicity.
As we will not yet concern ourselves with contractions via a metric, it does not particularly matter.
Regardless, the former may be substituted for the latter in subsequent definitions and results.

An operation that can be done with tensors is contraction.
This turns a type $(p, q)$ tensor into a type $(p-1, q-1)$ tensor, reducing the total rank by two.
For instance, a type $(1, 1)$ tensor $T^i_j$ may be contracted to a scalar by pairing the two indices together and writing $T^i_i$ to imply summation.
\begin{equation}
    T^i_i = \sum_i T^i_i = T^1_1 + \cdots + T^n_n
\end{equation}
In general, tensor contraction involves pairing a covariant index with a contravariant index and summing over them, denoted by using the same symbol for those two indices.
Whilst they can be defined like this, we can define them more generally using the universal property of the tensor product.
\begin{definition}[Tensor Contraction]\label{def:tensor_contraction}
    Let $V$ be a finite-dimensional $K$-vector space, and consider the type $(p, q)$ tensor space $T^p_q(V) = V^{\otimes p} \otimes (V^*)^{\otimes q}$ with $p, q \geq 1$.
    Denote $V^{\times p} \times (V^*)^{\times q}$ by $\times^p_q V$, and let $\otimes^p_q$ denote the canonical mapping between $\times^p_q V$ and $T^p_q(V)$; the latter space may be denoted by $\otimes^p_q V$.
    There is a canonical pairing $P \colon V \times V^* \to K$, $(v, \omega) \mapsto \omega(v)$ between a vector space and its dual.
    This bilinear map may be extended to a family of multilinear maps $P^{p, q}_{k, l} \colon \times^p_q V \to \otimes^{p-1}_{q-1} V$ via
    \begin{align}
        \begin{split}
        P^{p, q}_{k, l}(v_{(1)}, \ldots, v_{(p)}, \omega^{(1)}, \ldots, \omega^{(q)}) = \omega^{(l)}(v_{(k)}) v_{(1)} &\otimes \cdots \otimes \widehat{v}_{(k)} \otimes \cdots \otimes v_{(p)} \\
        \otimes \omega^{(1)} &\otimes \cdots \otimes \widehat{\omega}^{(l)} \otimes \cdots \otimes \omega^{(q)},
        \end{split}
    \end{align}
    where a hat above a symbol indicates omission.
    As these maps are multilinear, the universal property for the tensor product implies that there are unique linear maps $C^{p, q}_{k, l} \colon \otimes^p_q V \to \otimes^{p-1}_{q-1} V$, called contraction maps, making the following diagram commute.
    \[\begin{tikzcd}
        {\times^p_q V} && {\otimes^p_q V} \\
        \\
        && {\otimes^{p-1}_{q-1} V}
        \arrow["\otimes^p_q", from=1-1, to=1-3]
        \arrow["P^{p, q}_{k, l}"', from=1-1, to=3-3]
        \arrow["C^{p, q}_{k, l}", dashed, from=1-3, to=3-3]
    \end{tikzcd}\]
    Then, a contraction of the $k$-th contravariant index and the $l$-th covariant index of a type $(p, q)$ tensor $T \in \otimes^p_q V$ is defined to be $C^{p, q}_{k, l}(T)$.
\end{definition}
\begin{remark}
    The new tensor $C^{p, q}_{k, l}(T) = \widetilde{T}$ is of type $(p - 1, q - 1)$, and is explicitly given by
    \begin{equation}
        C^{p, q}_{k, l}(T) = \widetilde{T}^{a_1 \ldots \widehat{a}_k \ldots a_p}_{b_1 \ldots \widehat{b}_l \ldots b_q} = \sum_i T^{a_1 \ldots \widehat{a}_k i \ldots a_p}_{b_1 \ldots \widehat{b}_l i \ldots b_q},
    \end{equation}
    where a hat above a symbol denotes omission.
    If $p = q = 1$, the contraction of $T$ is called the trace of $T$ and is denoted by $\operatorname{tr}T$.
\end{remark}
The contraction of a tensor may be viewed as a generalization of the trace of a matrix.
As it has been defined entirely without the usage of a basis, contraction is independent of any choice of basis, being a fundamental operation on tensors.
Our definition of contraction involves just a single tensor, though we can extend it to contracting one tensor with another by combining contraction with the tensor product.
For instance, consider two type $(1, 1)$ tensors $A = A^a_b$ and $B = B^c_d$.
We can contract the index $a$ of $A$ with the index $d$ of $B$ by considering the contraction of the tensor product $A \otimes B = A^a_b B^c_d$.\footnote{Here, by juxtaposing the two tensor symbols to denote their tensor product, abstract index notation mimics the way tensor components are multiplied when taking the tensor product of two tensors.}
This is given by
\begin{equation}
    A^a_b B^c_a = \sum_a A^a_b B^c_a.
\end{equation}
In general, we are able to contract the covariant index of one tensor with the contravariant index of another tensor which may come before or afterwards; contraction is not a commutative operation, though.
The essence of tensor contractions is that we canonically pair vectors with dual vectors, cancelling out such a pair from the total rank.
We can take familiar objects that may be represented as tensors and see what contraction may be interpreted as.
In the case of matrices, the above sum is the component form of the product $\mathbf{BA}$, taking the upper index to be the row index, so contracting two matrices is just multiplying them.

The contraction of a vector with a dual vector to give a scalar is done by the usual pairing.
We cannot contract two vectors with each other, however, without a way to naturally pair them (they both only have a contravariant index).
A symmetric nondegenerate bilinear form provides such a way, and the resulting scalar is the inner product of the two vectors.
The same is the case for dual vectors.

\subsection{Tensor Algebra}\label{subsec:multilinear_algebra:tensor_algebra}

We will now introduce an important structure.
It is interesting in its own right, but our main use for it will be to define other important structures.
By considering all tensor powers of a vector space, we obtain the tensor algebra of the space.
However, just as the tensor product may be defined via a universal property, so too can the tensor algebra.
\begin{definition}[Tensor Algebra]\label{def:tensor_algebra}
    Let $V$ be a $K$-vector space.
    The tensor algebra of $V$ is the associative unital algebra $T(V)$, also denoted $T^\bullet(V)$, with a vector space inclusion $\iota : V \hookrightarrow T(V)$ such that for any linear map $f \colon V \to A$ from $V$ to an associative unital $K$-algebra $A$, there exists a unique algebra homomorphism $\tilde{f} \colon T(V) \to A$ making the following diagram commute.
    \[\begin{tikzcd}
        {V} && {T(V)} \\
        \\
        && {A}
        \arrow["\iota", hook, from=1-1, to=1-3]
        \arrow["f"', from=1-1, to=3-3]
        \arrow["\tilde{f}", dashed, from=1-3, to=3-3]
    \end{tikzcd}\]
\end{definition}
\begin{remark}
    An explicit construction of the tensor algebra of $V$ is given, as described before, by taking the direct sum of all tensor powers of $V$.
    Multiplication in $T(V)$ is given by the tensor product.
    \begin{equation}
        T(V) = \bigoplus_{k=0}^\infty V^{\otimes k} = K \oplus V \oplus (V \otimes V) \oplus (V \otimes V \otimes V) \oplus \cdots
    \end{equation}
    Being an algebra homomorphism, the induced map $\tilde{f}$ satisfies
    \begin{equation}
        \tilde{f}(v_{(1)} \otimes \cdots \otimes v_{(k)}) = \tilde{f}(v_{(1)}) \cdots f(v_{(k)}).
    \end{equation}
\end{remark}
It can be seen that the tensor algebra is a $\Z_{\geq 0}$-graded algebra.
That is, the algebra is split into different indexed parts, and elements from different parts being combined in a particular way: $T^k (V) \otimes T^l(V) \to T^{k+l}(V)$.
By defining $T^k(V) = \{0\}$ for negative integers $k$ and appending them to the tensor algebra, we get $\Z$ grading.
Another place where grading is present is with polynomials and their degrees.
The similarity more explicitly manifests itself in the fact that the tensor algebra is essentially the algebra of polynomials over $K$ in $n$ non-commuting variables.

The connection between the tensor algebra and non-commuting algebra comes directly from its universal property which defines it as the most free unital associative algebra.
It is free in the sense that it is the most general algebra generated by $V$, and any other algebra generated by $V$ can be linked to it.
Tensors in a certain tensor power are as unrelated to tensors in different tensor powers as polynomials with different degrees are.
It is hence reasonable to use the word degree to refer to the rank of a tensor in the tensor algebra, though this word may be used in place of rank in any situation.

Due to its generality, we can use the tensor algebra to define various other algebras as quotients of it.
These constructed algebras could also be defined via universal properties similar to the one for the tensor algebra.
However, due to the extreme similarity, they are not particularly illuminating and thus have not been used as definitions.
The first algebra we define from the tensor algebra will be the symmetric algebra.
\begin{definition}[Symmetric Algebra]\label{def:symmetric_algebra}
    Let $V$ be a finite-dimensional $K$-vector space.
    Let $\mathfrak{s}$ be the two-sided ideal of $T(V)$ generated by elements of the form $v \otimes w - w \otimes v$ for $v, w \in V$.
    The symmetric algebra of $V$, denoted $\Sym(V)$, is defined to be the quotient $T(V) / \mathfrak{s}$.
    The symmetric product $\odot \colon \Sym(V) \times \Sym(V) \to \Sym(V)$ is the operation descended from the tensor product, defined by
    \begin{equation}
        v \odot w \coloneq v \otimes w \pmod{\mathfrak{s}}.
    \end{equation}
    By letting $\mathfrak{s}_k = \mathfrak{s} \cap T^k(V)$, the $k$-th symmetric power of $V$ is defined to be the quotient $\Sym^k(V) = T^k(V) / \mathfrak{s}_k$.
    This is also denoted by $V^{\odot k}$.
    An element of the symmetric algebra is called a symmetric tensor.
\end{definition}
The symmetric algebra is what we get by forcing the tensor algebra to be commutative by collapsing necessary expressions to zero.
Elements of the quotient do not have well-defined components.
However, as $K$ has characteristic zero, $\sym$ is a projection of $T^k(V)$ onto the subspace of tensors $T$ with $\tau_\sigma T = T$ for all $\sigma \in S_k$, whose kernel is $\mathfrak{s}_k$, so $\Sym^k(V)$ may be identified with that subspace.
To see that the kernel is $\mathfrak{s}_k$, note that $\sym$ vanishes on the generators of $\mathfrak{s}_k$, and that for any tensor $T$, the difference $T - \sym(T) = \frac{1}{k!} \sum_{\sigma \in S_k} (T - \tau_\sigma T)$ lies in $\mathfrak{s}_k$, as each permutation is a product of adjacent transpositions and each of these changes a tensor by an element of $\mathfrak{s}_k$.
If $\sym(T) = 0$, then $T$ lies in $\mathfrak{s}_k$.
Such symmetric tensors are symmetric in their components, e.g. for $T \in T^2(V)$, we have $T^{ij} = T^{ji}$.
This can be thought of in a coordinate-independent way by considering them symmetric in their arguments, e.g. $T(v, w) = T(w, v)$.
In general, symmetric tensors are invariant under any permutation of their arguments.
Given a braiding map $\tau_\sigma$ associated with a permutation $\sigma \in S_k$, where $S_k$ is the symmetric group of degree $k$, $\tau_\sigma T = T$.

Although not all tensors are symmetric, we can turn them symmetric via symmetrization maps.
This is done by taking the average of all permutations of a tensor, so any further permutations will only rearrange terms of a sum.
\begin{definition}[Tensor Symmetrization]\label{def:tensor_symmetrization}
    Let $V$ be an $n$-dimensional $K$-vector space.
    Consider the braiding maps $\tau_\sigma \colon v_{(1)} \otimes \cdots \otimes v_{(k)} \mapsto v_{\sigma(1)} \otimes \cdots \otimes v_{\sigma(k)}$ for $\sigma \in S_k$.
    The symmetric part of a tensor $T \in T^k(V)$ is defined by
    \begin{equation}
        \sym(T) \coloneq \frac{1}{k!} \sum_{\sigma \in S_k} \tau_\sigma(T).
    \end{equation}
    This may be denoted by $T_{(i_1 \ldots i_k)}$.
    The symmetrization of $T$ over the slots in $I \subseteq \{1, \ldots, k\}$ is given by
    \begin{equation}
        \sym_I(T) = \frac{1}{\abs{S_I}} \sum_{\sigma \in S_I} \tau_\sigma(T),
    \end{equation}
    where $S_I$ is the subgroup of $S_k$ permuting only the slots in $I$.
    This may be denoted by placing parentheses between the indices to be symmetrized and vertical bars inside the parentheses to indicate indices to be excluded, e.g. $T_{(\alpha \lvert\beta\rvert \gamma)} = 1/2 (T_{\alpha \beta \gamma} + T_{\gamma \beta \alpha})$.
\end{definition}
\begin{remark}
    The factor $1/k!$ in the definition requires $k!$ to be invertible in $K$, which our assumption that $K$ has characteristic zero guarantees.
    The quotient construction of $\Sym(V)$ works over any field, which is a reason to prefer it, as we have done by choosing it as the definition.
\end{remark}
We can take our intrinsic definition of the symmetrization of a tensor and put it in terms of coordinates, obtaining the following equation.
\begin{equation}
    \sym{T} = \frac{1}{k!} \sum_{\sigma \in S_k} T^{i_{\sigma(1)} \ldots i_{\sigma(k)}} e_{i_1} \otimes \cdots \otimes e_{i_k}
\end{equation}
Note that our definition of symmetric tensors and tensor symmetrization has only been with respect to elements of a tensor power.
That is, we have only considered tensors of unmixed variance.
It does not make sense to swap the arguments of a tensor and obtain the same map.
However, we can still have a tensor such as $T^\alpha_{\beta \gamma}$ where the $\beta$ and $\gamma$ indices can be symmetrized.
Thus, we can discuss symmetry (or antisymmetry) when we have a tensor with at least two indices of the same variance.

Just as we constructed symmetric tensors, we can similarly construct antisymmetric tensors.
\begin{definition}[Antisymmetric Algebra]\label{def:antisymmetric_algebra}
    Let $V$ be a finite-dimensional $K$-vector space.
    Let $\mathfrak{a}$ be the two-sided ideal of $T(V)$ generated by elements of the form $v \otimes w + w \otimes v$ for $v, w \in V$.
    The antisymmetric algebra of $V$, denoted $\Asym(V)$, is defined to be the quotient $T(V) / \mathfrak{a}$.
    The multiplication operation is the one induced by the tensor algebra.
    By letting $\mathfrak{a}_k = \mathfrak{a} \cap T^k(V)$, the $k$-th antisymmetric power of $V$ is defined to be the quotient $\Asym^k(V) = T^k(V) / \mathfrak{a}_k$.
\end{definition}
In the same way, $\asym$ identifies $\Asym^k(V)$ with the subspace of antisymmetric tensors, whose components satisfy $T^{ij} = -T^{ji}$, which is the antisymmetry we wish for.
Again, this extends to general tensors (of unmixed variance) by considering permutations on all indices or permutations on all arguments.
The final result depends on whether the permutation is even or odd.
This is summarized by the equation $\tau_\sigma (T) = \sgn(\sigma) T$ for every braiding map.

Another similarity is that we can construct the antisymmetric part of any tensor.
Instead of taking the average, we include a factor that alternates signs with respect to the parity of the permutation.
\begin{definition}[Tensor Antisymmetrization]\label{def:tensor_antisymmetrization}
    Let $V$ be an $n$-dimensional $K$-vector space.
    Consider the braiding maps $\tau_\sigma \colon v_{(1)} \otimes \cdots \otimes v_{(k)} \mapsto v_{\sigma(1)} \otimes \cdots \otimes v_{\sigma(k)}$ for $\sigma \in S_k$.
    The antisymmetric part of a tensor $T \in T^k(V)$ is defined by
    \begin{equation}
        \asym(T) \coloneq \frac{1}{k!} \sum_{\sigma \in S_k} \sgn(\sigma) \tau_\sigma(T).
    \end{equation}
    This may be denoted by $T_{[i_1 \ldots i_k]}$.
    The antisymmetrization of $T$ over the slots in $I \subseteq \{1, \ldots, k\}$ is given by
    \begin{equation}
        \asym_I(T) = \frac{1}{\abs{S_I}} \sum_{\sigma \in S_I} \sgn(\sigma) \tau_\sigma(T),
    \end{equation}
    where $S_I$ is the subgroup of $S_k$ permuting only the slots in $I$.
    This may be denoted by placing square brackets between the indices to be antisymmetrized and vertical bars inside the brackets to indicate indices to be excluded, e.g. $T_{[\alpha \lvert\beta\rvert \gamma]} = 1/2 (T_{\alpha \beta \gamma} - T_{\gamma \beta \alpha})$.
\end{definition}
A very similar construction to that of the antisymmetric algebra is the alternating tensor algebra.
The properties of antisymmetry and being alternating are nearly identical, though they differ from each other in specific cases.
Regardless of that fact, our main reason for defining the following algebra is to have it be a subspace of the tensor algebra instead of a quotient.
It will then help us when dealing with the exterior algebra.
\begin{definition}[Alternating Multilinear Map]\label{def:alternating_multilinear_map}
    Let $R$ be a commutative ring and let $V$ and $W$ be $R$-modules.
    A $k$-multilinear map $f \colon V^{\times k} \to W$ is alternating if $f(x_1, \ldots, x_k) = 0$ whenever there exist $i, j \in \{1, \ldots, k\}$ such that $x_i = x_j$ and $i \neq j$.
    The space of such maps is denoted by $\Alt^k(V, W)$.
\end{definition}
\begin{definition}[Alternating Tensor Algebra]\label{def:alternating_tensor_algebra}
    Let $V$ be an $n$-dimensional $K$-vector space.
    Let $\Lambda^k(V)$ be the subspace of $T^k(V)$ consisting of alternating tensors of order $k$, i.e. elements of $\Alt^k(V^*, K)$.
    The alternating tensor algebra of $V$, denoted by $\Lambda(V)$, is defined by
    \begin{equation}
        \Lambda(V) \coloneq \bigoplus_{k=0}^n \Lambda^k(V).
    \end{equation}
    Multiplication on the algebra is called the wedge product and is defined using the alternation map $\alt \colon T(V) \to \Lambda(V)$ gradewise.
    \begin{equation}
        \alt(v_{(1)} \otimes \cdots \otimes v_{(k)}) \coloneq \frac{1}{k!} \sum_{\sigma \in S_k} \sgn(\sigma) v_{(\sigma(1))} \otimes \cdots \otimes v_{(\sigma(k))}
    \end{equation}
    We have the wedge product between two tensors $S$ and $T$ of rank $p$ and $q$ given by
    \begin{equation}
        S \wedge T \coloneq \frac{(p+q)!}{p!q!} \alt(S \otimes T).
    \end{equation}
\end{definition}
\begin{remark}
    There are two conventions with the definition of the wedge product.
    One involves the factor in front of the alternating map that we used.
    The other involves no such factor, so $S \wedge T = \alt(S \otimes T)$.
    Our choice of convention leads to cleaner formulae without any additional burden -- besides the initial definition.
\end{remark}

\subsection{Exterior Algebra}\label{subsec:multilinear_algebra:exterior_algebra}

The exterior algebra is the most important algebra we will construct from the tensor algebra, justifying its own section.
Whilst we will construct it as a quotient algebra, it is important enough to be defined via its universal property.
\begin{definition}[Exterior Algebra]\label{def:exterior_algebra}
    Let $V$ be a finite-dimensional $K$-vector space.
    The exterior algebra of $V$ is the associative unital algebra $\extp(V)$, also denoted $\extp^\bullet(V)$, with the vector space inclusion $\iota \colon V \hookrightarrow \extp(V)$ such that for any associative unital algebra $A$ and any multilinear map $f \colon V \to A$ satisfying $f(v)f(v) = 0$ for all $v \in V$, there exists a unique algebra homomorphism $\tilde{f} \colon \extp(V) \to A$ making the following diagram commute.
    \[\begin{tikzcd}
            {V} && {\extp(V)} \\
            \\
            && {A}
            \arrow["\iota", hook, from=1-1, to=1-3]
            \arrow["f"', from=1-1, to=3-3]
            \arrow["\tilde{f}", dashed, from=1-3, to=3-3]
    \end{tikzcd}\]
    The multiplication operation on $\extp(V)$ is called the exterior or wedge product, denoted by $\wedge$.
    The $k$-th exterior power of $V$ is the vector space $\extp^k(V)$ with an alternating multilinear map $\iota_k \colon V^{\times k} \to \extp^k(V)$ such that every alternating multilinear map $f \colon V^{\times k} \to W$ to a vector space $W$ factors uniquely as $f = \tilde{f} \circ \iota_k$ with $\tilde{f}$ linear.
    It is also denoted by $V^{\wedge k}$.
    \[\begin{tikzcd}
            {V^{\times k}} && {\extp^k(V)} \\
            \\
            && {W}
            \arrow["\iota_k", from=1-1, to=1-3]
            \arrow["f"', from=1-1, to=3-3]
            \arrow["\tilde{f}", dashed, from=1-3, to=3-3]
    \end{tikzcd}\]
    An element of the exterior algebra may be called an exterior tensor, an exterior form, or a multivector.
    The rank of an element is also commonly called its grade or degree.
    Grade two elements may be called bivectors, grade three elements may be called trivectors, and grade $k$ elements may be called $k$-vectors.
\end{definition}
Now, we will construct the exterior algebra as a quotient of the tensor algebra so that the descended product is alternating.
\begin{proposition}\label{prop:exterior_algebra_construction_satisfies_universal_property}
    Let $V$ be a finite-dimensional $K$-vector space.
    Let $\mathfrak{e}$ be the two-sided ideal of $T(V)$ generated by elements of the form $v \otimes v$ for $v \in V$.
    The exterior algebra of $V$, denoted $\extp(V)$, is defined to be the quotient $T(V) / \mathfrak{e}$.
    The exterior product or wedge product $\wedge \colon \extp(V) \times \extp(V) \to \extp(V)$ is the operation descended from the tensor product, defined by
    \begin{equation}
        v \wedge w \coloneq v \otimes w \pmod{\mathfrak{e}}.
    \end{equation}
    By letting $\mathfrak{e}_k = \mathfrak{e} \cap T^k(V)$, the $k$-th exterior power of $V$, denoted $\extp^k(V)$ or $V^{\wedge k}$, is defined to be the quotient $T^k(V) / \mathfrak{e}_k$.
    The exterior algebra and the $k$-th exterior power satisfy their universal properties.
\end{proposition}
\begin{proof}
    Let $f \colon V \to A$ be a linear map such that $f(v)f(v) = 0$ for all $v \in V$.
    By the universal property of the tensor algebra, there is a unique algebra homomorphism $\tilde{f} \colon T(V) \to A$ extending $f$.
    We have $\tilde{f}(v \otimes v) = f(v)f(v)$, and this must equal zero for all $v \in V$.
    Thus, $\mathfrak{e} \subseteq \ker \tilde{f}$.
    It follows that there is a unique algebra homomorphism $\varphi \colon T(V) / \mathfrak{e} \to A$ such that $\varphi(q(x)) = \tilde{f}(x)$ for all $x \in T(V)$, where $q(v \otimes w) = v \wedge w$ is just a label for the wedge product.
    
    The extension $\varphi$ is unique, and its uniqueness as a homomorphism satisfying the universal property follows $\varphi(\iota(v)) = f(v)$ by the universal property and the fact that $\iota(V)$ generated $\extp(V)$.
    The map $\iota$ is injective, since $\mathfrak{e}$ is generated by tensors of degree two and so has no component in degree one.
    Therefore, the constructed exterior algebra satisfies its universal property.
    Similar reasoning for a multilinear map on $V^{\times k}$ shows that the constructed $k$-th exterior power of $V$ satisfies its universal property as well.
\end{proof}
\begin{remark}
    In some texts, the ideal used to construct the exterior algebra is generated by elements of the form $v \otimes w + w \otimes v$ for $v, w \in V$.
    This makes the exterior algebra antisymmetric, not alternating.
    Over our field $K$ of characteristic zero the two conditions are equivalent.
    In characteristic two, the property of being alternating cannot be derived from antisymmetry, so it is the more fundamental property of the exterior algebra.
\end{remark}
The exterior algebra is essentially just the tensor algebra with the additional property of being alternating.
It inherits its grading from the tensor algebra, and it inherits a universal property.
The exterior algebra can thus be considered the most general unital associative algebra where multiplication is alternating.

Viewing the exterior algebra as a quotient is direct, whereas viewing it as a subspace of the tensor algebra requires more work.
The two constructions are related by the following isomorphism, which rescales by $k!$ in degree $k$ because of the convention for the wedge product on $\Lambda(V)$.
\begin{proposition}\label{prop:exterior_algebra_quotient_subspace_equivalence}
    Let $V$ be an $n$-dimensional $K$-vector space.
    Let $\extp(V) = T(V) / \mathfrak{e}$ carry the product induced by $\otimes$, and let $\Lambda(V)$ carry the wedge product of \cref{def:alternating_tensor_algebra}.
    The map $\Phi \colon \extp(V) \to \Lambda(V)$ given in degree $k$ by
    \begin{equation}
        \Phi(v_1 \wedge \cdots \wedge v_k) = k! \, \alt(v_1 \otimes \cdots \otimes v_k) = \sum_{\sigma \in S_k} \sgn(\sigma) \, v_{\sigma(1)} \otimes \cdots \otimes v_{\sigma(k)}
    \end{equation}
    is an isomorphism of graded algebras.
\end{proposition}
\begin{proof}
    We will use the fact that
    \begin{equation}\label{eq:alt_absorbs_alt}
        \alt(\alt(S) \otimes T) = \alt(S \otimes T) = \alt(S \otimes \alt(T))
    \end{equation}
    for $S \in T^p(V)$ and $T \in T^q(V)$.
    To see this, note that $\alt(S) \otimes T = \frac{1}{p!} \sum_{\rho \in S_p} \sgn(\rho) \, \tau_\rho(S \otimes T)$, where $S_p$ is embedded in $S_{p+q}$ as the permutations fixing the last $q$ slots.
    Since $\alt \circ \tau_\rho = \sgn(\rho) \alt$, the first equality follows, and the second is shown in the same way.
    Let $\tilde{\Phi} \colon T(V) \to \Lambda(V)$ be the linear map that is equal to $k! \, \alt$ on $T^k(V)$.

    First, we show that $\tilde{\Phi}$ is an algebra homomorphism, so we need $\tilde{\Phi}(S \otimes T) = \tilde{\Phi}(S) \wedge \tilde{\Phi}(T)$.
    By the definition of the wedge product and \eqref{eq:alt_absorbs_alt},
    \begin{equation*}
        \tilde{\Phi}(S) \wedge \tilde{\Phi}(T) = \frac{(p+q)!}{p!\,q!} \alt\bigl(p! \, \alt(S) \otimes q! \, \alt(T)\bigr) = (p+q)! \, \alt(S \otimes T) = \tilde{\Phi}(S \otimes T).
    \end{equation*}
    Next, we show that $\tilde{\Phi}$ descends to a map $\Phi$ on $T(V) / \mathfrak{e}$, so we need $\mathfrak{e} \subseteq \ker \tilde{\Phi}$.
    The ideal $\mathfrak{e}$ is spanned by the tensors $x \otimes (v \otimes v) \otimes y$.
    We have $\tilde{\Phi}(v \otimes v) = 2 \alt(v \otimes v) = v \otimes v - v \otimes v = 0$, so being an algebra homomorphism, $\tilde{\Phi}$ sends $x \otimes (v \otimes v) \otimes y$ to zero.
    It follows that $\Phi$ is a well-defined algebra homomorphism, and it preserves degree.

    Finally, we show that $\Phi$ is bijective.
    Every $t \in \Lambda^k(V)$ satisfies $\alt(t) = t$, so $t = \Phi\bigl(\tfrac{1}{k!}[t]\bigr)$, and $\Phi$ is surjective.
    For injectivity, define $\Psi \colon \Lambda^k(V) \to \extp^k(V)$ by $\Psi(t) = \frac{1}{k!}[t]$.
    Since $(v + w) \otimes (v + w) - v \otimes v - w \otimes w = v \otimes w + w \otimes v \in \mathfrak{e}$, swapping two adjacent factors changes the sign modulo $\mathfrak{e}$, so $[v_{\sigma(1)} \otimes \cdots \otimes v_{\sigma(k)}] = \sgn(\sigma) [v_1 \otimes \cdots \otimes v_k]$.
    Therefore
    \begin{equation*}
        \Psi(\Phi(v_1 \wedge \cdots \wedge v_k)) = \frac{1}{k!} \sum_{\sigma \in S_k} \sgn(\sigma)^2 [v_1 \otimes \cdots \otimes v_k] = v_1 \wedge \cdots \wedge v_k.
    \end{equation*}
    Thus $\Psi \circ \Phi$ is the identity on generators, so it is the identity on $\extp^k(V)$, and $\Phi$ is injective. \qedhere
\end{proof}
Now, we can view the exterior algebra in a less abstract sense.
Multivectors, as their name suggests, are a generalization of vectors.
It follows that they have myriad applications in mathematics and physics.
In the setting of an affine space, multivectors may be thought of as oriented volumes.
Vectors are known to be represented by oriented lines (arrows).
Bivectors represent oriented areas -- which may be used to model physical quantities such as torque -- and trivectors represent oriented volumes.
This extends to any dimension.

We often do not need to explicitly work with the tensor product when dealing with the exterior algebra; we may stick to just using the wedge product.
The wedge product is associative, distributive, and alternating, and it is therefore also antisymmetric.
The wedge product of two multivectors $v \in \extp^p(V)$ and $w \in \extp^q(V)$ belongs in $\extp^{p+q}(V)$ and exhibits graded antisymmetry: $v \wedge w = (-1)^{pq} w \wedge v$.

Given a basis $\{e_1, \ldots, e_n\}$ for $V$, we can construct a basis for $\extp^k(V)$.
It is formed by elements of the form $e_{i_1} \wedge \cdots \wedge e_{i_k}$ for $1 \leq i_1 < \cdots < i_k \leq n$.
Any $k$-vector $v$ can be expressed with respect to this basis as
\begin{equation}
    v = v^{i_1 \ldots i_k} e_{i_1} \wedge \cdots \wedge e_{i_k}.
\end{equation}
Here the sum is over the increasing multi-indices $i_1 < \cdots < i_k$ only, as for differential forms later on.
The dimension of $\extp^k(V)$ is $\binom{n}{k}$.
These elements span $\extp^k(V)$, since a product with a repeated factor vanishes and any other product can be reordered into increasing order up to sign.
They are linearly independent: for an increasing multi-index $I$, let $\varepsilon^I \colon V^{\times k} \to K$ be the alternating map $(v_1, \ldots, v_k) \mapsto \det(\varepsilon^{i_a}(v_b))$ built from the dual basis.
It induces a linear functional on $\extp^k(V)$, because it vanishes on the tensors $x \otimes v \otimes v \otimes y$ that span $\mathfrak{e}_k$, and this functional takes the value $\delta_{IJ}$ on $e_J$.
The dimension of $\extp(V)$ is the sum of the dimensions of the individual exterior powers, which we know to be $2^n$.
If we take the wedge product of $k$ vectors, the resulting $k$-vector $w$ is given by
\begin{equation}
    w = \det (A^{i_1 \ldots i_k}) e_{i_1} \wedge \cdots \wedge e_{i_k},
\end{equation}
where $\det(A_{i_1 \ldots i_k})$ is the determinant of the $k \times k$ submatrix formed by selecting the rows $i_1, \ldots, i_k$ from the matrix $\mathbf{A}$ given by
\begin{equation}
    \mathbf{A} = \begin{bmatrix}
v^1_1 & \cdots & v^1_n \\
\vdots & \ddots & \vdots \\
v^n_1 & \cdots & v^n_n 
\end{bmatrix}.
\end{equation}
Multivectors, similar to any other type of tensor, can be considered as multilinear maps.
From the isomorphism $\extp(V) \cong \Lambda(V)$, they are alternating tensors, which are alternating $k$-multilinear maps in $\Alt^k(V^*, K)$.
Considering a linear transformation $L \colon V \to V$, a type $(1, 1)$ tensor, we have the induced linear map $\extp^n(L) \colon \extp^n(V) \to \extp^n(V)$ given by $v^{(1)} \wedge \cdots \wedge v^{(n)} \mapsto Lv^{(1)} \wedge \cdots \wedge Lv^{(n)}$.
Since $\extp^n(V)$ is one dimensional, $\extp^n(L)$ is given by a scalar multiple of $v^{(1)} \wedge \cdots \wedge v^{(n)}$.
This coefficient is the determinant of the matrix of $L$ with respect to any basis of $V$.
In fact, we can use this to define the determinant in a general, intrinsic sense.
\begin{definition}[Determinant]\label{def:determinant}
    Let $V$ be a finite-dimensional $K$-vector space.
    Let $L \colon V \to V$ be a linear transformation of an $n$-dimensional vector space $V$.
    This induces the linear map $\extp^n(L) \colon \extp^n(V) \to \extp^n(V)$ given by $v^{(1)} \wedge \cdots \wedge v^{(n)} \mapsto Lv^{(1)} \wedge \cdots \wedge Lv^{(n)}$.
    The determinant of $L$ is the element in $K$ such that
    \begin{equation}\label{eq:determinant_definition}
        \extp^n(L)(v^{(1)} \wedge \cdots \wedge v^{(n)}) = \det(L) \cdot v^{(1)} \wedge \ldots \wedge v^{(n)}
    \end{equation}
    for all $v^{(i)} \in V$.
    Alternatively, once a basis $(e_1, \ldots, e_n)$ of $V$ is chosen, the determinant of a matrix is an alternating multilinear map of its columns, namely $\det(v_1, \ldots, v_n)$ defined by $v_1 \wedge \cdots \wedge v_n = \det(v_1, \ldots, v_n) \, e_1 \wedge \cdots \wedge e_n$.
\end{definition}
The determinant is one example of an alternating multilinear map, and we can consider them in general.
This is useful since they are the prototypical version of differential forms.
\begin{definition}[Exterior Form]\label{def:exterior_form}
    Let $V$ be a finite-dimensional $K$-vector space.
    An (exterior) $k$-form on $V$ is an alternating multilinear map $\omega \colon V^{\times k} \to K$, an element of $\Alt^k(V, K)$.
\end{definition}
Such a map is naturally associated with an element in the dual space $(\extp^k(V))^*$.
This amounts to a map of the wedge product of $k$ vectors being essentially the same as that map when considered to take a list of those $k$ vectors as input instead.
Furthermore, $\Alt^k(V, K) \cong (\extp^k(V))^*$ naturally.
If $V$ is finite-dimensional, then $(\extp^k(V))^* \cong \extp^k(V^*)$ naturally; the spaces have dimension $\binom{n}{k}$.
We can take two forms and consider the wedge product between them.
\begin{definition}[Wedge Product of Forms]\label{def:wedge_product_of_forms}
    Let $V$ be a finite-dimensional $K$-vector space.
    Let $\omega \in \Alt^k(V, K)$ and $\eta \in \Alt^l(V, K)$ be alternating multilinear maps.
    The wedge product between them is an alternating multilinear map $\omega \wedge \eta \in \Alt^{k+l}(V, K)$ defined by
    \begin{equation}
        \omega \wedge \eta \coloneq \frac{(k+l)!}{k!l!} \alt(\omega \otimes \eta),
    \end{equation}
    where
    \begin{equation}
        \alt(\omega)(x_1, \ldots, x_k) \coloneq \frac{1}{k!} \sum_{\sigma \in S_k} \sgn(\sigma) \omega(x_{\sigma(1)}, \ldots, x_{\sigma(k)}).
    \end{equation}
\end{definition}
This will be expanded upon when we introduce differential forms.

\section{Manifolds}\label{sec:manifolds}

\subsection{Topological Manifolds}\label{subsec:manifolds:topological_manifolds}

The core concept of a manifold is that it resembles Euclidean space near any particular point on the manifold.
The most basic type of manifold, which all other manifolds are special cases of, are topological manifolds.
They are sets with additional structure.
This structure makes it a topological space at a minimum, though it has further conditions to exclude certain pathological spaces from being considered manifolds, e.g. the long line \cite{nlab_long_line}.
Locally resembling Euclidean space can be made precise using homeomorphisms and neighbourhoods.

Note that for all definitions, unless otherwise stated, we consider $\R^n$ to be a topological space with its metric topology.
Also, $\R^0$ is defined to be $\{0\}$ with the discrete topology.
\begin{definition}[Locally Euclidean Space]\label{def:locally_euclidean_space}
    A topological space $X$ is locally Euclidean if for every $x \in X$, there are open neighbourhoods $U \subseteq X$, $V \subseteq \R^n$ for some $n \in \Z_{\geq 0}$, and a homeomorphism $\varphi \colon U \to V$.
    The neighbourhood $U$ may be called locally Euclidean.
\end{definition}
Thus, the precise description is that every point of a manifold has an open neighbourhood homeomorphic to $\R^n$.
We require that a topological manifold is Hausdorff.
This is a very sensible condition, for we very often wish to perform calculus with manifolds, and the uniqueness of limits is a necessity.
We also require second countability.
Together with the other conditions, it implies paracompactness \cite[Thm.~1.15]{Lee2012} and thereby the existence of partitions of unity \cite[Thm.~2.23]{Lee2012}, which are a very useful tool.
\begin{definition}[Topological Manifold]\label{def:topological_manifold}
    A topological manifold $(M, \mathcal{T})$ is a topological space which is locally Euclidean, Hausdorff, and second countable.
    If the locally Euclidean neighbourhoods $U \cong \R^n$ are all of dimension $n$ for some fixed $n \in \Z_{\geq 0}$, then $(M, \mathcal{T})$ is called an $n$-dimensional manifold, an $n$-manifold, or an $n$-fold.
    As with other topological spaces, and more generally mathematical spaces, we often call the underlying set $M$ the manifold.
\end{definition}
\begin{remark}
    Some authors require paracompactness in place of second countability.
    The two definitions agree for spaces with countably many connected components, since a connected, paracompact, locally Euclidean Hausdorff space is $\sigma$-compact and hence second countable.
    Paracompactness alone allows uncountably many components, as in an uncountable disjoint union of copies of $\R$.
\end{remark}
    To be able to describe points on a manifold with a coordinate system, we make use of charts.
The definition of charts comes directly from homeomorphisms with Euclidean space.
\begin{definition}[Coordinate Chart]\label{def:coordinate_chart}
    A coordinate chart $(U, \varphi)$ of an $n$-dimensional topological manifold $M$ is an open set $U \subseteq M$ with its corresponding homeomorphism $\varphi \colon U \to V$, where $V$ is the associated open subset of $\R^n$.
    We call $U$ a coordinate domain and a coordinate neighbourhood of each of its points.
    The homeomorphism $\varphi$ of a coordinate chart $(U, \varphi)$ is often referred to as the chart itself (specifying the domain specifies the open set), as too is $U$ itself.
    A chart is centred at $p$ if $\varphi(p) = 0$.
\end{definition}

\begin{figure}[!htbp]
    \centering
    \begin{tikzpicture}[scale=1.5, thick]
        \draw plot [smooth cycle] coordinates {(-1.4, -0.8) (0.1, -0.9) (0.9, -1.2) (1.5, -0.7) (1.6, 0.5) (0.4, 1.2) (-1.9, 0.7)};

        \draw[fill=red!40, opacity=0.5] (0,0) ellipse (0.5 and 0.35);
        \node at (0,0) [right] {$p$};
        \filldraw[black] (0,0) circle (0.05);
        \node at (0.55, 0) [right] {$U$};
        \node at (1.7, 0.7) [right] {$M$};

        \begin{scope}[shift={(-2.1, -5.25)}]
            \draw (0, 0) -- (3.5, 0) -- (4, 1.8) -- (0.5, 1.8) -- cycle;
            \node at (3.8, 1.9) [above right] {$\R^n$};

            \draw[fill=cyan!40, opacity=0.5] (2, 0.9) ellipse (0.7 and 0.5);
            \node at (2, 0.9) [above] {$\varphi(p)$};
            \filldraw[black] (2, 0.9) circle (0.05);
            \node at (2.75, 0.9) [right] {$\varphi(U)$};
        \end{scope}

        \draw[-{Stealth[length=4mm, width=3mm]}] (0, -1.2) -- +(0, -2.0) node [pos=0.5, right] {$\varphi$};
    \end{tikzpicture}
    \caption{A coordinate chart}\label{fig:coordinate_chart}
\end{figure}

As is shown in \Cref{fig:coordinate_chart}, the chart maps a part of the manifold to a part of Euclidean space, thus providing the manifold with a coordinate system.
Note that the coordinate system is not necessarily a global one, as a manifold may not necessarily be able to be covered using just one chart.
A cover of a manifold using charts is called an atlas, as it is akin to atlases of the Earth (a two-dimensional manifold).
The sphere is one such example of a manifold which requires more than one chart to form an atlas; it requires two.
\begin{definition}[Atlas]\label{def:atlas}
    An atlas of a topological manifold $M$ is a collection of coordinate charts $\{(U_i, \varphi_i)\}_{i \in I}$ such that $\{U_i \subseteq M\}_{i \in I}$ is an open cover for $M$.
\end{definition}
Two charts may overlap, meaning that a point on a manifold may be represented by multiple charts.
We may compare two overlapping charts via a transition function, which composes one chart with the inverse of the other.
This is shown in \Cref{fig:transition_map}.

\begin{figure}[!htbp]
    \centering
\begin{tikzpicture}
    % Manifold (save path for reuse in crossing masks)
    \draw[semithick, smooth cycle, tension=0.4, fill=white, pattern color=gray, pattern=spray, save path=\pathManifold]
        plot coordinates{(2,2) (-0.5,0) (3,-2) (5,1)} node at (3,2.3) {$M$};

    %\draw[help lines] (-3,-6) grid (8,6);

    % Subsets (save paths for reuse in crossing masks and border redraws)
    \fill[smooth cycle, pattern color=red, pattern=spray, save path=\pathU]
        plot coordinates {(1,0) (1.5, 1.2) (2.5,1.3) (2.6, 0.4)}
        node [label={[label distance=-0.3cm, xshift=-2cm, fill=white]:$U$}] {};
    \fill[smooth cycle, pattern color=blue, pattern=spray, save path=\pathV]
        plot coordinates {(4, 0) (3.7, 0.8) (3.0, 1.2) (2.5, 1.2) (2.2, 0.8) (2.3, 0.5) (2.6, 0.3) (3.5, 0.0)}
        node [label={[label distance=-0.8cm, xshift=.75cm, yshift=1cm, fill=white]:$V$}] {};
    % Intersection of U and V (purple, matching transition map regions)
    \begin{scope}
        \clip[use path=\pathU];
        \fill[white, use path=\pathV];
        \fill[pattern color=green!50!black, pattern=spray, use path=\pathV];
    \end{scope}
    % Redraw subset borders on top of intersection and label backgrounds
    \draw[semithick, use path=\pathU];
    \draw[semithick, use path=\pathV];

    % Curved arrows between manifold and R^m
    \path[->] (0.8, 0) edge [bend right] node[left, xshift=-2mm] {$\varphi$} (-1, -2.9);
    \path[->] (-0.7, -3.05) edge [bend right] node [right, yshift=-3mm] {$\varphi^{-1}$} (1.093, -0.11);
    \path[->] (5.8, -2.8) edge [bend left] node[midway, xshift=-5mm, yshift=-3mm] {$\psi^{-1}$} (3.8, -0.35);
    \path[->] (4.2, 0) edge [bend left] node[right, xshift=2mm] {$\psi$} (6.2, -2.8);

    % Crossing masks: hide arrows where they cross manifold boundary
    \foreach \x/\y in {0.06/-0.57, 0.95/-1.2, 4/-1.1, 4.54/-0.12} {
        \begin{scope}
            \clip (\x,\y) circle (.15cm);
            \fill[white] (\x,\y) circle (.15cm);
            \draw[semithick, fill=white, pattern color=gray, pattern=spray, use path=\pathManifold];
            \draw[semithick, pattern color=red, pattern=spray, use path=\pathU];
            \draw[semithick, pattern color=blue, pattern=spray, use path=\pathV];
            \begin{scope}
                \clip[use path=\pathU];
                \fill[white, use path=\pathV];
                \fill[pattern color=green!50!black, pattern=spray, use path=\pathV];
            \end{scope}
        \end{scope}
    }

    % Axes
    \draw[thick, ->] (-3,-5) -- (0, -5) node [label=above:$\varphi(U)$] {};
    \draw[thick, ->] (-3,-5) -- (-3, -2) node [label=right:$\mathbb{R}^m$] {};
    \draw[->] (0, -3.85) -- node[midway, above]{$\psi \circ \varphi^{-1}$} (4.5, -3.85);
    \draw[thick, ->] (5, -5) -- (8, -5) node [label=above:$\psi(V)$] {};
    \draw[thick, ->] (5, -5) -- (5, -2) node [label=right:$\mathbb{R}^m$] {};

    % phi_U region in R^m (red with purple corner)
    \path[save path=\pathPhiU] plot [smooth cycle] coordinates{(-2, -4.5) (-2, -3.2) (-0.8, -3.2) (-0.8, -4.5)};
    \begin{scope}
        \clip[use path=\pathPhiU];
        \fill[white, pattern color=red, pattern=spray, use path=\pathPhiU];
        \fill[white] (-0.67, -3.06) -- +(180:0.8) arc (180:270:0.8) -- cycle;
        \fill[pattern color=green!50!black, pattern=spray] (-0.67, -3.06) -- +(180:0.8) arc (180:270:0.8) -- cycle;
    \end{scope}
    \draw[semithick, use path=\pathPhiU];
    \draw[semithick] (-1.45, -3.06) arc (180:270:0.8);

    % phi_V region in R^m (blue with purple corner)
    \path[save path=\pathPhiV] plot [smooth cycle] coordinates{(7, -4.5) (7, -3.2) (5.8, -3.2) (5.8, -4.5)};
    \begin{scope}
        \clip[use path=\pathPhiV];
        \fill[white, pattern color=blue, pattern=spray, use path=\pathPhiV];
        \fill[white] (5.7, -3.06) -- +(-90:0.8) arc (-90:0:0.8) -- cycle;
        \fill[pattern color=green!50!black, pattern=spray] (5.7, -3.06) -- +(-90:0.8) arc (-90:0:0.8) -- cycle;
    \end{scope}
    \draw[semithick, use path=\pathPhiV];
    \draw[semithick] (5.69, -3.85) arc (-90:0:0.8);
\end{tikzpicture}
\caption{Two charts along with their transition map}\label{fig:transition_map}
\end{figure}

\begin{definition}[Transition Function]\label{def:transition_function}
    Given charts $(U, \varphi)$ and $(V, \psi)$ of an atlas of a topological manifold $M$ such that $U \cap V \neq \varnothing$, the induced homeomorphism $\psi \circ \varphi^{-1} \colon \varphi(U \cap V) \to \psi(U \cap V)$ is called a transition function or map.
\end{definition}

Something to note is that our current definition of a topological manifold does not allow for a manifold to have a boundary.
There are objects that we would want to classify as manifolds, such as the unit closed disk $\mathbb{D}^2 = \{(x, y) \in \R^2 : x^2 + y^2 \leq 1\}$, that cannot be manifolds.
A point on the boundary of $\mathbb{D}^2$ does not have a valid chart, i.e. there is no open neighbourhood of the point that is homeomorphic to an open neighbourhood of $\R^2$.
We resolve this issue by separately defining topological manifolds with boundary.
This is done by considering the closed upper half-space $\mathbb{H}^n \coloneq \{(x^1, \ldots, x^n) \in \R^n : x^n \geq 0\}$ instead of $\R^n$.

We will also define the interior and boundary of the upper half-space.
Let $\Int \mathbb{H}^n \coloneq \{(x^1, \ldots, x^n) \in \R^n : x^n > 0\}$ be the interior of $\mathbb{H}^n$.
Let $\partial \mathbb{H}^n \coloneq \{(x^1, \ldots, x^n) \in \R^n : x^n = 0\}$ be the boundary of $\mathbb{H}^n$.

\begin{definition}[Topological Manifold with Boundary]\label{def:topological_manifold_with_boundary}
    A topological manifold with boundary $(M, \mathcal{T})$ is a Hausdorff, second countable topological space where every point $p \in M$ has an open neighbourhood $U$ homeomorphic to an open subset $V \subseteq \mathbb{H}^n$.
    If all $U$ are of dimension $n$ for some fixed $n \in \Z_{\geq 0}$, then $(M, \mathcal{T})$ is said to be $n$ dimensional.
    Given this, a coordinate chart $(U, \varphi)$ of $M$ is an open set $U \subseteq M$ with its corresponding homeomorphism $\varphi \colon U \to V$, where $V$ is the associated open subset of $\mathbb{H}^n$.
    It is a boundary chart if $\varphi(U) \cap \partial \mathbb{H}^n \neq \varnothing$.
\end{definition}
As every open subset of $\R^n$ is homeomorphic to an open subset of $\mathbb{H}^n$ via $(x^1, \ldots, x^n) \mapsto (x^1, \ldots, x^{n-1}, e^{x^n})$, all topological manifolds are manifolds with boundary, though their boundary may be empty.
We will precisely define what the boundary, as well as the interior, of a manifold actually is.
\begin{definition}[Interior and Boundary]\label{def:interior_and_boundary}
    An interior chart $(U, \varphi)$ of a topological manifold $M$ with boundary is a chart of $M$ such that $\varphi(U)$ is an open subset of $\R^n$.
    A boundary chart $(U, \varphi)$ of $M$ is a chart of $M$ such that $\varphi(U)$ is an open subset of $\mathbb{H}^n$ and not an open subset of $\R^n$, equivalently $\varphi(U) \cap \partial \mathbb{H}^n \neq \varnothing$.
    
    A point $p \in M$ is an interior point if it is in the domain of an interior chart, and it is a boundary point if it is in the domain of a boundary chart such that $\varphi(p) \in \partial \mathbb{H}^n$.
    The set of all interior points of $M$ is called the interior of $M$, denoted by $\Int M$.
    The set of all boundary points of $M$ is called the boundary of $M$, denoted by $\partial M$.
\end{definition}
\begin{remark}
    If $p \in M$ is an interior point of a topological manifold $M$ with boundary, then it may be in the domain of a boundary chart $(U, \varphi)$ such that $\varphi(p) \notin \partial \mathbb{H}^n$.
    In this case, the restriction of $\varphi$ to $\varphi^{-1}(\Int \mathbb{H}^n)$ is an interior chart whose domain contains $p$.
\end{remark}
For the purpose of distinction, we may refer to manifolds with boundary whose boundary is empty, i.e. $\partial M = \varnothing$, as manifolds without boundary -- in spite of the fact that the term \enquote*{topological manifold} already describes such manifolds.
Furthermore, we may refer to a topological manifold with boundary, whose boundary might be empty, as a manifold with or without boundary.
We often say \enquote*{manifold} to refer to a topological manifold with or without boundary.

A point $p \in M$ is either an interior point or a boundary point.
Furthermore, it cannot be an interior point with respect to one chart and a boundary point with respect to another.
\begin{proposition}\label{prop:topological_invariance_of_the_boundary}
    Given a topological manifold $M$ with boundary, a point $p \in M$ is either an interior point or a boundary point, but it cannot be both, i.e. $\Int M \cap \partial M = \varnothing$ and $\Int M \cup \partial M = M$.
\end{proposition}
\begin{proof}
    This is the topological invariance of the boundary, which rests on invariance of domain; see \cite[Thm.~1.37]{Lee2012}.
\end{proof}
As topological manifolds with boundary are technically distinct objects from topological manifolds, and their distinction provides them with added complexities, we will have to mention them separately for some definitions.
This includes the definition of smooth manifolds.
We may omit their discussion for some concepts as well, with the extension to manifolds with boundary being implicit.

\subsection{Smooth Manifolds}\label{subsec:manifolds:smooth_manifolds}

By themselves, atlases are not particularly useful or significant to the study of topological manifolds.
Their usefulness comes when we wish to have manifolds with additional structure, as such structure is defined on the charts of an atlas.
We wish to introduce the concept of differential structure, which we do via defining special atlases.
To greatly simplify the following definitions, we will define both the space of continuously differentiable functions, and the space of (not necessarily continuously) differentiable functions.
The latter is unlikely to be used.

Let $D^k(A, B)$ be the set of $k$-times differentiable functions with domain $A$ and codomain $B$.
Let $C^k(A, B)$ be the set of $k$-times continuously differentiable functions with domain $A$ and codomain $B$, i.e. functions in $D^k(A, B)$ whose $k$-th derivative is continuous.
Lastly, let $C^\infty(A, B)$ be the set of all smooth (infinitely differentiable) functions with domain $A$ and codomain $B$.
We may omit the domain or the codomain to refer to the space of such functions irrespective of the removed detail.
We may omit both and refer to $D^k$, $C^k$, or $C^\infty$ as spaces of their respective functions with any domain.
    
Also, we may say that a function $f$ is $C^k$ to denote that it is in the space $C^k$.
Objects that are not functions, but are defined with relation to functions belonging to spaces mentioned prior, may be described using those spaces.
In this way, the notation for spaces is used as an adjective for related mathematical spaces, as well as being actual spaces of functions.

For what follows, unless otherwise stated, the term \enquote*{differentiable} attached to another term will denote continuous differentiability specifically.
\begin{definition}[Compatible Charts]\label{def:compatible_charts}
    Two charts $(U, \varphi)$ $(V, \psi)$ are $C^k$ compatible if $U \cap V = \varnothing$ or both transition functions $\psi \circ \varphi^{-1}$ and $\varphi \circ \psi^{-1}$ are $C^k$.
    If the aforementioned transition functions are instead $C^\infty$, then the two charts are smoothly compatible.
\end{definition}
\begin{definition}[Differentiable and Smooth Atlas]\label{def:differentiable_and_smooth_atlas}
    An atlas $\mathcal{A} = \{(U_i, \varphi_i)\}_{i \in I}$ of a topological manifold $M$ is $k$-times differentiable ($C^k$) if any two charts are $C^k$ compatible.
    If any two charts are smoothly compatible, then $\mathcal{A}$ is called a smooth atlas.
\end{definition}
Whilst not strictly necessary for the definition of differential structure, we will define maximal atlases.
\begin{definition}[Maximal Atlas]\label{def:maximal_atlas}
    A $C^k$ (or smooth) atlas $\mathcal{A}$ of a topological manifold $M$ is maximal if it is not contained within any other $C^k$ (or smooth) atlas.
\end{definition}
We could define differentiable structure via equivalence classes of differentiable atlases -- where equivalence of atlases means their union is a differentiable atlas -- but there is a natural isomorphism between those classes and maximal differentiable atlases.
Hence, we are able to use the latter for a more straightforward definition.
Any chart that is $C^k$ or smoothly compatible with every chart in a $C^k$ or smooth maximal atlas $\mathcal{A}$ will be in $\mathcal{A}$.
Otherwise, we could create a new smooth atlas which contains $\mathcal{A}$. 
\begin{definition}[Differential and Smooth Structure]\label{def:differential_and_smooth_structure}
    A $k$-differential ($C^k$) structure on a topological manifold $M$ is a $C^k$ maximal atlas $\mathcal{A}$ on $M$.
    A smooth structure on $M$ is a smooth maximal atlas $\mathcal{A}$ on $M$.
    Given a $C^k$ or smooth maximal atlas $\mathcal{A}$, we say that $M$ has $C^k$ or smooth structure.
\end{definition}
Every $C^k$ or smooth atlas $\mathcal{A}$ on a manifold $M$ is contained within a unique $C^k$ or smooth maximal atlas $\mathcal{A}_{\text{max}}$ on $M$ \cite[Prop.~1.17]{Lee2012}, which is stated there for smooth atlases; the $C^k$ case is identical.
This maximal atlas may be called the structure determined by $\mathcal{A}$.
The aforementioned equivalence class of atlases are those which determine the same structure on $M$.
These two facts means that we only need to specify some atlases in order to determine the structure on $M$.

We will define both differentiable and smooth manifolds, though it is the latter that we will mainly focus on.
The two are very similar concepts, however.
\begin{definition}[Differentiable and Smooth Manifold]\label{def:differentiable_and_smooth_manifold}
    A $k$-times, or $k$-fold, differentiable ($C^k$) manifold $(M, \mathcal{T}, \mathcal{A})$ is a topological manifold $(M, \mathcal{T})$ with $C^k$ structure.
    If the manifold instead has smooth structure, it is called a smooth manifold.
\end{definition}
\begin{remark}
    We could have defined manifolds, including those with extra structure, in a different way.
    Instead of starting with a topological space, we could have had a regular set which we define charts and atlases on.
    Charts and atlases would have a slightly different definition, but then a manifold could be defined as a set and an atlas.
    The atlas would induce the topology.
\end{remark}
A similar definition may be provided to define differentiable and smooth manifolds with boundary.
The reason this does not immediately follow from the previous definition is that the notion of smooth function between open subsets of $\R^n$ is not sufficient.
Thus, a definition of smooth functions whose domain is an open subset of $\mathbb{H}^n$ is required.
All other concepts are completely analogous, so they are all compiled into the following definition.
\begin{definition}[Differentiable and Smooth Manifold with Boundary]
    A smooth manifold with boundary $(M, \mathcal{T}, \mathcal{A})$ is a topological manifold with boundary which has smooth structure.
    Here, smooth structure is a maximal smooth atlas $\mathcal{A}$ on $M$, with smoothness being given by smooth compatibility of charts of $M$, which is defined by the smoothness of the transition functions.
    
    Let $U$ be an open subset of $\mathbb{H}^n$ and let $\widetilde{U}$ be an open subset of $\R^n$ containing a point $p \in U$.
    A function $f \colon U \to \R^m$ is smooth at $p$ if there exists such an open set $\widetilde{U}$ and a smooth (by the standard definition) function $\tilde{f} \colon \widetilde{U} \to \R^m$ such that $f(x) = \tilde{f}(x)$ for all $x \in U \cap \widetilde{U}$.
    The smoothness of $f$ on a subset is defined accordingly, and $f$ is a smooth function if it is smooth on its domain.
\end{definition}
Below is a lemma created to be a straightforward way of showing that a given set with structure is indeed a smooth manifold.
It is really a compilation and condensation of the information already given.
\begin{lemma}\label{lem:smooth_manifold_chart_lemma}
    Let $M$ be a set, and suppose we have a collection $\{U_i\}_{i \in I}$ of subsets of $M$ with maps $\varphi_i \colon U_i \to \R^n$ that satisfy the following properties.
    \begin{enumerate}
        \item For each $i$, $\varphi_i$ is a bijection between $U_i$ and an open subset $\varphi_i(U_i) \subseteq \R^n$.
        \item For each $i$ and $j$, the sets $\varphi_i(U_i \cap U_j)$ and $\varphi_j(U_i \cap U_j)$ are open in $\R^n$.
        \item If $U_i \cap U_j \neq \varnothing$, then the map $\varphi_j \circ \varphi^{-1}_i \colon \varphi_i(U_i \cap U_j) \to \varphi_j(U_i \cap U_j)$ is smooth.
        \item Countably many of the sets $U_i$ cover $M$.
        \item If $p, q \in M$ are distinct points, then there exists a $U_i$ containing $p$ and $q$ or there exist disjoint sets $U_i \ni p$ and $U_j \ni q$.
    \end{enumerate}
    Then, $M$ forms a unique second-countable smooth manifold with each $(U_i, \varphi_i)$ being a smooth chart.
\end{lemma}
\begin{proof}
    The proof of this follows the proof of \cite[Lem.~1.35]{Lee2012}.
    We define the topology of $M$ using the sets of the form $\varphi^{-1}_i(V)$, where $V \subseteq \R^n$ is open, as a basis.
    Given two basis sets $\varphi^{-1}_i(V)$ and $\varphi^{-1}_j(W)$, we wish to show that for any point $p \in \varphi^{-1}_i(V) \cap \varphi^{-1}_j(W)$, there exists a third basis set containing $p$ contained within $\varphi^{-1}_i(V) \cap \varphi^{-1}_j(W)$.
    It suffices to show that $\varphi^{-1}_i(V) \cap \varphi^{-1}_j(W)$ is a basis set.
    
    Property 3 states that $\varphi_j \circ \varphi^{-1}_i$ is smooth, so it is continuous.
    Its preimage $(\varphi_j \circ \varphi^{-1}_i)^{-1}(W)$ is thus an open subset of $\varphi_i(U_i \cap U_j)$. 
    Property 2 implies that the domain $\varphi_i(U_i \cap U_j)$ is open in $\R^n$, so the preimage is also open in $\R^n$.
    Thus,
    \begin{equation*}
        \varphi^{-1}_i(V) \cap \varphi^{-1}_j(W) = \varphi^{-1}_i(V \cap (\varphi_j \circ \varphi^{-1}_i)^{-1}(W))
    \end{equation*}
    is also a basis set.

    Each $\varphi_i$ is a homeomorphism onto $\varphi_i(U_i)$, since it maps the basis sets contained in $U_i$ onto the open subsets of $\varphi_i(U_i)$, so $M$ is locally Euclidean with dimension $n$.
    The space is second countable via property 4, and it is Hausdorff via property 5.
    This makes $M$ a topological manifold.
    Property 3 means that $\{U_i, \varphi_i\}_{i \in I}$ forms a smooth atlas.
    Therefore, the set $M$ with its topology and smooth structure is a second-countable smooth manifold.
\end{proof}
We mentioned before that charts specify a coordinate system, though now we will make that precise by defining local coordinates.
Indeed, their use comes with the smooth structure of a manifold.
\begin{definition}[Local Coordinates]\label{def:local_coordinates}
    Given an $n$-dimensional topological manifold $M$ and a chart $(U, \varphi)$, a local coordinate of the chart is a function $x^i \colon U \to \R$ given by $x^i = \proj_i \circ \varphi$, where $\proj_i \colon \R^n \to \R$ is the projection to the $i$-th component relative to the standard basis of $\R^n$.
    We often write $u^i = \proj_i$ and call it a Euclidean coordinate function.
    If the chart is within a smooth atlas of a smooth manifold, then $x^i$ may be called a smooth local coordinate.
    The collection of local coordinates forms a local coordinate system for $M$, and we may write
    \begin{equation*}
        \varphi = (x^1, \ldots, x^n)
    \end{equation*}
\end{definition}
Any point $p \in U$ can be represented using these local coordinates, as one would expect: $(x^1(p), \ldots, x^n(p)) \in \R^n$.
With this, and given a function $f \colon M \to \R$, we may write $f(x^1, \ldots, x^n)$ to mean $f(p)$.
This is an abuse of notation, as $\varphi(p) = (x^1, \ldots, x^n) \in \varphi(U) \subseteq \R^n$.
It is really an identification of $\eval*{f}{U}$ (the restriction of $f$ to $U$) with $f \circ \varphi^{-1} \colon \varphi(U) \to \R$ and a subsequent abuse of notation.
As it is convenient, however, we will consider $p \in M$ to be the actual Euclidean vector $(x^1, \ldots, x^n)$ at times when it is unambiguous.
\[\begin{tikzcd}[sep=large]
	U & M & \R \\
	{\varphi(U)}
	\arrow["\iota", hook, from=1-1, to=1-2]
	\arrow["\varphi"', from=1-1, to=2-1]
	\arrow["f", from=1-2, to=1-3]
	\arrow["f \circ \varphi^{-1}"', from=2-1, to=1-3]
\end{tikzcd}\]
The most important thing we can now do with manifolds is to define (differentiable and) smooth functions.
These are functions which allow us to perform calculus.
\begin{definition}[Differentiable and Smooth Function]\label{def:differentiable_and_smooth_function}
    A function $f \colon M \to N$ between smooth manifolds $M$ and $N$ is smooth at a point $p \in M$ if there are charts $(U, \varphi)$ for $M$ and $(V, \psi)$ for $N$ with $p \in U$ and $f(U) \subseteq V$ such that $\psi \circ f \circ \varphi^{-1} \colon \varphi(U) \to \psi(V)$ is smooth at $\varphi(p)$.
    If $f$ is smooth at all points $p \in S \subseteq M$, it is smooth on $S$.
    If $f$ is smooth at all points $p \in M$, it is a smooth function.
    Then, $k$-times differentiable ($C^k$) functions are defined analogously.
\[\begin{tikzcd}[sep=large]
	M && N \\
	{\varphi(U)} && {\psi(V})
	\arrow["f", from=1-1, to=1-3]
	\arrow["\varphi"', from=1-1, to=2-1]
	\arrow["\psi", from=1-3, to=2-3]
	\arrow["{\psi \circ f \circ \varphi^{-1}}"', from=2-1, to=2-3]
\end{tikzcd}\]
\end{definition}
Just as we have the idea of a homeomorphism as an isomorphism of topological manifolds, which is a more suitable concept than a plain bijection, we have the concept of a diffeomorphism for differentiable manifolds.
These are structure-preserving bijections, with the structure they preserve being differential or smooth.
Again, the one for smooth structures is what will be primarily utilized.
\begin{definition}[Diffeomorphism]\label{def:diffeomorphism}
    A function $f \colon M \to N$ between $C^k$ manifolds is a $C^k$ diffeomorphism if it is a $C^k$ bijection with a $C^k$ inverse.
    A function $f \colon M \to N$ between smooth manifolds is a smooth diffeomorphism if it is a smooth bijection with a smooth inverse.
    In either case, if such a function exists, the two manifolds are said to be diffeomorphic, denoted $M \cong N$.
\end{definition}
We have made sufficient effort to introduce basic additional structure on topological manifolds in a general sense, defining both $C^k$ and smooth manifolds.
Henceforth, we will almost exclusively consider the latter (if not one besides the two).
For anything that follows, unless otherwise stated, we consider manifolds to be smooth, and each component has dimension $n$ for some fixed $n \in \Z_{\geq 0}$.
If a second manifold is considered, it has dimension $m$ for some fixed $m \in \Z_{\geq 0}$.

\subsection{Tangent Spaces}\label{subsec:manifolds:tangent_spaces}

\subsubsection{Germs of Functions}\label{subsubsec:manifolds:tangent_spaces:germs_of_functions}

With smooth manifolds, and smooth functions defined on such manifolds, we may perform calculus.
An introduction to single-variable calculus or multivariable calculus will first present the concept of limits.
Limits can then be used, among other things, to define the derivative.
Introductory multivariable calculus really has two main derivatives which are the directional derivative and the Jacobian.
Given a smooth function $f \colon U \to \R^m$, where $U \subseteq \R^n$ is open, the directional derivative of $f$ along a vector $\mathbf{v} \in \R^n$ at a point $p$ is given by
\begin{equation}
    \nabla_\mathbf{v}f(p) \coloneq \lim_{h \to 0} \frac{f(p + h\mathbf{v}) - f(p)}{h}.
\end{equation}
The directional derivative satisfies a rule very similar to the product rule for regular derivatives, $\nabla_\mathbf{v}(fg) = (\nabla_\mathbf{v}f) g + f \nabla_\mathbf{v}g$, which we call the Leibniz rule or Leibniz identity.
This, along with the fact that it is a linear map, makes the directional derivative operator $\nabla_\mathbf{v}$ a derivation.
\begin{definition}[Derivation]\label{def:derivation}
    Given a nonassociative algebra $A$ over a commutative unital ring $R$, an $R$-linear map $D \colon A \to A$ is a \emph{derivation of $A$} if for all $a, b \in A$, it satisfies the \emph{Leibniz rule}.
    \begin{equation}
        D(ab) = D(a)b + aD(b)
    \end{equation}
    If the elements of $A$ are maps on a set $X$ to the ring $R$, then $D_x \colon A \to R$ is a derivation of $A$ at $x \in X$ if it is $R$-linear and satisfies
    \begin{equation}
        D_x(ab) = D_x(a)b(x) + a(x)D_x(b)
    \end{equation}
    for all $a, b \in A$.
    We may associate these two definitions via $D_x(a) = D(a)(x)$.
\end{definition}
\begin{example}\label{ex:derivation}
    Given the space $C^\infty(\R, \R^m)$ over the field $\R$, differentiation of functions is a derivation.
    Derivatives are linear, and the Leibniz rule manifests itself as the product rule.
    Given the space $C^\infty(U, \R^m)$ over the field $\R$, where $U \subseteq \R^n$ is open, the operator which takes the directional derivative along some vector in $\R^n$ inherits linearity and satisfaction of the Leibniz rule, making it a derivation.
\end{example}
With directional derivatives, we can introduce calculus on manifolds.
This may be done by considering the equivalence class of curves (smooth maps from a nonempty interval of $\R$ to $M$) related by their directional derivatives and value at zero.
However, there is no natural vector space structure on such a class.
This is resolved by considering the property that directional derivatives possess instead of the derivatives themselves.
That is, we wish to consider derivations.
Before we may do that, we will define a concept that allows us to consider functions locally.

Given two smooth functions $f \in C^\infty(U, \R)$ and $g \in C^\infty(V, \R)$, where $U$ and $V$ are open neighbourhoods of $p \in M$, the two functions may agree everywhere in an open neighbourhood $W \subseteq U \cap V$ of $p$.
Then, the two functions can be considered to be essentially the same on that open set, so they have the same local properties.
We could find the Taylor series for both functions and see that they would be the same within $W$.
From this, it is reasonable to define an equivalence relation for all such functions, given by
\begin{equation*}
    (U, f) \sim (V, g) \quad \iff \quad f(x) = g(x) \quad \forall x \in W.
\end{equation*}
With this defined notion of two functions being the same, we can consider the collection of all functions that are the same as $f$, i.e. we have the equivalence class
\begin{equation*}
    [(U, f)]_x = \{g \in C^\infty(V, \R) : (U, f) \sim (V, g)\}.
\end{equation*}
All functions in this equivalence class share all local properties with each other, such as continuity and differentiability.
Thus, these properties of a function can be studied by studying the equivalence class instead.
This is the motivation for the definition of germs of functions.
\begin{definition}[Germ of Smooth Function]\label{def:germ_of_smooth_function}
    Let M be a smooth manifold, and let $p \in M$ be a point.
    Consider pairs $(U, f)$, where $U \subseteq M$ is an open neighbourhood of $p$ and $f \colon U \to \R$ is a smooth function.
    The germ of $f$ at $p$ is the equivalence class
    \begin{equation}
    [(U, f)]_p = \{g \in C^\infty(V, \R) : (U, f) \sim (V, g)\},
    \end{equation}
    with the equivalence relation given by
    \begin{equation}
    (U, f) \sim (V, g) \quad \iff \quad f(x) = g(x) \quad \forall x \in W,
    \end{equation}
    where the relation holds if there is some open neighbourhood $W \subseteq U \cap V$ of $p$ on which $f$ and $g$ agree.
    We may write $[U, f]_p$ as $[f]_p$, which is justified as $[U', \eval*{f}{U'}]_p = [(U, f)]$ for an open neighbourhood $U' \subseteq U$ of $p$.
    When unambiguous, we may even write it as $[f]$.
    The set of germs of smooth real-valued functions on $M$ at $p$ is denoted by $\mathcal{C}^\infty_p(M, \R)$.
\end{definition}
It follows that $\mathcal{C}^\infty_p(M, \R)$ forms an associative algebra over $\R$, and thus a vector space over $\R$.
\begin{proposition}\label{prop:germs_form_associative_algebra}
    Let $M$ be a smooth manifold and $p \in M$ be a point.
    The set $\mathcal{C}^\infty_p (M, \R)$ forms an associative algebra over the ring $\R$ under the following operations.
    Note that $\lambda \in \R$, $f \in C^\infty(U, \R)$, and $g \in C^\infty(V, \R)$, where $U, V \subseteq M$ are open and $p \in U \cap V$.
    \begin{align}
        \begin{split}
        \lambda [(U, f)]_p &= [(U, \lambda f)]_p \\
        [(U, f)]_p + [(V, g)]_p &= [(U \cap V, \eval{f}{U \cap V} + \eval{g}{U \cap V})]_p \\
        [(U, f)]_p \cdot [(V, g)]_p &= [(U \cap V, \eval{f}{U \cap V} \cdot \eval{g}{U \cap V})]_p
        \end{split}
    \end{align}
\end{proposition}
\begin{proof}
    The operations are well defined on germs: if $f$ and $\tilde{f}$ agree on a neighbourhood $W$ of $p$, and $g$ and $\tilde{g}$ agree on a neighbourhood $W'$ of $p$, then $\lambda f + \mu g$ and $fg$ agree with $\lambda \tilde{f} + \mu \tilde{g}$ and $\tilde{f}\tilde{g}$ on $W \cap W'$.
    Each algebra axiom involves finitely many germs, and it holds on the intersection of their domains because it holds pointwise there.
    The germ of the constant function $1$ is the unit.
\end{proof}
As stated before, we capture the idea of directional derivatives using derivations.
However, derivatives are also local objects, and we can capture the idea of local objects using germs.
Hence, we can combine these two to obtain a definition of tangent vectors and the tangent space of a manifold at a point.
\begin{definition}[Tangent Space]\label{def:tangent_space}
    Given a smooth manifold $M$ and a point $p \in M$, the tangent space of $M$ at $p$ is the set of all derivations of $\mathcal{C}^\infty_p(M, \R)$ at $p$, that is, the linear maps $v \colon \mathcal{C}^\infty_p(M, \R) \to \R$ satisfying the Leibniz rule below.
    The tangent space is denoted by $T_p M$, and an element in $T_p M$ is called a tangent vector.
    A tangent vector $v \in T_p M$ acts on germs $[f]_p$, though we define it to also act on representative functions $f$ of germs $[f]_p$ via $v(f) = v[f]_p$.
    Being a derivation, $v$ acts on smooth functions $f$ and $g$ as follows.
    \begin{equation}
        v(fg) = v(f)g(p) + f(p)v(g)
    \end{equation}
\end{definition}
Thus, a tangent vector is a map that takes a smooth function to a real number.
Intuitively, a tangent vector $v$ says how much some function $f$ changes in a particular direction.
Calculus is built on the concept of linearization, with functions -- and more generally entire spaces -- being flat locally.
That is what tangent vectors help to do.
For tangent vectors to be vectors, the tangent space must be a vector space.
\begin{proposition}\label{prop:tangent_space_is_vector_space}
    Let $M$ be a smooth manifold and $p \in M$ be a point.
    The tangent space of $M$ at $p$ is a vector subspace of the dual space $\mathcal{C}^\infty_p(M, \R)^*$.
\end{proposition}
\begin{proof}
    The dual space of $\mathcal{C}^\infty_p(M, \R)$ is a vector space over $\R$, so it suffices to show that $T_p M$ is closed under linear combinations.
    Given tangent vectors $u, v \in T_p M$ and scalars $\lambda, \mu$, it holds that $w = \lambda u + \mu v$ is in $T_p M$.
    Being a linear combination of linear maps, $w$ is linear.
    Also, $w$ satisfies the Leibniz rule; consider functions $f, g \in C^\infty(U, \R)$.
    \begin{align*}
        w(fg) &= (\lambda u + \mu v)(fg) \\
        &= \lambda u(fg) + \mu v(fg) \\
        &= \lambda u(f) g(p) + \lambda f(p) u(g) + \mu v(f) g(p) + \mu f(p) v(g) \\
        &= (\lambda u + \mu v)(f) g(p) + f(p) (\lambda u + \mu v)(g) \\
        &= w(f) g(p) + f(p) w(g) \qedhere
    \end{align*}
\end{proof}
This is not the only way to define tangent spaces.
Besides velocities of curves mentioned before, we can consider a tangent space to be the set of derivations, where a derivation is a function -- instead of an equivalence class of functions.
By not using germs, it is a less abstract definition, but the fact that tangent vectors are local objects is not immediately noticeable.
The two are related via an isomorphism.
\begin{proposition}\label{prop:tangent_space_definition_equivalence}
    Let $M$ be a smooth manifold, and let $p \in M$ be a point.
    Let $D_p M$ be space of derivations of $C^\infty(M, \R)$ at $p$.
    The map $\Phi \colon T_p M \to D_p M$ defined by $(\Phi v) f = v([f]_p)$ is an isomorphism (linear bijection) between the two vector spaces, so $T_p M \cong D_p M$.
\end{proposition}
\begin{proof}[Sketch]
    The map $\Phi$ is linear, and $\Phi v$ is a derivation at $p$ because $v$ is.
    It is injective: suppose that $\Phi v = 0$, and let $g \in C^\infty(U)$ represent a germ at $p$.
    Choose a smooth bump function $\psi \in C^\infty(M)$ that equals $1$ near $p$ and has support in $U$, which exists by the existence of partitions of unity \cite[Thm.~2.23]{Lee2012}.
    Then $\psi g$ extends by zero to an element of $C^\infty(M)$ with $[\psi g]_p = [g]_p$, so $v[g]_p = (\Phi v)(\psi g) = 0$.
    It is surjective: for $w \in D_p M$, define $v[g]_p = w(\psi g)$ with $\psi$ and $g$ as above.
    This is independent of the choice of $\psi$ and of the representative $g$, because $w$ vanishes on any $h \in C^\infty(M)$ that is zero near $p$: choosing a bump function $\psi_0$ that equals $1$ near $p$ and has support where $h$ vanishes, we have $h = (1 - \psi_0)h$, so $w(h) = w(1 - \psi_0)h(p) + (1 - \psi_0)(p)w(h) = 0$.
    The resulting $v$ is a derivation of germs with $\Phi v = w$.
\end{proof}
The tangent space of a manifold at a point may be thought of as the space that the manifold resembles at that point.
This follows from the characterization of a manifold as locally Euclidean.
The tangent space of a circle may be thought of as the line literally tangent to the circle.
This cannot be a way to think of all tangent spaces, as it assumes an ambient space in which the manifold is embedded.
Hence, it must remain as -- still valuable -- intuition.

As tangent spaces are vector spaces, they have a basis.
Given that $M$ is $n$ dimensional, we expect $T_p M$ to be $n$ dimensional.
This will be shown by demonstrating that $T_p M$ has a basis of $n$ elements.
Just as a Euclidean vector may be expressed as a linear combination of Euclidean basis vectors, a directional derivative may be expressed as the linear combination of partial derivatives.
Partial derivatives are directional derivatives along a basis vector.
We extend this idea to use partial derivatives to form a basis for $T_p M$.
\begin{definition}[Partial Derivative of Germ]\label{def:partial_derivative_of_germ}
    Let $M$ be an $n$-dimensional smooth manifold, $p \in M$ be a point, and $(U, \varphi)$ be a chart containing $p$.
    Let $\varphi = (x^1, \ldots, x^n)$, where $x^i \colon U \to \R$ is a local coordinate of $U$.
    The partial derivative with respect to $x^i$ evaluated at $p$, denoted by $\eval*{\partial/\partial x^i}{p}$ or $\eval*{\partial_i}{p}$, is an operator of germs in $\mathcal{C}^\infty_p(M, \R)$ defined by
    \begin{equation}
        \eval{\frac{\partial}{\partial x^i}}{p} [f]_p = \frac{\partial(f \circ \varphi^{-1})}{\partial u^i}(\varphi(p)),
    \end{equation}
    where $u^i$ are Euclidean coordinates and $f$ is a representative of the germ $[f]_p$ defined on an open neighbourhood of $p$ contained in $U$.
    We may also denote $\eval*{\partial/\partial x^i}{p} [f]_p$ as $\eval*{\partial_i}{p} [f]_p$.
    If appropriate and unambiguous, reference to the point $p$ may be omitted from the notation; we may write $\partial_i [f]_p$, $\eval*{\partial_i}{p}f$, or $\partial_i f$.
\end{definition}
The partial derivative operator $\eval*{\partial_i}{p}$ is a derivation, which follows from its nature as a derivative, meaning that it is an element of $T_p M$.
Furthermore, we have the following orthogonality relations.
\begin{equation}\label{eq:tangent_space_orthogonality_relations}
    \eval{\frac{\partial}{\partial x^i}}{p} [x^j]_p = \frac{\partial u^j}{\partial u^i}(\varphi(p)) = \delta^i_j = \begin{cases}
        1 & \text{if } i = j \\
        0 & \text{if } i \neq j
    \end{cases}
\end{equation}
It follows from this that
\begin{equation*}
    \left\{\eval{\frac{\partial}{\partial x^1}}{p}, \ldots, \eval{\frac{\partial}{\partial x^n}}{p}\right\}
\end{equation*}
is linearly independent.
That this set also spans $T_p M$ is shown in \cref{prop:tangent_space_basis_and_dimension}, after the following two lemmas.
\begin{lemma}\label{lem:tangent_vector_on_constant}
    Let $M$ be a smooth manifold, and let $p \in M$ be a point.
    A tangent vector $v \in T_p M$ acts on a constant to give zero, i.e. $v(\lambda) = 0$ for all $\lambda \in \R$.
\end{lemma}
\begin{proof}
    We may write $\lambda = \lambda \cdot 1$ and $1 = 1 \cdot 1$.
    Applying $v$ to the latter, using the Leibniz rule, gives
    \begin{equation*}
        v(1) = v(1 \cdot 1) = v(1) \cdot 1 + 1 \cdot v(1) = 2v(1),
    \end{equation*}
    which means that
    \begin{equation*}
        v(1) = 0.
    \end{equation*}
    Consider this, along with the linearity of $v$.
    \begin{equation*}
        v(\lambda) = v(\lambda \cdot 1) = \lambda v(1) = \lambda \cdot 0 = 0 \qedhere
    \end{equation*}
\end{proof}
\begin{lemma}[Hadamard's Lemma]\label{lem:hadamard's_lemma}
    Let $M$ be a smooth manifold with a coordinate chart $(U, \varphi)$ such that $\varphi(U)$ is star-convex about $\varphi(p)$, and let $p \in U$ be a point.
    Let $f \in C^\infty(U, \R)$.
    Then, there exist $g_i \in C^\infty(U, \R)$ such that for $q \in U$,
    \begin{equation}
        f(q) = f(p) + \sum_{i=1}^n (x^i(q) - x^i(p)) g_i(q).
    \end{equation}
    Furthermore,
    \begin{equation}
        g_i(p) = \frac{\partial (f \circ \varphi^{-1})}{\partial u^i}(\varphi(p)).
    \end{equation}
\end{lemma}
\begin{proof}
    Write $a = \varphi(p)$ and $\hat{f} = f \circ \varphi^{-1} \colon \varphi(U) \to \R$.
    For $y \in \varphi(U)$ the segment from $a$ to $y$ lies in $\varphi(U)$, so the fundamental theorem of calculus and the chain rule give
    \begin{equation*}
        \hat{f}(y) - \hat{f}(a) = \int_0^1 \frac{d}{dt} \hat{f}(a + t(y - a)) \d{t} = \sum_{i=1}^n (y^i - a^i) \int_0^1 \frac{\partial \hat{f}}{\partial u^i}(a + t(y - a)) \d{t}.
    \end{equation*}
    The functions $\hat{g}_i(y) = \int_0^1 \frac{\partial \hat{f}}{\partial u^i}(a + t(y - a)) \d{t}$ are smooth by differentiation under the integral sign, and $g_i = \hat{g}_i \circ \varphi$ satisfy the identity with $y = \varphi(q)$.
    Finally, $\hat{g}_i(a) = \partial \hat{f}/\partial u^i (a)$, which gives the value of $g_i(p)$.
\end{proof}
\begin{proposition}\label{prop:tangent_space_basis_and_dimension}
    Let $M$ be a smooth manifold, and let $p \in M$ be a point.
    Let $(U, \varphi)$ be a chart of $M$ containing $p$ with $\varphi = (x^1, \ldots, x^n)$.
    The set
    \begin{equation}\label{eq:tangent_space_basis}
        \left\{\eval{\frac{\partial}{\partial x^1}}{p}, \ldots, \eval{\frac{\partial}{\partial x^n}}{p}\right\}
    \end{equation}
    forms a basis for $T_p M$.
\end{proposition}
\begin{proof}
    Since tangent vectors act on germs, we may shrink $U$ so that $\varphi(U)$ is an open ball about $\varphi(p)$, which is star-convex about $\varphi(p)$.
    The candidate set \eqref{eq:tangent_space_basis} is linearly independent by \eqref{eq:tangent_space_orthogonality_relations}.
    To see this, consider $\lambda^1, \ldots, \lambda^n \in \R$ such that $\lambda^i\partial_i = 0$ (by symmetry, $\delta^i_j = \delta^j_i$).
    Evaluating both sides at $[x^j]$ gives $\lambda^i\delta^j_i = 0$.
    This implies that $\lambda^i = 0$.
    We now need to show that, for $v \in T_p M$,
    \begin{equation}\label{eq:tangent_vector_basis_sum}
        v = \sum_{i=1}^n v[x^i] \eval*{\partial_i}{p}.
    \end{equation}
    For this, we evaluate both sides of \eqref{eq:tangent_vector_basis_sum} at an arbitrary germ $[f]_p$.
    By Hadamard's lemma (\Cref{lem:hadamard's_lemma}),
    \begin{equation*}
        f = f(p) + \sum_{i=1}^n(x^i - x^i(p))g_i,
    \end{equation*}
    where
    \begin{equation*}
        g_i(p) = \frac{\partial (f \circ \varphi^{-1})}{\partial u^i}(\varphi(p)).
    \end{equation*}
    Then, making use of \Cref{lem:tangent_vector_on_constant}, we have
    \begin{align*}
        v[f] &= v\left[f(p) + \sum_{i=1}^n (x^i - x^i(p)) g_i\right] \\
        &= v[f(p)] + \sum_{i=1}^n v[(x^i - x^i(p)) g_i] \\
        &= 0 + \sum_{i=1}^n \left(v[x^i - x^i(p)] g_i(p) + (x^i(p) - x^i(p))v[g_i]\right) \\
        &= \sum_{i=1}^n v[x^i - x^i(p)]g_i(p) \\
        &= \sum_{i=1}^n v[x^i] g_i(p) \\
        &= \sum_{i=1}^n v[x^i] \frac{\partial (f \circ \varphi^{-1})}{\partial u^i}(\varphi(p)) \\
        &= \sum_{i=1}^n v[x^i] \eval{\partial_i}{p} [f] \\
        &= \left(\sum_{i=1}^n v[x^i] \eval{\partial_i}{p}\right) [f].
    \end{align*}
    As this is true for arbitrary germs $[f]$, we have
    \begin{equation*}
        v = \sum_{i=1}^n v[x^i] \eval*{\partial_i}{p}. \qedhere
    \end{equation*}
\end{proof}
We define $v^i = v[x^i]$, so 
\begin{equation}
    v = v^i \partial_i = \sum_{i=1}^n v^i \eval{\partial_i}{p}.
\end{equation}
\begin{remark}
    As a tangent vector is a linear combination of partial derivatives, just as a directional derivative is, there is an isomorphism between the tangent space $T_p M$ and the space of derivations of smooth functions on $M$ at $p$.
\end{remark}

\subsubsection{Pushforwards}\label{subsubsec:manifolds:tangent_spaces:pushforwards}

Given that tangent vectors extend the idea of directional derivatives to smooth manifolds, it would make sense for there to be a corresponding concept that extends the idea of the Jacobian.
We call this second concept the pushforward or differential of a function.
Essentially, given a smooth function $f \colon M \to N$ between smooth manifolds, the pushforward of $f$ is the best linear approximation of $f$ near a point.
It linearizes the map.
As it is convenient and correct, we will often consider tangent vectors acting on the representative functions of germs.
This will be the case for pushforwards and other concepts henceforth.
\begin{definition}[Pushforward]\label{def:pushforward}
    Let $f \colon M \to N$ be a smooth function between smooth manifolds $M$ and $N$.
    The pushforward of $f$ at $p$, or of $p$ by $f$, is a linear map $Df_p \colon T_p M \to T_{f(p)} N$ defined by the action of the tangent vector $Df_p(v) \in T_{f(p)} N$ on a smooth function $g \in C^\infty(V)$ such that $V \subseteq N$ is open and contains $f(p)$.
    \begin{equation}
        Df_p(v)(g) = v(g \circ f)
    \end{equation}
    The pushforward may alternatively be written as $(f_*)_p$, and if unambiguous, as $Df$ or $f_*$.
\end{definition}
\begin{remark}
    As we have restricted the definition of the pushforward to consider only smooth functions, we need not worry about its existence.
    However, with other models of the tangent space, such as equivalence classes of curves, we could have defined it for a broader class of functions.
    With derivations of smooth germs, $Df_p(v)(g) = v(g \circ f)$ requires $g \circ f$ to be smooth, so $f$ must be smooth.
\end{remark}
The pushforward is linear in $v$, since $Df_p(\lambda v + \mu w)(g) = (\lambda v + \mu w)(g \circ f) = \lambda \, Df_p(v)(g) + \mu \, Df_p(w)(g)$, and $Df_p(v)$ is a derivation at $f(p)$ because $Df_p(v)(gh) = v((g \circ f)(h \circ f)) = Df_p(v)(g) \, h(f(p)) + g(f(p)) \, Df_p(v)(h)$.
The pushforward gains its name as it pushes tangent vectors in $T_p M$ to tangent vectors in $T_{f(p)} N$.
Now, let us consider how it acts on the basis tangent vectors.
Given a chart $(U, \varphi)$ of $M$ around $p$ and a chart $(V, \psi)$ of $N$ around $f(p)$, with $\varphi = (x^1, \ldots, x^n)$ and $\psi = (y^1, \ldots, y^m)$, we can express the pushforward $Df_p \colon T_p M \to T_{f(p)} N$ using local coordinates.
It acts on a basis tangent vector $\partial_i \in T_p M$ to give $Df_p(\partial_i) \colon \mathcal{C}^\infty_{f(p)}(N, \R) \to \R$, which acts on a smooth function $g \in C^\infty(V)$ to give
\begin{equation}\label{eq:pushforward_at_basis_tangent_vector}
    Df_p \left(\eval{\frac{\partial}{\partial x^i}}{p}\right)(g) = \eval{\frac{\partial}{\partial x^i}}{p} (g \circ f).
\end{equation}
By the multivariable chain rule, we can write
\begin{equation}\label{eq:multivariable_chain_rule}
    \frac{\partial(g \circ f)}{\partial x^i} = \sum_{j=1}^m \frac{\partial g}{\partial y^j} \frac{\partial(y^j \circ f)}{\partial x^i}
\end{equation}
Substituting \eqref{eq:multivariable_chain_rule} into \eqref{eq:pushforward_at_basis_tangent_vector} gives
\begin{equation}
    Df_p\left(\eval{\frac{\partial}{\partial x^i}}{p}\right)(g) = \sum_{j=1}^m \frac{\partial(y^j \circ f)}{\partial x^i}(p) \eval{\frac{\partial g}{\partial y^j}}{f(p)}.
\end{equation}
This is true for all such $g$.
We may associate $f$ with $\psi \circ f \circ \varphi^{-1}$, and define the $j$-th \enquote*{component} of $f$ as $f^j = (y^j \circ f \circ \varphi^{-1}) \circ \varphi = y^j \circ f$.
With this (as well as making use of Einstein summation convention and omitting points), it follows that
\begin{equation}
    Df_p(\partial_i) = Df \left(\frac{\partial}{\partial x^i}\right) = \frac{\partial f^j}{\partial x^i} \frac{\partial}{\partial y^j} = \partial_{i} f^j \partial_{j},
\end{equation}
where $\partial_i = \partial/\partial x^i$ and $\partial_j = \partial/\partial y^j$.
The linearity of the pushforward can be used to evaluate it on a general tangent vector $v = v^i \partial_i \in T_p M$.
\begin{equation}
    Df_p(v) = Df(v^i \partial_i) = v^i Df(\partial_i)
\end{equation}
There are various different ways, notationally, to express this -- each having some level of compactness and clarity.
\begin{equation}\label{eq:pushforward_in_coordinates}
    Df_p(v) = v^i \partial_i f^j \partial_j = \sum_{i=1}^n v^i \partial_i f^j \partial_j = \sum_{i=1}^n \sum_{j=1}^m v^i \frac{\partial f^j}{\partial x^i} \frac{\partial}{\partial y^j}
\end{equation}
Regardless, this shows that the pushforward is of the form of an $m \times n$ matrix when expressed with local coordinates.
It may be written by stating only its generic component.
Note that finding the generic component only required considering the action of the pushforward on a generic basis tangent vector, by virtue of pushforward being a linear map.
\begin{equation}
    (Df_p)^j_i = \partial_i f^j = \frac{\partial f^j}{\partial x^i}
\end{equation}
This is of the exact same form as the Jacobian, which is to be expected.
It is perhaps not entirely correct to state a priori the two as equal, since the Jacobian only works for functions $f \colon \R^n \to \R^m$.
This may be resolved by extending the definition of the Jacobian to functions $f \colon M \to N$.
Writing it out in full, without missing details, we have
\begin{equation}
    Df_p = \mathbf{J}_f (p) = \begin{bmatrix}
        \eval{\dfrac{\partial f^1}{\partial x^1}}{p} & \cdots & \eval{\dfrac{\partial f^1}{\partial x^n}}{p} \\[0.5em]
        \vdots & \ddots & \vdots \\[0.5em]
        \eval{\dfrac{\partial f^m}{\partial x^1}}{p} & \cdots & \eval{\dfrac{\partial f^m}{\partial x^n}}{p}
    \end{bmatrix}.
\end{equation}
Being an extension and generalization of the Jacobian, which itself generalizes the ordinary derivative, properties of the derivative take shape as analogous properties of the pushforward.
One such example is the chain rule.
\begin{proposition}[Chain Rule for Pushforwards]\label{prop:chain_rule_for_pushforwards}
    Let $M$, $N$, and $P$ be manifolds, let $f \in C^\infty(M, N)$ and $g \in C^\infty(N, P)$ be functions, and let $p \in M$ be a point.
    Then,
    \begin{equation}
        D(g \circ f)_p = Dg_{f(p)} \circ Df_p.
    \end{equation}
\end{proposition}
\begin{proof}
    For $v \in T_p M$ and a smooth function $h$ defined near $g(f(p))$,
    \begin{equation*}
        D(g \circ f)_p(v)(h) = v(h \circ g \circ f) = Df_p(v)(h \circ g) = Dg_{f(p)}(Df_p(v))(h). \qedhere
    \end{equation*}
\end{proof}
The pushforward can be used to express tangent vectors after a change of coordinates.
If $T_p M$ is given using coordinates $x^1, \ldots, x^n$, and we change to coordinates $y^1, \ldots, y^n$ given by another chart, the $\partial/\partial y^i \in T_p M$.
Thus,
\begin{equation*}
    \eval{\frac{\partial}{\partial y^i}}{p} = \sum_{j=1}^n \lambda^j \eval{\frac{\partial}{\partial x^j}}{p},
\end{equation*}
for some $\lambda^j \in \R$.
Applying both sides to $x^k$ gives
\begin{equation*}
    \eval{\frac{\partial}{\partial y^i}}{p}(x^k) = \frac{\partial x^k}{\partial y^i}(p) = \sum_{j=1}^n \lambda^j \eval{\frac{\partial}{\partial x^j}}{p} (x^k) = \sum_{j=1}^n \lambda^j \delta^k_j = \lambda^k,
\end{equation*}
so we get
\begin{equation}
    \eval{\frac{\partial}{\partial y^i}}{p} = \sum_{j=1}^n \frac{\partial x^j}{\partial y^i}(p) \eval{\frac{\partial}{\partial x^j}}{p}.
\end{equation}
This is the standard change-of-coordinate formula.

Now, we may define curves and velocities of curves.
\begin{definition}[Curve]\label{def:curve}
    A curve is a map $\gamma \colon I \to M$, where $I \subseteq \R$ is a nonempty interval.
    If $\gamma$ is smooth, it is a smooth curve.
\end{definition}
\begin{definition}[Derivative of Curve]\label{def:derivative_of_curve}
    Given a curve $\gamma \colon I \to M$ and a point $t_0 \in I$, the velocity of $\gamma$ at $t_0$ is a tangent vector $\gamma'(t_0) \in T_{\gamma(t_0)} M$ defined by
    \begin{equation}
        \gamma'(t_0) = D\gamma_{t_0} \left(\eval{\frac{d}{dt}}{t_0}\right),
    \end{equation}
    where $\eval*{\frac{d}{dt}}{t_0}$ is the standard basis vector of $T_{t_0} \R$ (we use the regular derivative symbol since it is one dimensional).
    The velocity can also be written as $\dot{\gamma}(t_0)$, $d\gamma/dt(t_0)$, or $\eval*{d\gamma/dt}{t_0}$.
\end{definition}
The velocity $\gamma'(t_0)$ acts on a function $f \colon M \to \R$ as follows.
\begin{equation}
    \gamma'(t_0)(f) = D\gamma_{t_0} \left(\eval{\frac{d}{dt}}{t_0}\right)(f) = \eval{\frac{d}{dt}}{t_0} (f \circ \gamma) = (f \circ \gamma)'(t_0)
\end{equation}
As the composite $f \circ \gamma$ is a function from $\R$ to $\R$, its derivative is ordinary.
If $t_0$ is an endpoint of $I$, then the velocity at $t_0$ still exists as long as we interpret the derivative as a one-sided derivative.

The velocity $\gamma'(t_0)$ can be expressed using the formula for the pushforward with coordinates \eqref{eq:pushforward_in_coordinates}.
Here, $\gamma^i = x^i \circ \gamma$, as is usually the case when we write a function in terms of coordinates.
\begin{equation}
    \gamma'(t_0) = \sum_{i=1}^n \frac{d\gamma^i}{dt}(t_0) \eval{\frac{\partial}{\partial x^i}}{\gamma(t_0)}
\end{equation}
By the above formula, the velocity of a curve is a tangent vector whose components are the derivatives of the components of the curve.
This makes it not too dissimilar from the velocity of a curve in Euclidean space.

Furthermore, every tangent vector on a manifold is the velocity vector of some curve in the manifold.
This allows us to reinterpret various results and concepts regarding tangent vectors in terms of velocities.
\begin{proposition}\label{prop:tangent_vectors_are_velocities_of_curves}
    Given a smooth manifold $M$ with or without boundary and a point $p \in M$, every tangent vector $v \in T_p M$ is the velocity of some smooth curve in $M$.
\end{proposition}
\begin{proof}
    This proof is modelled after the proof of \cite[Prop.~3.23]{Lee2012}.

    First, suppose that $p \in \Int M$; it may still be the case that $\partial M \neq \varnothing$.
    Let $(U, \varphi)$ be a chart of $M$ with $\varphi(p) = 0$ and $\varphi = (x^1, \ldots, x^n)$ so that $v = v^i \eval{\partial_i}{p}$.
    We define a curve $\hat{\gamma} \colon (-\varepsilon, \varepsilon) \to \R^n$ by
    \begin{equation}\label{eq:euclidean_curve_coordinate_representation}
        \hat{\gamma}(t) = (tv^1, \ldots, tv^n),
    \end{equation}
    with $\varepsilon > 0$ being a real number small enough so that the image of $\hat{\gamma}$ is within $\varphi(U)$.
    Then, we have the associated curve $\gamma \colon (-\varepsilon, \varepsilon) \to M$ given by
    \begin{equation*}
        \gamma(t) = \varphi^{-1}(\hat{\gamma}(t))
    \end{equation*}
    Since $\hat{\gamma}(0) = 0$, it follows that $\gamma(0) = \varphi^{-1}(0) = p$.
    The velocity at zero is given by
    \begin{align*}
        \gamma'(0) &= \sum_{i=1}^n \frac{d(x^i \circ \varphi^{-1} \circ \hat{\gamma})}{dt}(0) \eval{\frac{\partial}{\partial x^i}}{p} \\
        &= \sum_{i=1}^n \frac{d(tv^i)}{dt}(0) \eval{\frac{\partial}{\partial x^i}}{p} \\
        &= \sum_{i=1}^n v^i \eval{\frac{\partial}{\partial x^i}}{p} \\
        &= v.
    \end{align*}
    
    Now suppose that $p \in \partial M$.
    Let $(U, \varphi)$ be a smooth boundary chart with $\varphi(p) = 0$ with $\varphi$ and $v \in T_p M$ being expressed the same way as before.
    We want a curve $\gamma$ whose Euclidean version has representation shown in \eqref{eq:euclidean_curve_coordinate_representation}, $\hat{\gamma}$ may only map to $\mathbb{H}^n$.
    The solution is to restrict the domain of $\hat{\gamma}$ to $[0, \varepsilon)$ if $v^n > 0$ or to $(-\varepsilon, 0]$ if $v^n < 0$; if $v^n = 0$, then we may keep the domain the same.
    In any of these cases, we similarly arrive at $\gamma(0) = p$ and $\gamma'(0) = v$.
\end{proof}
The velocity of the composite of a curve with a smooth function is given by the following formula.
\begin{proposition}\label{prop:pushforward_of_velocity}
    Let $f \colon M \to N$ be a smooth function between smooth manifolds, and let $\gamma \colon I \to M$ be a smooth curve.
    At any $t_0 \in I$, the velocity of $f \circ \gamma \colon I \to N$ at $t_0$ is given by
    \begin{equation}
        (f \circ \gamma)'(t_0) = Df_{\gamma(t_0)}(\gamma'(t_0)).
    \end{equation}
\end{proposition}
\begin{proof}
    By expanding out the definition of the velocity of a curve and using the chain rule, we have
    \begin{equation*}
        (f \circ \gamma)'(t_0) = D(f \circ \gamma)_{t_0}\left(\eval{\frac{d}{dt}}{t_0}\right) = Df_{\gamma(t_0)} \circ D\gamma_{t_0} \left(\eval{\frac{d}{dt}}{t_0}\right) = Df_{\gamma(t_0)}(\gamma'(t_0)). \qedhere
    \end{equation*}
\end{proof}
Whilst this shows how to evaluate the velocity of a composite curve, it can also be used to evaluate a pushforward.
Since every tangent vector is a velocity by \Cref{prop:tangent_vectors_are_velocities_of_curves}, this allows us to evaluate $Df_p(v)$ for any $v \in T_p M$.
\begin{corollary}\label{cor:pushforward_formula_using_velocities}
    Let $f \colon M \to N$ be a smooth function between smooth manifolds, $p \in M$ be a point, and $v \in T_p M$ be a tangent vector.
    Then,
    \begin{equation}
        Df_p(v) = (f \circ \gamma)'(0)
    \end{equation}
    for any smooth curve $\gamma \colon I \to M$ such that $0 \in I$, $\gamma(0) = p$, and $\gamma'(0) = v$.
\end{corollary}
\begin{proof}
    By \cref{prop:pushforward_of_velocity}, $(f \circ \gamma)'(0) = Df_p(\gamma'(0)) = Df_p(v)$.
\end{proof}
\begin{remark}
    Note the fact that we do not need to carefully construct a curve for this corollary as we had to for \Cref{prop:tangent_vectors_are_velocities_of_curves}.
    The statement is universal rather than existential, and it follows immediately from the preceding proposition.
\end{remark}
We can use pushforwards to categorize smooth functions between manifolds as being immersions or submersions.
\begin{definition}[Immersion and Submersion]\label{def:immersion_and_submersion}
    Let $M$ and $N$ be smooth manifolds, and let $f \colon M \to N$ be a smooth function.
    The function $f$ is an immersion at $p \in M$ if $Df_p \colon T_p M \to T_{f(p)} N$ is injective, and it is a submersion at $p$ if $Df_p$ is surjective.
    It is an immersion or submersion if it is an immersion or submersion at every point in $M$.
\end{definition}
\begin{definition}[Embedded Submanifold]\label{def:embedded_submanifold}
    A smooth immersion $f \colon M \to N$ is a smooth embedding if it is a homeomorphism onto its image $f(M)$ with the subspace topology.
    An embedded submanifold of $N$ is a subset $S \subseteq N$ that is a manifold in the subspace topology, with a smooth structure for which the inclusion $S \hookrightarrow N$ is a smooth embedding.
    For $p \in S$, we identify $T_p S$ with its image in $T_p N$ under the pushforward of the inclusion.
\end{definition}
The boundary of a manifold with boundary is an example.
\begin{proposition}\label{prop:boundary_is_embedded_submanifold}
    Let $M$ be a smooth $n$-dimensional manifold with boundary.
    Then $\partial M$ is an embedded submanifold of $M$ of dimension $n - 1$.
\end{proposition}
\begin{proof}[Sketch]
    By \cref{prop:topological_invariance_of_the_boundary}, the transition maps between boundary charts send points of $\partial \mathbb{H}^n$ to points of $\partial \mathbb{H}^n$.
    Hence they restrict to smooth maps with smooth inverses between open subsets of $\partial \mathbb{H}^n \cong \R^{n-1}$.
    It follows that the restrictions of the boundary charts to $\partial M$ form a smooth atlas for the subspace topology, which is Hausdorff and second countable.
    In these charts, the inclusion of $\partial M$ into $M$ is the standard inclusion of $\R^{n-1}$ into $\mathbb{H}^n$, so it is an immersion and a homeomorphism onto its image.
\end{proof}
This may be used to define a type of submanifold more general than embedded submanifolds.
\begin{definition}[Immersed Submanifold]\label{def:immersed_submanifold}
    Let $M$ be a smooth manifold.
    An immersed submanifold of $M$ is a subset $S \subseteq M$ with a topology making it a topological manifold and a smooth structure such that the inclusion map $\iota \colon S \hookrightarrow M$ is a smooth immersion.
\end{definition}

\subsection{Tangent Bundles}\label{subsec:manifolds:tangent_bundles}

With tangent spaces, we are able to generalize differentiation to functions defined on smooth manifolds.
However, a tangent space is only a local object, and we are currently unable to talk about anything on a global scale.
This is the motivation for defining the tangent bundle, which is the collection of all tangent spaces of a manifold.
\begin{definition}[Tangent Bundle]\label{def:tangent_bundle}
    Given a smooth manifold $M$, the tangent bundle $TM$ is the disjoint union of the tangent spaces of $M$ at each point $p \in M$.
    \begin{equation}
        TM = \coprod_{p \in M} T_p M
    \end{equation}
    An element of $TM$ is an ordered pair $(p, v)$, where $v \in T_p M$.
    We may write the tangent vector, as $v_p$.
    We may identify the tangent vector $v$ with the ordered pair $(p, v)$ and set them as equal by abuse of notation; elements of $TM$ may be called tangent vectors by abuse of language.
    There is a natural projection map $\pi \colon TM \to M$ defined by $\pi(p, v) = p$.
\end{definition}
\begin{remark}
    The elements of the disjoint union $\coprod_{i \in I} A_i$ are usually defined to be $(a_i, i)$ instead of $(i, a_i)$, with the latter being the way elements of $TM$ are written.
    This distinction does not really matter.
\end{remark}
Whilst being defined as a mere collection, the tangent bundle has much more structure.
It is itself a manifold, with a natural topology and smooth structure.
\begin{proposition}\label{prop:tangent_bundle_is_manifold}
    The tangent bundle of an $n$-dimensional smooth manifold is a $2n$-dimensional smooth manifold.
\end{proposition}
\begin{proof}
    We use \Cref{lem:smooth_manifold_chart_lemma}.
    Let $(U, \varphi)$ be a smooth chart of $M$.
    The preimage $\pi^{-1}(U) \subseteq TM$ is the set of all tangent vectors to $M$ at all points in $U$, which is also the tangent bundle of $U$.
    \begin{equation}
        TU = \coprod_{p \in U} T_p U = \coprod_{p \in U} T_p M = \pi^{-1} (U)
    \end{equation}
    Furthermore, we have the vector space isomorphisms $\Phi_p \colon T_p M \to \R^n$ for all $p \in U$, given by $v^i \partial_i \mapsto (v^1, \ldots, v^n)$ with respect to the coordinate basis of the chart.
    The collection of these maps gives the bijection $\Phi \colon TU \to \varphi(U) \times \R^n$ defined by $\Phi(p, v) = (\varphi(p), \Phi_p(v))$ for $v \in T_p U$.
    The images of these bijections, $\Phi(TU) = \varphi(U) \times \R^n$ are open subsets of $\R^{2n}$.
    Thus, by considering the collection of sets $TU$ for each chart domain $U \subseteq M$, property 1 of the lemma is satisfied.
    We will consider a pair $(TU, \Phi)$ to be a chart for $TM$.
    \[\begin{tikzcd}[sep=large]
        TU & {\varphi(U) \times \R^n} \\
        U & {\varphi(U)}
        \arrow["{\Phi}", from=1-1, to=1-2]
        \arrow["\eval{\pi}{TU}"', from=1-1, to=2-1]
        \arrow["{(u, v) \mapsto u}", from=1-2, to=2-2]
        \arrow["\varphi"', from=2-1, to=2-2]
    \end{tikzcd}\]
    Now, consider two smooth charts $(U, \varphi)$ and $(V, \psi)$ for $M$ with the corresponding charts $(TU, \Phi)$ and $(TV, \Psi)$ for $TM$.
    The sets
    \begin{equation*}
     \Phi(TU \cap TV) = \varphi(U \cap V) \times \R^n
    \end{equation*}
    and
    \begin{equation*}
    \Psi(TU \cap TV) = \psi(U \cap V) \times \R^n
    \end{equation*}
    are open in $\R^{2n}$, so property 2 is satisfied.
    The transition map $\Psi \circ \Phi^{-1} \colon \varphi(U \cap V) \times \R^n \to \psi(U \cap V) \times \R^n$ can be written as
    \begin{equation}
        \Psi \circ \Phi^{-1}(x, a) = (\psi \circ \varphi^{-1}(x), D(\psi \circ \varphi^{-1})_x(a)),
    \end{equation}
    which is smooth, as the second component is $J(x)a$, where the entries of the Jacobian matrix $J$ of $\psi \circ \varphi^{-1}$ are smooth in $x$; property 3 is satisfied.
    Note that $a \in \R^n$ would be a Euclidean vector corresponding to some tangent vector.

    We get a countable cover $\{TU_i\}_{i \in I}$ of $TM$ from a countable cover $\{U_i\}_{i \in I}$ of $M$ using smooth coordinate domains.
    This satisfies property 4 of the lemma.
    Consider the fact that two points of $TM$ are in the same chart $TU$ if they are in the same tangent space.
    Otherwise, since $M$ is Hausdorff, there are disjoint smooth coordinate domains $U \ni p$ and $V \ni q$ of $M$, and $TU$ and $TV$ are disjoint coordinate neighbourhoods containing $(p, v)$ and $(q, w)$, respectively.
    This shows that $TM$ is Hausdorff, satisfying property 5.
    By \Cref{lem:smooth_manifold_chart_lemma}, the tangent bundle $TM$ with the given topology and the smooth atlas $\mathcal{A} = \{TU_i\}_{i \in I}$ forms a second-countable smooth manifold.
\end{proof}
\begin{remark}
    Given a chart $(U, \varphi)$ of $M$ and the corresponding chart $(TU, \Phi)$ of $TM$, the coordinate representation of the projection map $\pi$ is $\pi(p, v) = x$, so the projection map is smooth.
\end{remark}
\begin{example}
    Consider the unit circle $\mathbb{S}^1 \coloneq \{(x, y) \in \R^2 : x^2 + y^2 = 1\}$.
    The tangent bundle is given by
    \begin{equation}
        T\mathbb{S}^1 = \coprod_{p \in \mathbb{S}^1} T_p \mathbb{S}^1 = \bigcup_{p \in \mathbb{S}^1} \{p\} \times T_p \mathbb{S}^1.
    \end{equation}
    We may parametrize the circle by $\theta \mapsto (\cos\theta, \sin\theta)$ with $\theta \in [0, 2\pi)$.
    Then, the tangent space $T_p \mathbb{S}^1$ is spanned by the tangent vector $\partial_\theta = (-\sin\theta, \cos\theta)$, which is in terms of the standard basis tangent vectors of $T_p \R^2$.
    Our basis tangent vector works for any point on the circle -- this is an important detail.
    Given a tangent vector $v \in T_p \mathbb{S}^1$, we define the map $\Phi \colon T\mathbb{S}^1 \to \mathbb{S}^1 \times \R$ by $\Phi(p, v) = (p, \lambda)$, where $\lambda \in \R$ is the unique real number such that $v = \lambda\partial_\theta$.
    This map is a diffeomorphism, meaning that $T\mathbb{S}^1 \cong \mathbb{S}^1 \times \R$.
\end{example}
We can consider the global object we get if we take the pushforwards $Df_p \colon T_p M \to T_{f(p)} N$ for each $p \in M$ and combine them.
This gives us the definition of the global pushforward, which allows us to more meaningfully talk about the derivative of a function in the sense of how the function changes globally.
\begin{definition}[Global Pushforward]\label{def:global_pushforward}
    Given a smooth function $f \colon M \to N$ between smooth manifolds $M$ and $N$, the global pushforward of $f$ is the map $Df \colon TM \to TN$ defined by
    \begin{equation}
        Df(p, v) = (f(p), Df_p(v)).
    \end{equation}
    The global pushforward may alternatively be denoted $f_*$.
    It is a bundle map induced by the function $f$, making the following diagram commute.
    \[\begin{tikzcd}[sep=large]
        TM & TN \\
        M & N
        \arrow["Df", from=1-1, to=1-2]
        \arrow["{\pi_M}"', from=1-1, to=2-1]
        \arrow["{\pi_N}", from=1-2, to=2-2]
        \arrow["f"', from=2-1, to=2-2]
    \end{tikzcd}\]
\end{definition}
The global pushforward of a smooth function is itself smooth.
\begin{proposition}\label{prop:global_pushforward_is_smooth}
    Given a smooth function $f \colon M \to N$ between smooth manifolds $M$ and $N$, its global pushforward $Df \colon TM \to TN$ is smooth.
\end{proposition}
\begin{proof}
    In the usual coordinates $(x^1, \ldots, x^n, v^1, \ldots, v^n)$ for $TM$ and $(y^1, \ldots, y^m, w^1, \ldots, w^m)$ for $TN$, \eqref{eq:pushforward_in_coordinates} gives the action of $Df$ (with the usual abuse of notation).
    \begin{equation}
        Df(x^1, \ldots, x^n, v^1, \ldots, v^n) = \left(f^1(x), \ldots, f^m(x), \sum_{i=1}^n \frac{\partial f^1}{\partial x^i}(x) v^i, \ldots, \sum_{i=1}^n \frac{\partial f^m}{\partial x^i}(x) v^i\right)
    \end{equation}
    The components $f^j$ and their partial derivatives are smooth because $f$ is smooth, and the last $m$ entries are linear in $v$, so this expression is smooth.
\end{proof}
\begin{corollary}\label{cor:global_pushforward_properties}
    Given smooth functions $f \colon M \to N$ and $g \colon N \to P$, where $M$, $N$, and $P$, are smooth manifolds, the following statements are true.
    \begin{enumerate}
        \item $D(g \circ f) = Dg \circ Df$
        \item $D(\id_M) = \id_{TM}$
        \item If $f$ is a diffeomorphism, then $Df \colon TM \to TN$ is a diffeomorphism, and $(Df)^{-1} = D(f^{-1})$.
    \end{enumerate}
\end{corollary}
\begin{proof}
    Part 1 is \cref{prop:chain_rule_for_pushforwards} applied at each $p \in M$, and part 2 holds because $D(\id_M)_p(v)(g) = v(g \circ \id_M) = v(g)$.
    For part 3, parts 1 and 2 give $Df \circ D(f^{-1}) = D(f \circ f^{-1}) = \id_{TN}$ and $D(f^{-1}) \circ Df = \id_{TM}$.
    Both maps are smooth by \cref{prop:global_pushforward_is_smooth}, so $Df$ is a diffeomorphism with the stated inverse.
\end{proof}
\begin{remark}
    By property 3, the notation $Df^{-1}$ is unambiguous and thus may freely be used when $f$ is a diffeomorphism.
\end{remark}

\subsection{Vector Bundles}\label{subsec:manifolds:vector_bundles}

The tangent bundle assigns a vector space to each point of a manifold, in a way that varies smoothly from point to point.
The cotangent bundle, and the bundles whose sections are differential forms, do the same.
These are all examples of vector bundles.
\begin{definition}[Vector Bundle]\label{def:vector_bundle}
    Let $M$ be a smooth manifold.
    A smooth vector bundle of rank $r$ over $M$ is a smooth manifold $E$ with a smooth surjective map $\pi \colon E \to M$, such that each fibre $E_p = \pi^{-1}(p)$ has the structure of an $r$-dimensional real vector space, and every point of $M$ has an open neighbourhood $U$ with a diffeomorphism $\Phi \colon \pi^{-1}(U) \to U \times \R^r$ satisfying $\proj_1 \circ \Phi = \pi$ which restricts to a linear isomorphism $E_p \to \{p\} \times \R^r$ for each $p \in U$.
    We call $E$ the total space, $M$ the base, and $\Phi$ a local trivialization of $E$ over $U$.
\end{definition}
\begin{definition}[Section]\label{def:section}
    A section of a vector bundle $\pi \colon E \to M$ is a smooth map $\sigma \colon M \to E$ such that $\pi \circ \sigma = \id_M$.
    The space of all sections of $E$ is denoted by $\Gamma(E)$.
\end{definition}
Constructing a vector bundle from its fibres and its trivializations is the analogue of constructing a manifold from its charts.
\begin{lemma}[Vector Bundle Construction Lemma]\label{lem:vector_bundle_construction_lemma}
    Let $M$ be a smooth manifold, and let $E = \coprod_{p \in M} E_p$ be a disjoint union of $r$-dimensional real vector spaces with the projection $\pi \colon E \to M$.
    Suppose that $\{U_\alpha\}$ is an open cover of $M$ with bijections $\Phi_\alpha \colon \pi^{-1}(U_\alpha) \to U_\alpha \times \R^r$ which restrict to linear isomorphisms $E_p \to \{p\} \times \R^r$, and that for all $\alpha$ and $\beta$ there is a smooth map $\tau_{\alpha\beta} \colon U_\alpha \cap U_\beta \to \mathrm{GL}_r(\R)$ such that
    \begin{equation}
        \Phi_\alpha \circ \Phi_\beta^{-1}(p, v) = (p, \tau_{\alpha\beta}(p) v).
    \end{equation}
    Then $E$ has a unique topology and smooth structure making it a smooth vector bundle over $M$ with each $\Phi_\alpha$ a local trivialization.
\end{lemma}
\begin{proof}[Sketch]
    For each $\alpha$ and each smooth chart $(V, \varphi)$ of $M$, consider the set $\pi^{-1}(U_\alpha \cap V)$ with the map $(\varphi \times \id) \circ \Phi_\alpha$ to $\R^n \times \R^r$.
    By hypothesis, the transition maps $(x, v) \mapsto (\varphi' \circ \varphi^{-1}(x), \tau(\varphi^{-1}(x)) v)$ between these are smooth.
    As $M$ is second countable, countably many of these sets cover $E$.
    Two points of $E$ in the same fibre lie in a common set, and two points in different fibres lie in the sets $\pi^{-1}(U)$ and $\pi^{-1}(U')$ for disjoint open sets $U, U' \subseteq M$, since $M$ is Hausdorff.
    Hence \cref{lem:smooth_manifold_chart_lemma} applies, and it gives the topology and smooth structure.
\end{proof}
\begin{example}\label{ex:tangent_bundle_is_a_vector_bundle}
    The tangent bundle $TM$ of an $n$-dimensional smooth manifold is a vector bundle of rank $n$.
    The charts $(TU, \Phi)$ in the proof of \cref{prop:tangent_bundle_is_manifold}, composed with $\varphi^{-1} \times \id$, are local trivializations, and the transition functions $\tau_{\alpha\beta}$ are the Jacobian matrices of the transition functions of $M$.
    The sections of $TM$ are the vector fields, which we study in the next subsection.
\end{example}
The cotangent spaces of a manifold assemble into a vector bundle in the same way.
\begin{definition}[Cotangent Space]
    Let $M$ be a smooth manifold and let $p \in M$ be a point.
    The cotangent space $T^*_p M$ of $M$ at $p$ is the dual of the tangent space $T_p M$.
    An element of $T^*_p M$ is called a cotangent vector or covector.
\end{definition}
\begin{remark}
    A covector may be thought of as a type $(0, 1)$ tensor.
    It is also called a one-form, and it may be denoted by $\omega$ or $\omega_p$.
    The latter notation may be preferred to distinguish it from differential one-forms.
\end{remark}
\begin{definition}[Cotangent Bundle]
    Given a smooth manifold $M$, the cotangent bundle $T^*M$ is the disjoint union of the cotangent spaces of $M$ at each point $p \in M$.
    \begin{equation}
        T^*M = \coprod_{p \in M} T^*_p M
    \end{equation}
    An element of $T^*M$ is an ordered pair $(p, \omega_p)$, where $\omega_p \in T^*_p M$. We may write the covector as $\omega_p$.
    We may identify the covector $\omega_p$ with the ordered pair $(p, \omega_p)$ and set them as equal by abuse of notation; elements of $T^*M$ may be called covectors by abuse of language.
    There is a natural projection map $\pi \colon T^*M \to M$ defined by $\pi(p, \omega_p) = p$.
\end{definition}
Given a basis for the tangent space $T_p M$, we have a dual basis for $T^*_p M$.
The $i$-th basis covector is dual to the $i$-th basis tangent vector.
\begin{equation}
    \d{x}^j(\partial_i) = \delta^j_i
\end{equation}
Strictly, the basis covector is $\eval*{\d{x}}{p}$, though we often shorten it for convenience just as we shorten $\omega_p$ to $\omega$.
With respect to the dual basis, we can express a covector $\omega_p$ as
\begin{equation}
    \omega_p = (\omega_p)_i \eval{\d{x}^i}{p}
\end{equation}
\begin{proposition}\label{prop:cotangent_bundle_and_exterior_powers_are_vector_bundles}
    Let $M$ be an $n$-dimensional smooth manifold.
    The cotangent bundle $T^*M$ and, for each $k$, the bundle $\extp^k(T^*M) = \coprod_{p \in M} \extp^k(T^*_p M)$ of $k$-covectors are smooth vector bundles over $M$, of ranks $n$ and $\binom{n}{k}$ respectively.
\end{proposition}
\begin{proof}
    We use \cref{lem:vector_bundle_construction_lemma}.
    A smooth chart $(U, \varphi)$ of $M$ gives, at each point $p \in U$, the basis $\d{x}^i|_p$ of $T^*_p M$, and the basis $\d{x}^I|_p = \d{x}^{i_1}|_p \wedge \cdots \wedge \d{x}^{i_k}|_p$, for increasing multi-indices $I$, of $\extp^k(T^*_p M)$.
    Taking coefficients with respect to these bases gives bijections $\pi^{-1}(U) \to U \times \R^{\binom{n}{k}}$ which are linear on each fibre.
    For a second chart with coordinates $y^j$, we have $\d{y}^j = (\partial y^j/\partial x^i) \d{x}^i$.
    Hence the transition matrix on $T^*M$ is given by the transpose of the Jacobian matrix of the transition function of $M$ (or its inverse, in the other direction), and the transition matrix on $\extp^k(T^*M)$ is given by the corresponding $k \times k$ minors.
    These are smooth functions on the overlap, and the matrices are invertible, as the inverse is given by the same construction for the inverse transition function.
\end{proof}
We can naturally pair vectors in $TM$ to covectors in $T^*M$, but they are not naturally isomorphic to each other.
For that to be the case, we require additional structure such as a Riemannian metric or a symplectic form.


\subsection{Vector Fields}\label{subsec:manifolds:vector_fields}

With the introduction of vector bundles, we can consider the sections of the tangent bundle.
These are the objects we call vector fields.
\begin{definition}[Vector Field]\label{def:vector_field}
   Let $M$ be a smooth manifold, and let $TM$ be its tangent bundle with the projection map $\pi$.
   A vector field $X \colon M \to TM$ on $M$ is a section of $TM$, i.e. a smooth right inverse of $\pi$.
   \begin{equation}
    \pi \circ X = \id_M
   \end{equation}
   We may denote $X(p)$ by $X_p$.
   A map $X \colon M \to TM$ with $\pi \circ X = \id_M$ that is not necessarily smooth, or even continuous, may be called a rough vector field.
   The space of all vector fields on $M$ is denoted by $\mathfrak{X}(M)$ or $\Gamma(TM)$.
\end{definition}
\begin{remark}
    Vector fields may equivalently be defined as derivations $X \colon C^\infty(M) \to C^\infty(M)$ with $(Xf)(p) = X_p(f)$.
    The definition is equivalent in the sense that the space of all such derivations is isomorphic to $\mathfrak{X}(M)$.
    The fact that vector fields are considered derivations is a useful property and may be utilized when required, e.g. it allows us to verify whether an object is a vector field by seeing if it is a derivation.
\end{remark}
Vector fields are an object commonly first found in the study of multivariable calculus.
We may visualize vector fields using the same methods found in such studies -- as a collection of arrows attached to every point on a manifold (that vary smoothly from point to point).

Given a smooth coordinate chart $(U, \varphi)$ for $M$, a vector field $X \in \mathfrak{X}(M)$ may be expressed at a point $p \in M$ in terms of the basis tangent vectors $\partial_i \in T_p M$ as
\begin{equation}
    X_p = X^i(p) \eval{\frac{\partial}{\partial x^i}}{p},
\end{equation}
where we have $n$ functions $X^i \colon U \to \R$ called the components of $X$ with respect to our chart.
By our usual associations, the natural coordinate representation of $X$ is given by $\widehat{X}(x) = (x^1, \ldots, x^n, X^1(x), \ldots, X^n(x))$, where $x \in U$.
Here, $X^i$ is the $i$-th component of $X$ with respect to the coordinates $x^i$.
Furthermore, the chart determines a vector field $\partial_i \in \mathfrak{X}(U)$ on $U$ called the $i$-th coordinate vector field.
These vector fields have constant components and are thus smooth.
The vector field $X$ may be written using the coordinates vector fields $\partial_i$ and the components $X^i$.
\begin{equation}
    X = X^i \frac{\partial}{\partial x^i}
\end{equation}
We can consider collections of vector fields defined on some subset $A \subseteq M$ and whether they perform a frame, which is akin to a basis for vector spaces.
This is done by considering bases for the tangent spaces of each point within the subset.
\begin{definition}[Frame of Manifold]\label{def:frame_of_manifold}
    Let $M$ be a smooth manifold and consider a subset $A \subseteq M$.
    A $k$-tuple $(X_{(1)}, \ldots, X_{(k)})$, also denoted by $(X_{(i)})$ of vector fields on a subset $A \subseteq M$ is linearly independent if $(\eval*{X_{(1)}}{p}, \ldots, \eval*{X_{(k)}}{p})$ a linearly independent $k$-tuple in $T_p M$ for each $p \in A$.
    It spans $TA$ if $(\eval*{X_{(1)}}{p}, \ldots, \eval*{X_{(k)}}{p})$ spans $T_p M$ for each $p \in A$.
    
    If $A$ is an open set and $(X_{(1)}, \ldots, X_{(k)})$ is linearly independent and spans $TA$, so that $k = n$, it is called a local frame for $M$.
    If $A$ is $M$, then a local frame $(X_{(1)}, \ldots, X_{(k)})$ is called a global frame for $M$.
    If all $X_{(i)}$ are smooth, then a local frame $(X_{(1)}, \ldots, X_{(k)})$ is called a smooth frame for $M$.
    If $M$ is $\R^n$ and the vectors $(\eval*{X_{(1)}}{p}, \ldots, \eval*{X_{(k)}}{p})$ are orthonormal with respect to the Euclidean inner product, then $(X_{(1)}, \ldots, X_{(k)})$ is called an orthonormal frame.
    If there exists a global frame for $M$, it is said to be parallelizable.
\end{definition}
In general, a smooth function $f \colon M \to N$ between smooth manifolds admits a pushforward $Df \colon TM \to TN$.
From this, we may obtain a vector $Df_p(X_p) \in T_{f(p)} N$.
However, we may not be able to obtain a vector field $Df(X) \in \mathfrak{X}(N)$ since $f$ may not be injective.
When $f$ is not surjective, there is no canonical way to assign vectors to points outside $f(M)$.
We characterize when we can by using the property of $f$-relatedness.
\begin{definition}[$f$-Related]\label{f-related}
    Let $M$ and $N$ be smooth manifolds and let $f \colon M \to N$ be a smooth function.
    We say that two vector fields $X \in \mathfrak{X}(M)$ and $Y \in \mathfrak{X}(N)$ are $f$-related if the following diagram commutes.
    \[\begin{tikzcd}[sep=large]
        TM & TN \\
        M & N
        \arrow["Df", from=1-1, to=1-2]
        \arrow["{X}", from=2-1, to=1-1]
        \arrow["{Y}"', from=2-2, to=1-2]
        \arrow["f"', from=2-1, to=2-2]
    \end{tikzcd}\]
    If $f$ is a diffeomorphism, for any vector $X \in \mathfrak{X}(M)$, we denote the unique $f$-related vector field $Df(X)$ and call it the pushforward of $f$ at $X$ or of $X$ by $f$.
    This is also denoted by $f_*X$.
\end{definition}
\begin{remark}
    The pushforward at $p$ may be explicitly computed via the equation
    \begin{equation}
        (f_*X)_p = Df_{f^{-1}(p)}(X_{f^{-1}(p)}).
    \end{equation}
\end{remark}
Given we have two vector fields $X, Y \in \mathfrak{X}(M)$, their sum $X + Y$ given by $(X + Y)(f) = X(f) + Y(f)$ is again a vector field, since $(X + Y)(fg) = (X + Y)(f) \, g + f \, (X + Y)(g)$.
However, their product $XY$, acting on functions via $XY(f) = X(Y(f))$, is not a vector field which can be seen from the fact that it is not a derivation.
A way we can turn it into a vector field is by instead considering the commutator of $X$ and $Y$; from this, we define the Lie bracket.
\begin{definition}[Lie Bracket]\label{def:lie_bracket}
    Let $M$ be a smooth manifold, and let $X, Y \in \mathfrak{X}(M)$ be vector fields.
    The Lie bracket of $X$ and $Y$, denoted by $[X, Y]$, is defined by its action on a smooth function $f \in C^\infty(M)$ given by
    \begin{equation}
        [X, Y](f) = X(Y(f)) - Y(X(f)).
    \end{equation}
    We can thus write $[X, Y] = XY - YX$.
\end{definition}
The bracket $[X, Y]$ is a derivation, since the cross terms $(Xf)(Yg)$ and $(Yf)(Xg)$ cancel in $[X, Y](fg) = [X, Y](f) \, g + f \, [X, Y](g)$, and it is smooth when $X$ and $Y$ are smooth by the coordinate formula below.
The value of $[X, Y]$ at a point $p \in M$ is given by
\begin{equation}
    [X, Y]_p(f) = X_p(Y(f)) - Y_p(X(f)),
\end{equation}
but a more useful formula using coordinate representations $X = X^i \partial_i$ and $Y = Y^j \partial_j$ is
\begin{equation}
    [X, Y] = (XY^j - YX^j) \frac{\partial}{\partial x^j}.
\end{equation}
From this, it follows that the Lie bracket of any two coordinate vector fields is zero, since their component functions are constant.
This essentially restates the fact that mixed partial derivatives of smooth functions commute.

The Lie bracket is bilinear, anticommutative -- and thus alternating, though this is immediately apparent from its definition -- and it satisfies what is known as the \emph{Jacobi identity}.
\begin{equation}
    [X, [Y, Z]] + [Y, [Z, X]] + [Z, [X, Y]] = 0
\end{equation}
Also, the Lie bracket is, in a sense, \enquote*{natural}.
This can be considered in general via natural transformations of functors, but we present it with the following proposition.
\begin{proposition}\label{prop:naturality_of_lie_bracket}
    Let $M$ and $N$ be smooth manifolds and let $f \colon M \to N$ be a smooth function.
    Let $X_{(1)}, X_{(2)} \in \mathfrak{X}(M)$ and $Y_{(1)}, Y_{(2)} \in \mathfrak{X}(N)$ be vector fields such that $X_{(i)}$ is $f$-related to $Y_{(i)}$ for $i \in \{1, 2\}$.
    Then, $[X_{(1)}, X_{(2)}]$ is $f$-related to $[Y_{(1)}, Y_{(2)}]$.
\end{proposition}
\begin{proof}
    Consider a point $p \in M$ and a smooth function $g \in C^\infty(U, \R)$, where $U$ is a neighbourhood of $f(p)$.
    We have
    \begin{equation*}
        X(g \circ f)(p) = X_p(g \circ f) = Df_p(X_p)(g)
    \end{equation*}
    and
    \begin{equation*}
        (Yg) \circ f(p) = (Yg)(f(p)) = Y_{f(p)}(g).
    \end{equation*}
    Thus, the equation
    \begin{equation}\label{eq:f-related_composition}
        X(g \circ f) = (Yg) \circ f
    \end{equation}
    is true for all $g$ if and only if $Df_p(X_p) = Y_{f(p)}$ for all $p$.
    That is, it is true if and only if $X$ and $Y$ are $f$-related.
    Using \eqref{eq:f-related_composition}, we have
    \begin{equation*}
        X_{(1)}X_{(2)}(g \circ f) = X_{(1)}((Y_{(2)}g) \circ f) = (Y_{(1)}Y_{(2)}g) \circ f
    \end{equation*}
    and, similarly,
    \begin{equation*}
        X_{(2)}X_{(1)}(g \circ f) = (Y_{(2)}Y_{(1)}g) \circ f.
    \end{equation*}
    Therefore,
    \begin{align*}
        [X_{(1)}, X_{(2)}](g \circ f) &= X_{(1)} X_{(2)}(g \circ f) - X_{(2)} X_{(1)} (g \circ f) \\
        &= (Y_{(1)} Y_{(2)} g) \circ f - (Y_{(2)} Y_{(1)} g) \circ f \\
        &= ([Y_{(1)}, Y_{(2)}]g) \circ f. \qedhere
    \end{align*}
\end{proof}
\begin{corollary}\label{cor:pushforward_of_lie_bracket}
    Let $M$ and $N$ be smooth manifolds and let $f \colon M \to N$ be a diffeomorphism.
    Let $X_{(1)}, X_{(2)} \in \mathfrak{X}(M)$ be vector fields.
    Then,
    \begin{equation}
        f_*[X_{(1)}, X_{(2)}] = [f_*X_{(1)}, f_*X_{(2)}].
    \end{equation}
\end{corollary}
\begin{proof}
    By definition, $X_{(i)}$ is $f$-related to $f_*X_{(i)}$.
    \Cref{prop:naturality_of_lie_bracket} then shows that $[X_{(1)}, X_{(2)}]$ is $f$-related to $[f_*X_{(1)}, f_*X_{(2)}]$.
    Since $f$ is a diffeomorphism, the vector field $f$-related to $[X_{(1)}, X_{(2)}]$ is unique, and it is $f_*[X_{(1)}, X_{(2)}]$.
\end{proof}
Vector fields allow us to talk more about curves on manifolds.
Given a smooth curve $\gamma \colon I \to M$, the velocity vector $\gamma'(t)$ is tangent to the manifold for each $t \in I$.
In the other direction, we may be given a vector field $X$ and consider a curve whose velocity is tangent to $X$ at each point.
\begin{definition}[Integral Curve]\label{def:integral_curve}
    Let $M$ be a smooth manifold and let $X \in \mathfrak{X}(M)$ be a vector field.
    An integral curve of $X$ is a differentiable curve $\gamma \colon I \to M$, where $I \subseteq \R$, such that
    \begin{equation}
        \gamma'(t) = X_{\gamma(t)}
    \end{equation}
    for all $t \in I$.
\end{definition}
We can write a smooth curve $\gamma \colon I \to M$ on a smooth coordinate domain $U \subseteq M$ as $\gamma(t) = (\gamma^1(t), \ldots, \gamma^n(t))$.
Then, $\gamma$ is an integral curve if
\begin{equation}
    \dot{\gamma}^i \eval{\frac{\partial}{\partial x^i}}{\gamma(t)} = X^i(\gamma(t)) \eval{\frac{\partial}{\partial x^i}}{\gamma(t)},
\end{equation}
which gives us the following system of ordinary differential equations.
\begin{align}
    \begin{split}
        \dot{\gamma}^1(t) &= X^1(\gamma^1(t), \ldots, \gamma^n(t)) \\
        \hphantom{\dot{\gamma}^1(t)} & \;\: \vdots \\
        \dot{\gamma}^n(t) &= X^n(\gamma^1(t), \ldots, \gamma^n(t))
    \end{split}
\end{align}
That is, finding an integral curve for some vector field involves solving a system of ordinary differential equations, hence giving reason for the name \enquote*{integral curve}.

We may wish to consider how points on a manifold flow along an integral curve.
One may first consider the usual definition of the directional derivative,
\begin{equation}
    \nabla_v f(p) = \lim_{t \to 0} \frac{f(p + tv) - f(p)}{t}.
\end{equation}
However, extending this to vector fields does not really work since $X_{p+tv}$ and $X_p$ are tangent vectors in different tangent spaces; taking their difference thus does not make sense.
A solution is found by defining flows of vector fields.
\begin{definition}[Flow]\label{def:flow}
    Let $M$ be a smooth manifold and let $X \in \mathfrak{X}(M)$ be a vector field.
    Consider an open interval $D$ of $\R$ containing $0$ and an open neighbourhood $W$ of $p \in M$, and let $s, t \in \R$.
    A local flow of $X$ about $p$ is a function $\theta \colon D \times W \to M$ that is continuous and satisfies the rules for a left action of the additive group $\R$ on $M$, wherever they are defined, i.e. by writing $\theta_t(q) = \theta(t, q)$, we have
    \begin{equation}
        \theta_0(q) = q
    \end{equation}
    for all $q \in W$, and
    \begin{equation}
        \theta_t(\theta_s(q)) = \theta_{t+s}(q)
    \end{equation}
    whenever both sides are defined.
    We also require that $t \mapsto \theta_t(q)$ is an integral curve of $X$ for each $q \in W$.
    If $D \times W = \R \times M$, we call $\theta$ a global flow of $X$.
\end{definition}
After time $t$, a point $p$ is moved by the flow $\theta$ of $X$ to the point $\theta_t(p)$.
To be able to compare $Y_{\theta_t(p)}$ to $Y_p$ (be able to subtract them), we use the pushforward $(D(\theta_t)_p)^{-1} = D(\theta_{-t})_{\theta_t(p)}$ to move the flow back to $p$.
This gives us our definition of the Lie derivative.
\begin{definition}[Lie Derivative of Vector Field]\label{def:lie_derivative_of_vector_field}
    Let $M$ be a smooth manifold, and let $X, Y \in \mathfrak{X}(M)$ be vector fields.
    Let $\theta \colon D \times W \to M$ be a local flow of $X$.
    The Lie derivative of $Y$ with respect to $X$ at $p$ is defined by
    \begin{equation}
        (\mathcal{L}_X Y)_p = \lim_{t \to 0} \frac{D(\theta_{-t})_{\theta_t(p)}(Y_{\theta_t(p)}) - Y_p}{t}
    \end{equation}
    when the limit exists.
\end{definition}
In practice, however, we can compute the Lie derivative of a vector field as a Lie bracket.
\begin{proposition}\label{prop:lie_derivative_is_lie_bracket}
    Let $M$ be a smooth manifold and let $X, Y \in \mathfrak{X}(M)$ be vector fields.
    Then,
    \begin{equation}
        \mathcal{L}_X Y = [X, Y].
    \end{equation}
\end{proposition}
\begin{proof}
    See \cite[Thm.~9.38]{Lee2012}.
\end{proof}

\section{Differential Forms}\label{sec:differential_forms}

\subsection{Differential Forms on Manifolds}\label{subsec:differential_forms:differential_forms_on_manifolds}

Differential forms combine the multilinear algebra of \cref{sec:multilinear_algebra} with the tangent bundles of \cref{sec:manifolds}.
A differential $k$-form assigns to each point $p \in M$ an alternating $k$-multilinear map taking tangent vectors to real numbers.
A differential one-form may be considered to be a field of cotangent vectors (covectors), which are linear functionals of tangent vectors.
This is succinctly expressed using the cotangent bundle of \cref{subsec:manifolds:vector_bundles}.
A covector field evaluates to a real number in two steps.
First, it is given a point to turn it into a covector, and then it is given a tangent vector.
However, we are able to define the action of a covector field $\omega$ on a vector field $X$ by
\begin{equation}
    \omega(X)(p) \coloneq \omega_p(X_p).
\end{equation}
A differential one-form assigns to each point $p \in M$ a covector, so it may also be called a covector field.
It is a section of the cotangent bundle.
By instead considering sections of exterior powers of the cotangent bundle, we obtain general differential forms.
\begin{definition}[Differential Form]
    Let $M$ be a smooth $n$-dimensional manifold, and let $T^*M$ be its cotangent bundle with the projection map $\pi$.
    A differential $k$-form on $M$ is a smooth section of the $k$-th exterior power of the cotangent bundle, $\extp^k(T^*M)$, i.e. it is a smooth right inverse of $\pi$.
    The space of all differential $k$-forms on $M$ is denoted by $\Omega^k(M)$ or $\Gamma(\extp^k(T^*M))$.
    The space of all differential forms on $M$ is defined by $\Omega^\bullet(M) = \bigoplus_{k=0}^n \Omega^k(M)$.
\end{definition}
\begin{remark}
    A differential $k$-form is equivalently an alternating type $(0, k)$ tensor field; it assigns to each point an exterior $k$-form (also called a $k$-covector), an element of $\extp^k(T^*_p M)$.
    With the assumption that $M$ is finite-dimensional, the $k$-form is associated with a linear function from $\extp^k T_p M$ to $\R$:
    \begin{equation}
        \extp^k (T^*_p M) \cong (\extp^k T_p M)^*.
    \end{equation}
    That is then naturally isomorphic to the space of alternating $k$-multilinear maps, so $\extp^k T^*_p M \cong \Alt^k(T_p M)$.
    These associations allow us to consider differential forms with both their elegant definition and with easy ways to work with them in computations.
    We often call differential $k$-forms just $k$-forms for convenience.
    Also, it follows from our definition that a zero-form is a smooth function $f \colon M \to \R$.
    A differential $n$-form may be called a top-dimensional form or a top-form.
\end{remark}
We can consider the wedge product $\wedge \colon \Omega^k(M) \times \Omega^l \to \Omega^{k+l}$ between differential forms using \Cref{def:wedge_product_of_forms}.
\begin{align}
    (\omega \wedge \eta)(v_{(1)}, \ldots, v_{(k+l)}) &= \frac{(k+l)!}{k!l!} \alt(\omega \otimes \eta) \\
    &= \frac{1}{k!l!} \sum_{\sigma \in S_{k+l}} \sgn(\sigma) \omega(v_{(\sigma(1))}, \ldots, v_{(\sigma(k))}) \eta(v_{(\sigma(k+1))}, \ldots, v_{(\sigma(k+l))}) \nonumber
\end{align}
This turns $\Omega(M)$ into an associative, alternating graded algebra.
The properties of the wedge product are the same as those of the tensor product with the addition of being alternating; it is associative and distributive over addition.
The wedge product is skew commutative, which is a type of graded commutativity characterized by the equation
\begin{equation}
    \omega \wedge \eta = (-1)^{kl} \eta \wedge \omega.
\end{equation}
In coordinates, we can express a differential one-form as
\begin{equation}
    \omega = \omega_i \d{x}^i,
\end{equation}
and we can express a differential $k$-form as
\begin{equation}
    \omega = \omega_{i_1 \ldots i_k} \d{x}^{i_1} \wedge \cdots \wedge \d{x}^{i_k} = \omega_I \d{x}^I,
\end{equation}
where we use multi-index notation to simplify our equation.
Note that we are summing over increasing multi-indices above only.

Being alternating multilinear maps, differential forms are related to determinants of matrices.
If we consider the differential $k$-form made by the wedge product of $k$ one-forms, it acts on $k$ vectors (or a $k$-vector) to give the determinant of the matrix whose entries are the pairings of the $i$-th one-form with the $j$-th vector.
\begin{equation}
    (\omega^{(1)} \wedge \cdots \wedge \omega^{(k)})(v_{(1)}, \ldots, v_{(k)}) = \begin{vmatrix}
        \omega^{(1)}(v_{(1)}) & \cdots & \omega^{(1)}(v_{(k)}) \\
        \vdots & \ddots & \vdots \\
        \omega^{(k)}(v_{(1)}) & \cdots & \omega^{(k)}(v_{(k)})
    \end{vmatrix}
\end{equation}
Differential forms are covariant tensor fields, so their components transform covariantly (with the inverse Jacobian).
Given a change in coordinates from $x^\mu$ to $\bar{x}^a$, we have the following transformation law for a differential $k$-form.
\begin{equation}
    \bar{\omega}_{\alpha_1 \ldots \alpha_k} = \frac{\partial x^{\mu_1}}{\partial \bar{x}^{\alpha_1}} \cdots \frac{\partial x^{\mu_k}}{\partial \bar{x}^{\alpha_k}} \omega_{\mu_1 \ldots \mu_k}
\end{equation}
In elementary multivariable calculus, we are able to consider various differential operators acting on functions $f$ defined on $\R^n$.
One such operator is the gradient of $f$, a vector field whose components are given by the partial derivatives of $f$.
With the terminology that we have defined, this can be expressed as
\begin{equation}
    \grad f = \sum_{i=1}^n \frac{\partial f}{\partial x^i} \frac{\partial}{\partial x^i}.
\end{equation}
However, this is a coordinate-dependent definition.
It does not work with Einstein summation convention, and it is not clear how it can be made coordinate independent.
From the clue of index conventions, we instead consider a covector field whose components are given by the aforementioned partial derivatives.
We call this the differential of $f$.
\begin{definition}[Differential of Function]\label{def:differential_of_function}
    Let $M$ be a smooth manifold and let $f \colon M \to \R$ be a smooth function.
    The differential of $f$, denoted $\d{f}$, is a differential one-form defined by
    \begin{equation}
        \d{f}_p(v) = v(f),
    \end{equation}
    where $p \in M$ and $v \in T_p M$.
\end{definition}
\begin{remark}
    It immediately follows that the coordinate representation of the differential is
    \begin{equation}
        \d{f}_p = \frac{\partial f}{\partial x^i} \d{x}^i,
    \end{equation}
    justifying our definition as being inspired by our previous discussion.
    It is also worth noting the connection between the differential of a function and the pushforward of the function.
    The codomains of these functions are $\R$ and $T_{f(p)} \R$, respectively, but we know that $\R \cong T_{f(p)} \R$ canonically.
    The tangent vector $Df_p(v) \in T_{f(p)}$ is associated with the scalar $\d{f}_p(v)$, so these two operations are essentially the same thing in the case of functions $f \colon M \to \R$.
\end{remark}
If we take the differential of a coordinate function $x^i$, then $\d{x^i}(\partial_j) = \partial_j x^i = \delta^i_j$, so the differential of $x^i$ equals the basis covector $\d{x}^i$.
This justifies our choice of notation for the latter object.

The differential of a function acts very similarly to the ordinary derivative of a function.
The differential is linear, it satisfies the product, quotient and chain rules, and the differential of a constant is zero.
For the latter fact, it is actually an equivalence: the differential $df = 0$ if and only if $f$ is a locally constant function.

The following proposition is analogous to \Cref{prop:pushforward_of_velocity}, but for real-valued functions and the differential.
\begin{proposition}\label{prop:differential_of_velocity}
    Let $f \in C^\infty(M)$ be a smooth function on a smooth manifold $M$, and let $\gamma \colon I \to M$ be a smooth curve.
    At any $t_0 \in I$,
    \begin{equation}
        (f \circ \gamma)'(t_0) = \d{f}_{\gamma(t_0)}(\gamma'(t_0)).
    \end{equation}
\end{proposition}
\begin{proof}
    By the definition of the velocity and of the pushforward, $\gamma'(t_0)(f) = D\gamma_{t_0}\left(\eval*{\frac{d}{dt}}{t_0}\right)(f) = \eval*{\frac{d}{dt}}{t_0}(f \circ \gamma) = (f \circ \gamma)'(t_0)$.
    By the definition of the differential, the left-hand side is $\d{f}_{\gamma(t_0)}(\gamma'(t_0))$. \qedhere
\end{proof}
An operation that can be performed on differential forms, which makes use of the differential of a function, is the pullback.
As suggested by its name, it is the concept dual to the pushforward, taking a form back to the domain of some map.
\begin{definition}[Pullback of Differential Forms]\label{def:pullback_of_differential_forms}
    Let $M$ and $N$ be smooth manifolds, let $f \colon M \to N$ be a smooth function, and let $\omega \in \Omega^k(N)$ be a differential $k$-form.
    The pullback map $f^* \colon \Omega^k(N) \to \Omega^k(M)$ takes $\omega$ to a map $f^*\omega \colon M \to \extp^k T^* M$ called the pullback of $\omega$ by $f$.
    At a point $p \in M$, this is a map $(f^* \omega)_p \colon (T_p M)^{\times k} \to \R$ acting on $k$ tangent vectors to give
    \begin{equation}
        (f^*\omega)_p(v_{(1)}, \ldots, v_{(k)}) = \omega_{f(p)}(Df_p(v_{(1)}), \ldots, Df_p(v_{(k)})).
    \end{equation}
\end{definition}
We can view the pullback of a form in a more abstract way.
At a point $p \in M$, the pushforward $Df_p \colon T_p M \to T_{f(p)} N$ has a dual map $(Df_p)^* \colon T^*_{f(p)} N \to T^*_p M$.
The pullback of a differential $k$-form $\omega$ by $f$ is then given at $p$ by the exterior power $\extp^k (Df_p)^*$ applied to $\omega_{f(p)}$,
\begin{equation}
    (f^*\omega)_p = \extp^k (Df_p)^* (\omega_{f(p)}).
\end{equation}
There is no dual map defined on the whole of $T^*N$ unless $f$ is injective, as the covectors at $f(p)$ would have to be pulled back to every point of the fibre of $f$ over it.
The pullback is linear, and it distributes over the wedge product.
It is functorial in the sense that
\begin{equation}
    (g \circ f)^* = f^* \circ g^*,
\end{equation}
and
\begin{equation}
    \id^*_M = \id_{\Omega(M)}.
\end{equation}
We can express the pullback of a generic differential $k$-form in coordinates as
\begin{equation}
    f^*\omega = (\omega_I \circ f) \frac{\partial f^I}{\partial x^J} \d{x}^J,
\end{equation}
where $\partial f^I / \partial x^J$ denotes the $k \times k$ minor $\det(\partial f^{i_a}/\partial x^{j_b})$, and the sum is over increasing multi-indices $I$ and $J$.
In the case where $k = m = n$, then we have
\begin{equation}\label{eq:pullback_top-form_coordinates_formula}
    f^*\omega = (\omega_I \circ f) (\det Df) \d{x}^I
\end{equation}
where $\det Df$ is the determinant of the Jacobian matrix of $f$ in these coordinates.
Thus, the pullback can be seen as a change of variables in this case.

\subsection{Exterior Derivative}\label{subsec:differential_forms:exterior_derivative}

The differential turns a function into a one-form.
This can be considered a special case of an operation that takes a $k$-form to a $(k+1)$-form.
The exterior derivative generalizes the differential of a function and thus generalizes the regular derivative.
It is an example of an antiderivation, which is a special type of graded derivation.
\begin{definition}[Graded Derivation]\label{def:graded_derivation}
    Let $A$ be a graded algebra.
    A linear map $D \colon A \to A$ is a graded derivation on $A$ of grade or degree $\abs{D}$ if
    \begin{equation}
        D(ab) = D(a)b + (-1)^{\abs{a}\abs{D}}aD(b)
    \end{equation}
    for all $a, b \in A$.
    If $\abs{D}$ is odd, then $D$ is called an antiderivation.
\end{definition}
\begin{remark}
    Note that the degree of a graded operator is determined by the way it affects the grade of a particular element.
    That is, if $a \in A$ is of grade $\abs{a}$, then $D(a)$ is of grade $\abs{a} + \abs{D}$.
\end{remark}
\begin{definition}[Exterior Derivative]\label{def:exterior_derivative}
    Let $M$ be a smooth manifold.
    The exterior derivative of a differential $k$-form on $M$ is a degree $+1$ antiderivation $\d \colon \Omega^k(M) \to \Omega^{k+1}(M)$ such that $\d^2 = 0$, and it agrees with the differential of a function when $k = 0$.
\end{definition}
\begin{remark}
    The exterior derivative exists and is unique on any smooth manifold, and it does not depend on a choice of charts \cite[Thm.~14.24]{Lee2012}.
\end{remark}

Our requirement that the exterior derivative is a degree $+1$ antiderivation amounts to the following graded product rule.
Here, $\omega$ is a $k$-form and $\eta$ is an $l$-form.
\begin{equation}
    \d(\omega \wedge \eta) = \d{\omega} \wedge \eta + (-1)^k \omega \wedge \d{\eta}
\end{equation}
The exterior derivative commutes with the pullback of a form by any smooth map.
\begin{proposition}[Naturality of the Exterior Derivative]\label{prop:naturality_of_exterior_derivative}
    Let $f \colon M \to N$ be a smooth map between smooth manifolds, and let $\omega \in \Omega^k(N)$.
    Then
    \begin{equation}
        \d(f^*\omega) = f^*\d{\omega}.
    \end{equation}
\end{proposition}
\begin{proof}
    Both sides are local, so we may take $\omega = \omega_I \d{y}^{i_1} \wedge \cdots \wedge \d{y}^{i_k}$ in a chart of $N$.
    For a function $g$ on $N$, $(f^*\d{g})_p(v) = \d{g}_{f(p)}(Df_p(v)) = Df_p(v)(g) = v(g \circ f) = \d(g \circ f)_p(v)$, so $f^*\d{g} = \d(g \circ f)$.
    Since $f^*$ preserves wedge products, $f^*\omega = (\omega_I \circ f) \, \d(y^{i_1} \circ f) \wedge \cdots \wedge \d(y^{i_k} \circ f)$.
    The exterior derivative is an antiderivation with $\d^2 = 0$, so the factors $\d(y^{i_j} \circ f)$ are closed and
    \begin{equation*}
        \d(f^*\omega) = \d(\omega_I \circ f) \wedge \d(y^{i_1} \circ f) \wedge \cdots \wedge \d(y^{i_k} \circ f) = f^*(\d{\omega_I} \wedge \d{y}^{i_1} \wedge \cdots \wedge \d{y}^{i_k}) = f^*\d{\omega}. \qedhere
    \end{equation*}
\end{proof}
This property is called the naturality of the exterior derivative, and Palais \cite{palais1959natural} shows that the exterior derivative defined using its naturality is unique up to a constant multiple.

Given a differential form $\omega = \omega_I \d{x}^I$ expressed in local coordinates, the exterior derivative is given by
\begin{equation}
    \d{\omega} = \d{\omega_I} \wedge \d{x^I} + \omega_I \d{\d{x^I}} = \d{\omega_I} \wedge \d{x^I} = \frac{\partial \omega_I}{\partial x^j} \d{x^j} \wedge \d{x^I}.
\end{equation}
We can also provide a formula for the exterior derivative of a differential $k$-form that does not depend on any choice of coordinates.
To do this, we will need both the interior product and the Lie derivative of a differential form.
As opposed to the wedge product or exterior product, which increases the degree of a differential form, the interior product decreases the degree by pairing it with a vector field.
\begin{definition}[Interior Product]\label{def:interior_product}
    Let $M$ be a smooth manifold, and let $\omega \in \Omega^k(M)$ be a differential $k$-form.
    The interior product with a vector field $V \in \mathfrak{X}(M)$ is a degree $-1$ antiderivation $\iota_V \colon \Omega^k(M) \to \Omega^{k-1}$ defined by
    \begin{equation}
        (\iota_V \omega)(X_{(1)}, \ldots, X_{(k-1)}) \coloneq \omega(V, X_{(1)}, \ldots, X_{(k-1)}).
    \end{equation}
    For $k = 0$, $\iota_V \omega$ is defined to be zero.
\end{definition}
\begin{remark}
    When $k = 1$, the interior product of $\omega$ with $X$ is the duality pairing between the differential form and the vector field.
    \begin{equation}
        \iota_X \omega = \omega(X)
    \end{equation}
\end{remark}
Being an antiderivation, the interior product satisfies
\begin{equation}
    \iota_X (\omega \wedge \eta) = (\iota_X \omega) \wedge \eta + (-1)^k \omega \wedge (\iota_X \eta)
\end{equation}
for $\omega \in \Omega^k(M)$ and $\eta \in \Omega^l(M)$.

The Lie derivative of a differential form may be defined similarly to that of a vector field found in \Cref{def:lie_derivative_of_vector_field}.
Instead of pushing a vector field forward, we pull a $k$-form back.
\begin{definition}[Lie Derivative of Differential Form]\label{def:lie_derivative_of_differential_form}
    Let $M$ be a smooth manifold, let $X \in \mathfrak{X}(M)$ be a vector field, and let $\omega \in \Omega(M)$ be a differential $k$-form.
    Let $\theta \colon D \times W \to M$ be a local flow of $X$.
    The Lie derivative of $\omega$ with respect to $X$ at $p$ is defined by
    \begin{equation}
        (\mathcal{L}_X \omega)_p = \lim_{t \to 0} \frac{\theta^*_t(\omega_{\theta_t(p)}) - \omega_p}{t}
    \end{equation}
    when the limit exists.
\end{definition}
We then have some relationships between the interior product, Lie derivative, and the exterior derivative.
The interior product with respect to their Lie bracket $[X, Y]$ satisfies
\begin{equation}
    \iota_{[X, Y]} = [\mathcal{L}_X, \iota_Y] = [\iota_X, \mathcal{L}_Y].
\end{equation}
A more important result that relates the three quantities, which can be taken as the definition of the Lie derivative of a differential form, is named after \'{E}lie Cartan.
\begin{theorem}[Cartan's Magic Formula]\label{thm:cartan's_magic_formula}
    Let $M$ be a smooth manifold, let $X \in \mathfrak{X}(M)$ be a vector field.
    Then,
    \begin{equation}\label{eq:cartan's_magic_formula}
        \mathcal{L}_X = \d \circ \iota_X + \iota_X \circ \d.
    \end{equation}
\end{theorem}
\begin{proof}
    Let $P = \d \circ \iota_X + \iota_X \circ \d$.

    First, we show that $\mathcal{L}_X$ and $P$ are both derivations of degree zero which commute with $\d$.
    Being the anticommutator of two antiderivations, $P$ is a derivation: expanding $P(\alpha \wedge \beta)$ with the antiderivation rules leaves $P\alpha \wedge \beta + \alpha \wedge P\beta$, as the terms $\iota_X\alpha \wedge \d{\beta}$ and $\d{\alpha} \wedge \iota_X\beta$ each occur twice with opposite signs.
    It commutes with $\d$, since $\d P = \d \iota_X \d = P \d$ by $\d^2 = 0$.
    For $\mathcal{L}_X$, differentiating $\theta_t^*(\alpha \wedge \beta) = \theta_t^*\alpha \wedge \theta_t^*\beta$ at $t = 0$ gives the product rule.
    Differentiating $\d(\theta_t^*\omega) = \theta_t^*\d{\omega}$ at $t = 0$ shows that it commutes with $\d$, which uses \cref{prop:naturality_of_exterior_derivative} and the smoothness of $(t, p) \mapsto (\theta_t^*\omega)_p$.

    Next, we show that $\mathcal{L}_X f = Pf$ for every smooth function $f$.
    Since $t \mapsto \theta_t(p)$ is an integral curve of $X$, \cref{prop:differential_of_velocity} gives $(\mathcal{L}_X f)(p) = \eval*{\frac{d}{dt}}{0} f(\theta_t(p)) = \d{f}_p(X_p) = (Xf)(p)$.
    On the other hand, $Pf = \iota_X\d{f} = \d{f}(X) = Xf$.

    Finally, we show that $\mathcal{L}_X \omega = P\omega$ for every differential form $\omega$.
    Both operators are local, so we may work in a chart, where $\omega$ is a sum of forms $g \, \d{x}^{i_1} \wedge \cdots \wedge \d{x}^{i_k}$.
    Being derivations, it suffices for $\mathcal{L}_X$ and $P$ to agree on functions and on the one-forms $\d{x}^i$.
    They agree on functions as shown above, and on $\d{x}^i$ because $\mathcal{L}_X \d{x}^i = \d(\mathcal{L}_X x^i) = \d(P x^i) = P\d{x}^i$. \qedhere
\end{proof}
\begin{corollary}\label{cor:lie_derivative_commutes_with_exterior_derivative}
    Let $M$ be a smooth manifold, $X \in \mathfrak{X}(M)$ be a vector field, and let $\omega \in \Omega^k(M)$ be a differential $k$-form.
    Then,
    \begin{equation}
        \mathcal{L}_X(\d{\omega}) = \d(\mathcal{L}_X \omega).
    \end{equation}
\end{corollary}
\begin{proof}
    By \cref{thm:cartan's_magic_formula} and $\d^2 = 0$, both sides equal $\d \iota_X \d{\omega}$.
\end{proof}
We first need a formula for the Lie derivative of a form evaluated on vector fields.
\begin{lemma}\label{lem:lie_derivative_of_form_evaluated_on_vector_fields}
    Let $V, X_{(1)}, \ldots, X_{(k)} \in \mathfrak{X}(M)$ and let $\omega \in \Omega^k(M)$.
    Then
    \begin{equation}
        (\mathcal{L}_V \omega)(X_{(1)}, \ldots, X_{(k)}) = V(\omega(X_{(1)}, \ldots, X_{(k)})) - \sum_{i=1}^k \omega(X_{(1)}, \ldots, [V, X_{(i)}], \ldots, X_{(k)}).
    \end{equation}
\end{lemma}
\begin{proof}
    Let $\theta$ be the flow of $V$.
    By the definitions of the pullback and the pushforward, at each point $p$,
    \begin{equation*}
        \omega(X_{(1)}, \ldots, X_{(k)})(\theta_t(p)) = (\theta_t^*\omega)_p\bigl(D(\theta_{-t})(X_{(1)}), \ldots, D(\theta_{-t})(X_{(k)})\bigr),
    \end{equation*}
    where the vector fields are evaluated at $\theta_t(p)$ as in \cref{def:lie_derivative_of_vector_field}.
    Differentiating at $t = 0$ with the product rule for the multilinear evaluation gives
    \begin{equation*}
        V(\omega(X_{(1)}, \ldots, X_{(k)})) = (\mathcal{L}_V\omega)(X_{(1)}, \ldots, X_{(k)}) + \sum_{i=1}^k \omega(X_{(1)}, \ldots, \mathcal{L}_V X_{(i)}, \ldots, X_{(k)}),
    \end{equation*}
    and $\mathcal{L}_V X_{(i)} = [V, X_{(i)}]$ by \cref{prop:lie_derivative_is_lie_bracket}. \qedhere
\end{proof}
With the interior product and Lie derivative, we are able to prove an invariant formula for the exterior derivative of a differential form.
\begin{proposition}\label{prop:exterior_derivative_invariant_formula}
    Let $M$ be a smooth manifold, and let $\omega \in \Omega^k(M)$ be a differential $k$-form.
    Given $k+1$ smooth vector fields $X_{(0)}, \ldots, X_{(k)}$, the exterior derivative of $\omega$ acts on the vector fields to give
\begin{align}\label{eq:exterior_derivative_invariant_formula}
    \begin{split}
            \d{\omega}(X_{(0)}, \ldots, X_{(k)}) &= \sum_i (-1)^i X_{(i)}(\omega(X_{(0)}, \ldots, \widehat{X}_{(i)}, \ldots, X_{(k)})) \\
    &+ \sum_{i<j} (-1)^{i+j} \omega([X_{(i)}, X_{(j)}], X_{(0)}, \ldots, \widehat{X}_{(i)}, \ldots, \widehat{X}_{(j)}, \ldots, X_{(k)}),
    \end{split}
\end{align}
where a hat above a symbol denotes it omission.
\end{proposition}
\begin{proof}
    We prove the formula by induction on $k$.

    For $k = 0$, we have $\d{f}(X_{(0)}) = X_{(0)}f$ for a function $f$, which is \eqref{eq:exterior_derivative_invariant_formula} in this case.

    Now suppose that $k \geq 1$ and that the formula holds for $(k-1)$-forms.
    By the definition of the interior product and \cref{thm:cartan's_magic_formula},
    \begin{align*}
        \d{\omega}(X_{(0)}, \ldots, X_{(k)}) &= (\iota_{X_{(0)}}\d{\omega})(X_{(1)}, \ldots, X_{(k)}) \\
        &= (\mathcal{L}_{X_{(0)}} \omega)(X_{(1)}, \ldots, X_{(k)}) - \d(\iota_{X_{(0)}}\omega)(X_{(1)}, \ldots, X_{(k)}).
    \end{align*}
    The first term is given by \cref{lem:lie_derivative_of_form_evaluated_on_vector_fields} with $V = X_{(0)}$.
    For the second term, apply the induction hypothesis to the $(k-1)$-form $\eta = \iota_{X_{(0)}}\omega$ and the vector fields $X_{(1)}, \ldots, X_{(k)}$, noting that $\eta(Y_{(1)}, \ldots) = \omega(X_{(0)}, Y_{(1)}, \ldots)$.
    Together, this gives
    \begin{align*}
        \d{\omega}(X_{(0)}, \ldots, X_{(k)}) &= X_{(0)}(\omega(X_{(1)}, \ldots, X_{(k)})) - \sum_{i=1}^k \omega(X_{(1)}, \ldots, [X_{(0)}, X_{(i)}], \ldots, X_{(k)}) \\
        &\quad - \sum_{i=1}^k (-1)^{i-1} X_{(i)}(\omega(X_{(0)}, \ldots, \widehat{X}_{(i)}, \ldots, X_{(k)})) \\
        &\quad - \sum_{1 \leq i < j \leq k} (-1)^{i+j} \omega(X_{(0)}, [X_{(i)}, X_{(j)}], X_{(1)}, \ldots, \widehat{X}_{(i)}, \ldots, \widehat{X}_{(j)}, \ldots, X_{(k)}).
    \end{align*}
    We now match these terms with those of \eqref{eq:exterior_derivative_invariant_formula}.
    The first term and the third line together give $\sum_{i=0}^k (-1)^i X_{(i)}(\omega(X_{(0)}, \ldots, \widehat{X}_{(i)}, \ldots, X_{(k)}))$.
    Moving $[X_{(0)}, X_{(i)}]$ to the first slot passes $i - 1$ arguments, so the second term equals $\sum_{i=1}^k (-1)^i \omega([X_{(0)}, X_{(i)}], X_{(1)}, \ldots, \widehat{X}_{(i)}, \ldots, X_{(k)})$, which is the part of the second sum with $i = 0$.
    Swapping the first two arguments in the last line changes its sign, so it equals $\sum_{1 \leq i < j \leq k} (-1)^{i+j} \omega([X_{(i)}, X_{(j)}], X_{(0)}, \ldots, \widehat{X}_{(i)}, \ldots, \widehat{X}_{(j)}, \ldots, X_{(k)})$, which is the part of the second sum with $i \geq 1$. \qedhere
\end{proof}
There are ways we can use the exterior derivative to classify differential forms.
These classifications are of interest to us in many of the applications of forms.
\begin{definition}[Closed and Exact Differential Forms]\label{def:closed_and_exact_differential_forms}
    Let $M$ be a smooth manifold.
    A differential $k$-form $\omega \in \Omega^k(M)$ is closed if $\d{\omega} = 0$.
    It is exact if $\omega = \d{\eta}$ for some differential $(k-1)$-form $\eta \in \Omega^{k-1}(M)$.
    An exact zero-form is defined to be zero.
\end{definition}
An immediate observation is that every exact form is closed via the fact that $\d^2 = 0$.
A counterexample to the converse of this statement is the following one-form $\d{\theta}$ defined on $\R^2 \setminus \{0\}$.
\begin{equation}\label{eq:angle-form}
    \d{\theta} = \frac{-y \d{x} + x \d{y}}{x^2 + y^2}
\end{equation}
The derivative of this form is zero, so it is closed.
However, there exists no smooth function of which it is the derivative -- the specification of smoothness is the important detail here.
If we were to instead define $\d{\theta}$ on the slit plane $\R^2 \setminus \{(x, 0) : x \leq 0\}$, then $\d{\theta}$ is the exterior derivative of the argument function $\theta$ with values in $(-\pi, \pi)$.
The key difference between $\R^2 \setminus \{0\}$ and the slit plane is that the former contains a hole; this leads to not all closed forms being exact on $\R^2 \setminus \{0\}$, while they are for spaces such as $\R^n$ and $\mathbb{H}^n$.
A more detailed discussion of this topic leads to the study of de Rham cohomology.

As we had previously discussed, the gradient of a function is related to its differential in that they have the same components with respect to some basis.
The gradient takes functions to vectors, whilst the differential takes a function to a one-form.
Without additional structure on our manifold, we cannot go further with this comparison.
We can make similar comparisons to the other differential operators of multivariable calculus.
The divergence of a vector field $X$ on $\R^n$ is given by
\begin{equation}
    \divg X = \sum_{i=1}^n \frac{\partial X^i}{\partial x^i}.
\end{equation}
Thus, the divergence operator takes vector fields to smooth functions.
We can consider the $(n-1)$-form $\omega_V$ corresponding to $X$ given by
\begin{equation}
    \omega_V = \sum_{i=1}^n (-1)^{i-1} X^i \d{x}^1 \wedge \cdots \wedge \widehat{\d{x}^i} \wedge \cdots \wedge \d{x}^n,
\end{equation}
where a hat above a symbol denotes its omission.
The exterior derivative of this has the divergence of $X$ as its components.

When $n = 3$, we can also consider the curl of a vector field, given by
\begin{equation}
    \curl X = \left(\frac{\partial X^3}{\partial x^2} - \frac{\partial X^2}{\partial x^3}\right) \frac{\partial}{\partial x^1} + \left(\frac{\partial X^1}{\partial x^3} - \frac{\partial X^3}{\partial x^1}\right) \frac{\partial}{\partial x^2} + \left(\frac{\partial X^2}{\partial x^1} - \frac{\partial X^1}{\partial x^2}\right) \frac{\partial}{\partial x^3}.
\end{equation}
This takes a vector field to a vector field.
Our corresponding one-form $\eta_V$ is expressed as
\begin{equation}
    \eta_V = \sum_{i=1}^n X^i \d{x}^i,
\end{equation}
and the exterior derivative $\d{\eta}_V$ is a two-form whose components are those of the curl of $X$.
When we cover Riemannian manifolds and Hodge duality, we will be able to provide a connection between these operators and the exterior derivative in a more enlightening way.

\subsection{Integration on Manifolds}\label{subsec:differential_forms:integration_on_manifolds}

The theory of integration is initially developed with the definition of the Riemann integral, allowing one to integrate over $n$-dimensional rectangles.
This may be generalized to the integration of more general spaces using measure theory and the Lebesgue integral.
An alternative way to generalize integration is to consider it on manifolds.
It happens to be the case that functions are not the natural object to integrate over a manifold, but rather it is differential forms that we should consider.

\subsubsection{Orientation of Manifolds}\label{subsubsec:differential_forms:integration_on_manifolds:orientation_of_manifolds}

We first consider orientations, which are important for our theory of integration.
On an $n$-dimensional real vector space $V$, an orientation is given by an equivalence class of ordered bases for $V$.
If $(e_1, \ldots e_n)$ is an ordered basis, then it determines the orientation $[e_1, \ldots, e_n]$ where two ordered bases are equivalent if they are consistently ordered, i.e. the transition matrix between $(e_1, \ldots, e_n)$ and $(\tilde{e}_1, \ldots, \tilde{e}_n)$, defined by
\begin{equation}
   e_i = A^j_i \tilde{e}_j,
\end{equation}
has positive determinant.
If $V$ is zero dimensional, then its orientation is the choice between the numbers $+1$ and $-1$.
An oriented vector space is a vector space with a choice of orientation.
Given an oriented vector space, an ordered basis with the same orientation is called positively oriented; otherwise, it is negatively oriented and denoted by $-[e_1, \ldots, e_n]$.
This is meaningful because there are only two orientation equivalence classes for any vector space.

An orientation for a vector space may also be specified by a nonzero exterior $n$-form $\omega \in \extp^n(V^*)$.
If $n \geq 1$, then the orientation determined by $\omega$ is $[e_1, \ldots, e_n]$ such that $\omega(e_1, \ldots, e_n) > 0$.
If $n = 0$, then the orientation is $+1$ if $\omega > 0$ and $-1$ if $\omega < 0$.
It follows that two $n$-forms determine the same orientation on $V$ if and only if they are positive multiples of each other.
An $n$-form, with respect to an oriented vector space $V$, is called positively oriented if it determines the orientation of $V$.

We can extend the idea of orientation to manifolds.
\begin{definition}[Orientation of Manifold]\label{def:orientation_of_manifold}
    Let $M$ be a smooth manifold.
    A pointwise orientation of $M$ is a choice of orientation on each of its tangent spaces.
    If $(E_{(i)})_{i \in I}$ is a local frame for $TM$ with respect to some open subset $U \subseteq M$, then $(E_{(i)})_{i \in I}$ is positively oriented if $(\eval*{E_{(i)}}{p})_{i \in I}$ is a positively-oriented basis for $T_p M$ at each point $p \in U$.
    A negatively-oriented frame is similarly defined.
    We may denote the orientation at each tangent space of a pointwise orientation $\mathcal{O}$ by $\mathcal{O}_p$.

    A pointwise orientation $\mathcal{O}$ of $M$ is continuous at a point $p \in M$ if there exists a neighbourhood $U$ for $p$ for which the orientation is represented by a (continuous by definition) local frame, i.e. $\mathcal{O}_p = [\eval*{E_{(i)}}{p}]_{i \in I}$ for (continuous by definition) vector fields $E_{(i)}$.
    It is continuous on $M$ if it is continuous at all points $p \in M$.
    An orientation of $M$ is a continuous pointwise orientation of $M$.
    If $M$ is zero dimensional, then its orientation is choice of $\pm 1$ for each point.
    A manifold $M$ is orientable if there exists an orientation for it.
    An oriented manifold is the ordered $(M, \mathcal{O})$ of a manifold with a chosen orientation.
\end{definition}
An orientation of $M$ is determined by a choice of a nonvanishing differential $n$-form $\omega \in \Omega^n(M)$ called an orientation form or a volume form.
If $\omega$ determines the orientation of an oriented manifold, then it is said to be positively oriented.
Two positively-oriented forms $\omega$ and $\tilde{\omega}$ are related by $\tilde{\omega} = f\omega$ for some strictly-positive smooth function $f \colon M \to \R$.
In the case where $n = 0$ and $\omega$ is a smooth real-valued function, it assigns $+1$ to points where it is positive and $-1$ to points where it is negative.

We may pull a differential back via some smooth function and consider how the orientation determined by the form is affected.
\begin{definition}[Orientation-Preserving and -Reversing Functions]
    Let $M$ and $N$ be oriented smooth $n$-dimensional manifolds, and let $\omega \in \Omega^n(M)$ be a positively-oriented volume form.
    A local diffeomorphism $f \colon N \to M$ is orientation preserving if $f^*\omega$ is a positively-oriented volume form on $N$.
    If $-f^*\omega$ is positively oriented, then it is orientation reversing.
\end{definition}

Orientation may be defined for coordinate charts and for atlases.
This provides an alternative way to specify an orientation of a manifold.
\begin{definition}[Orientation of Charts and Atlases]\label{def:orientation_of_charts_and_atlases}
    Let $M$ be a smooth manifold.
    A smooth chart on an oriented manifold is positively oriented if the coordinate frame $(\partial/\partial x^i)_{i \in I}$ is positively oriented, and the chart is negatively oriented if the coordinate frame is negatively oriented.
    A smooth atlas $\{(U_i, \varphi_i)\}_{i \in I}$ is consistently oriented if each transition map $\varphi_j \circ \varphi_i^{-1}$ has positive Jacobian determinant everywhere on $\varphi_i(U_i \cap U_j)$.
\end{definition}
The following is a connection between smooth atlases and smooth charts.
Given a smooth manifold $M$, a consistently-oriented smooth atlas for $M$ induces a unique orientation of $M$ such that every chart in the given atlas is positively oriented.
This justifies the name of a consistently-oriented atlas.
Conversely, if $M$ is oriented and $\partial M = \varnothing$ or $\dim M > 1$, then the collection of all oriented smooth charts is a consistently-oriented atlas for $M$.

We can consider the equivalence class of consistently-oriented smooth atlases with two such atlases being equivalent if all transition maps between their charts have positive Jacobian determinant.
There is a one-to-one relation between these equivalence classes and orientations of $M$, so one can specify an orientation of $M$ by specifying an equivalence class of consistently-oriented smooth atlases.
We may denote a manifold with the opposite orientation of some oriented manifold $M$ by $-M$ and denote the atlas corresponding to $\{U_i, x^1_i, \ldots, x^n_i\}_{i \in I}$ by $\{(U_i, -x^1_i, x^2_i \ldots, x^n_i)\}_{i \in I}$.

We can consider orientations of manifolds with boundary.
The boundary of an oriented manifold is itself an oriented manifold.
For this, we define inward- and outward-pointing vector fields.
\begin{definition}[Inward- and Outward-Pointing Vector Fields]\label{def:inward-_and_outward-pointing_vector_fields}
    Let $M$ be a smooth manifold with boundary, and let $p \in M$ be a point.
    A tangent vector $v \in T_p M$ is inward pointing if $v \notin T_p \partial M$ and there exist a positive real number $\varepsilon > 0$ and a curve $\gamma \colon [0, \varepsilon) \to M$ such that $\gamma(0) = p$, $\gamma([0, \varepsilon)) \subseteq M$, and $\gamma'(0) = v$.
    It is outward pointing if the domain of $\gamma$ is instead $(-\varepsilon, 0]$.
    Inward- and outward-pointing vector fields are defined analogously.
\end{definition}
\begin{remark}
    A tangent vector (field) is outward pointing if and only if its negative is inward pointing and vice versa.
    A tangent vector $v$ is outward pointing if and only if $v^n < 0$, i.e. outward-pointing tangent vectors are those whose $x^n$-components are negative.
\end{remark}
\begin{definition}[Boundary Orientation]\label{def:boundary_orientation}
    Let $M$ be an oriented smooth manifold with boundary of dimension greater than zero.
    The boundary orientation is the orientation of $\partial M$ for which an ordered basis $(E_1, \ldots, E_{n-1})$ of $T_p \partial M$ is positively oriented if and only if $(N, E_1, \ldots, E_{n-1})$ is a positively-oriented basis of $T_p M$, where $N$ is any outward-pointing vector at $p$.
    If $n = 1$, then $\partial M$ is zero dimensional, and a point $p$ is positively oriented if $N$ is positively oriented and negatively oriented otherwise.
\end{definition}

\subsubsection{Integration of Differential Forms}\label{subsubsec:differential_forms:integration_on_manifolds:integration_of_differential_forms}

First, we will consider the integration of differential forms on subsets of $\R^n$.
A domain of integration in $\R^n$ is a bounded subset whose boundary has measure zero.
Let $D \subseteq \R^n$ be a domain of integration, and let $\omega \in \Omega^n(\overbar{D})$ be a differential $n$-form.
If $\omega = f \d{x}^1 \wedge \cdots \wedge \d{x}^n$ for some continuous function $f \colon \overbar{D} \to \R$, the integral of $\omega$ over $D$ is defined to be
\begin{equation}
    \int_D \omega \coloneq \int_D f \d V.\footnote{Note that we use an upright \enquote*{$\d$} for all integrals in this text, including Riemann integrals, for consistency.}
\end{equation}
This could be seen as getting rid of the wedge products to obtain a Riemann integral.
More generally, for an open subset $U$ of $\R^n$ or $\mathbb{H}^n$, the integral of a compactly-supported differential $n$-form $\omega$ over $U$ is defined by
\begin{equation}
    \int_U \omega \coloneq \int_D \omega,
\end{equation}
where $D$ is a domain of integration in $\R^n$ or $\mathbb{H}^n$ containing $\supp \omega$.
This definition is independent of the choice of domain of integration: if $D$ and $D'$ both contain $\supp \omega$, then $f$ vanishes outside $D \cap D'$, so $\int_D f \d V = \int_{D \cap D'} f \d V = \int_{D'} f \d V$.
Our definition is sensible since it naturally transforms via pullbacks.
Given open subsets $U$ and $V$ of $\R^n$ or $\mathbb{H}^n$ and a diffeomorphism $f \colon U \to V$, we have
\begin{equation}
    \int_V \omega = \pm \int_U f^* \omega,
\end{equation}
where $\omega$ is compactly-supported on $V$ and the $\pm$ sign is positive if $f$ is orientation preserving and negative if $f$ is orientation reversing.

Now, we may consider smooth manifolds.
We first define the integral of a differential form over a single coordinate chart on which it is compactly supported.
Then, we use partitions of unity to define the integral on a manifold in general.
\begin{definition}[Integral of Differential Form]\label{def:integral_of_differential_form}
    Let $M$ be an oriented smooth manifold.
    Let $\omega \in \Omega^n(M)$ be a differential $n$-form with compact support in the domain of an oriented smooth coordinate chart $(U, \varphi)$ of $M$.
    The integral of $\omega$ over $U$ is defined to be
    \begin{equation}
        \int_U \omega \coloneq \pm \int_{\varphi(U)} (\varphi^{-1})^* \omega,
    \end{equation}
    where the $\pm$ sign is positive or negative depending on whether the chart is positively or negatively oriented.
    Now, suppose that $\omega$ is compactly supported on $M$.
    Let $\{U_i\}_{i \in I}$ be a finite open cover of $\supp \omega$ by oriented smooth charts, and let $\{\psi_i\}_{i \in I}$ be a subordinate smooth partition of unity, which exists by \cite[Thm.~2.23]{Lee2012}.
    The integral of $\omega$ over $M$ is defined to be
    \begin{equation}
        \int_M \omega = \sum_i \int_M \psi_i \omega.
    \end{equation}
    If the manifold is zero dimensional, then the integral of a compactly-supported zero-form $f \in \Omega^0(M)$ over $M$ is defined to be
    \begin{equation}
        \int_M f = \sum_i f(p_i) - \sum_j f(q_j),
    \end{equation}
    where $p_i$ are points where the orientation is positive and $q_j$ are points where it is negative.
\end{definition}
Each term in the integral is well-defined from the fact that each $\psi_i \omega$ is compactly supported on $U_i$.
The integral over a smooth chart does not depend on our choice of oriented coordinates on the chart: two such choices differ by an orientation-preserving diffeomorphism, to which the change-of-variables formula above applies.
Similarly, the integral over a manifold does not depend on our choice of partition of unity or open cover: if $\{\psi_i\}$ and $\{\tilde{\psi}_j\}$ are two choices, then both resulting sums equal $\sum_{i, j} \int_M \psi_i \tilde{\psi}_j \omega$, since each term is supported in the intersection of a chart from each cover, where the chart independence applies.
This means that the integrals are well-defined.

If $S \subseteq M$ is an oriented immersed $k$-dimensional submanifold of $M$, and $\omega \in \Omega^k(M)$ is a differential $k$-form, then the integral $\int_S \omega$ is interpreted as $\int_S \iota^*_S \omega$, where $\iota_S \colon S \hookrightarrow M$ is the inclusion map.
If $S$ is the boundary $\partial M$, then $\int_{\partial M} \omega$ is unambiguously interpreted as $\int_{\partial M} \iota^*_{\partial M} \omega$, with $\partial M$ being assumed to have the induced orientation.
Also, if $\partial M = \varnothing$, then we interpret $\int_{\partial M} \omega$ as zero.

There are some basic properties of integration of differential forms.
First, integration is linear.
\begin{equation}
    \int_M \lambda \omega + \mu \eta = \lambda \int_M \omega + \mu \int_M \eta
\end{equation}
Reversing the orientation of the manifold $M$ switches the sign of the integral.
\begin{equation}
    \int_{-M} \omega = - \int_M \omega
\end{equation}
If $\omega$ is a positively-oriented form with compact support and $M \neq \varnothing$, then $\int_M \omega > 0$.
Integration is invariant under diffeomorphisms.
That is, given a diffeomorphism $f \colon N \to M$, we have
\begin{equation}
    \int_M \omega = \pm \int_N f^* \omega,
\end{equation}
with the $\pm$ symbol being a positive or negative sign depending on whether $f$ preserves or reverses orientation.
This can be seen as a generalization of the ordinary change-of-variables formula for integration, with the pullback $f^*\omega$ being given by \eqref{eq:pullback_top-form_coordinates_formula} in coordinates.

To actually compute an integral of a differential form, partitions of unity are not particularly useful.
Instead, we split up a manifold into domains of integration and compute the integral on those domains via parametrizations.
Suppose we have $k$ domains of integration $D_1, \ldots, D_k$ in $\R^n$ along with $k$ smooth maps $f_i \colon \overbar{D}_i \to M$ such that
\begin{enumerate}
    \item $f_i$ restricts to an orientation-preserving diffeomorphism from $D_i$ to an open subset $W_i \subseteq M$;
    \item $W_i \cap W_j = \varnothing$ when $i \neq j$; and
    \item $\supp \omega \subseteq \overbar{W}_1 \cup \cdots \cup \overbar{W}_k$.
\end{enumerate}
Then, we can express the integral of $\omega$ over $M$ as
\begin{equation}
    \int_M \omega = \sum_{i=1}^k \int_{D_i} f^*_i \omega.
\end{equation}
We can elucidate the connection between differential forms and the differential operators found in multivariable calculus by the considering the important integral theorems found in such studies.
These theorems all involve integrating some differential operator over a manifold and equating it to an integral over the boundary of the manifold.
With the theory of differential forms, we can generalize these theorems into one result.
\begin{theorem}[Stokes' Theorem]\label{thm:stokes'_theorem}
    Let $M$ be an oriented smooth $n$-dimensional manifold with or without boundary, and let $\omega \in \Omega^{n-1}(M)$ be a differential $(n-1)$-form with compact support on $M$.
    Then,
    \begin{equation}\label{eq:stokes'_theorem}
        \int_M \d{\omega} = \int_{\partial M} \omega.
    \end{equation}
\end{theorem}
\begin{proof}
    Consider an atlas $\{(U_i, \varphi_i)\}_{i \in I}$ for $M$ in which each $U_i$ is diffeomorphic to either $\R^n$ or $\mathbb{H}^n$ via an orientation-preserving diffeomorphism, which for $n \geq 2$ we can always do since any open disk is diffeomorphic to $\R^n$ and any half-disk containing its boundary diameter is diffeomorphic to $\mathbb{H}^n$, after composing with the reflection $x^1 \mapsto -x^1$ if necessary.
    If $n = 1$, we also allow charts onto $\mathbb{H}^1$ that reverse orientation, as an orientation-preserving chart need not exist near a boundary point such as $1 \in [0, 1]$.
    
    Let $\{\psi_i\}_{i \in I}$ be a partition of unity subordinate to $\{U_i\}_{i \in I}$ \cite[Thm.~2.23]{Lee2012}.
    The $(n-1)$-form $\psi_i \omega$ has compact support in $U_i$.
    If Stokes' theorem holds for $\R^n$ and $\mathbb{H}^n$, then it holds for all the charts in our atlas.
    To see this, note that the integrals on both sides are invariant under an orientation-preserving diffeomorphism, by the change-of-variables formula, and that the diffeomorphism commutes with $\d$ by \cref{prop:naturality_of_exterior_derivative}.
    Moreover, it carries outward-pointing vectors to outward-pointing vectors, so it carries the boundary orientation of $\partial U_i$ to that of $\partial \mathbb{H}^n$.
    Also, since $\partial M \cap U_i = \partial U_i$ and $\sum_i \psi_i = 1$, we have
    \begin{align}
        \begin{split}
            \int_{\partial M} \omega &= \int_{\partial M} \sum_i \psi_i \omega \\
        &= \sum_i \int_{\partial M} \psi_i \omega \\
        &= \sum_i \int_{\partial U_i} \psi_i \omega \\
        &= \sum_i \int_{U_i} \d(\psi_i \omega) \\
        &= \sum_i \int_M \d(\psi_i \omega) \\
        &= \int_M \d \left( \sum_i \psi_i \omega \right) \\
        &= \int_M \d{\omega}.
        \end{split}
    \end{align}
    Hence, it is sufficient to show that Stokes' theorem holds for $\R^n$ and for $\mathbb{H}^n$.
    First, consider $\mathbb{H}^n$.
    Since $\omega$ has compact support, there exists a real number $R > 0$ such that $\supp \omega$ is contained in the interior of the rectangle $A = [-R, R]^{\times (n-1)} \times [0, R]$, relative to $\mathbb{H}^n$.
    The form $\omega$ may be written as
    \begin{equation}
        \omega = \sum_{i=1}^n \omega_i \d{x}^1 \wedge \cdots \wedge \widehat{\d{x}^i} \wedge \cdots \wedge \d{x}^n,
    \end{equation}
    where a hat above a symbol denotes its omission.
    Thus,
    \begin{align}
        \begin{split}
            \d{\omega} &= \sum_{i=1}^n \d{\omega}_i \wedge \d{x}^1 \wedge \cdots \wedge \widehat{\d{x}^i} \wedge \cdots \wedge \d{x}^n \\
        &= \sum_{i, j = 1}^n \frac{\partial \omega_i}{\partial x^j} \d{x}^j \wedge \d{x}^1 \wedge \cdots \wedge \widehat{\d{x}^i} \wedge \cdots \wedge \d{x}^n \\
        &= \sum_{i=1}^n (-1)^{i-1} \frac{\partial \omega_i}{\partial x^i} \d{x}^1 \wedge \cdots \wedge \d{x}^n.
        \end{split}
    \end{align}
    We use this to get
    \begin{align}
        \begin{split}
            \int_{\mathbb{H}^n} \d{\omega} &= \sum_{i=1}^n (-1)^{i-1} \int_A \frac{\partial \omega_i}{\partial x^i} \d{x}^1 \wedge \cdots \wedge \d{x}^n \\
            &= \sum_{i=1}^n (-1)^{i-1} \int_0^R \int_{-R}^R \cdots \int_{-R}^R \frac{\partial \omega_i}{\partial x^i} (x) \d{x}^1 \cdots \d{x}^n \\
        \end{split}
    \end{align}
    For the terms from $i=1$ to $i=n-1$, we change the order of integration and use the fundamental theorem of calculus to obtain
    \begin{align}
        \begin{split}
            \sum_{i=1}^{n-1} & (-1)^{i-1} \int_0^R \int_{-R}^R \cdots \int_{-R}^R \frac{\partial \omega_i}{\partial x^i} (x) \d{x}^1 \cdots \d{x}^n \\
            &= \sum_{i=1}^{n-1} (-1)^{i-1} \int_0^R \int_{-R}^R \cdots \int_{-R}^R \frac{\partial \omega_i}{\partial x^i} (x) \d{x}^i \d{x}^1 \cdots \widehat{\d{x}^i} \cdots \d{x}^n \\
            &= \sum_{i=1}^{n-1} (-1)^{i-1} \int_0^R \int_{-R}^R \cdots \int_{-R}^R \left[\omega_i(x) \right]_{x^i = -R}^{x^i = R} \d{x}^1 \cdots \widehat{\d{x}^i} \cdots \d{x}^n \\
            &= 0,
        \end{split}
    \end{align}
    since $\omega_i = 0$ when $x^i = \pm R$, as $\supp \omega$ lies in the interior of $A$.
    We are left with
    \begin{align}\label{eq:Stokes'_theorem_half-space_LHS}
        \begin{split}
            \int_{\mathbb{H}^n} \d{\omega} &= (-1)^{n-1} \int_0^R \int_{-R}^R \cdots \int_{-R}^R \frac{\partial \omega_n}{\partial x^n} (x) \d{x}^1 \cdots \d{x}^n \\
            &= (-1)^{n-1} \int_{-R}^R \cdots \int_{-R}^R \int_0^R \frac{\partial \omega_n}{\partial x^n} (x) \d{x}^n \d{x}^1 \cdots \d{x}^{n-1} \\
            &= (-1)^{n-1} \int_{-R}^R \cdots \int_{-R}^R \left[\omega_n(x) \right]_{x^n = 0}^{x^n = R} \d{x}^1 \cdots \d{x}^{n-1} \\
            &= (-1)^n \int_{-R}^R \cdots \int_{-R}^R \omega_n (x^1, \ldots, x^{n-1}, 0) \d{x}^1 \cdots \d{x}^{n-1},
        \end{split}
    \end{align}
    since $\omega_n = 0$ when $x^n = R$.
    Then, we compute the RHS of \eqref{eq:stokes'_theorem}.
    \begin{align}
        \begin{split}
            \int_{\partial \mathbb{H}^n} \omega &= \sum_i \int_{A \cap \mathbb{H}^n} \omega_i (x^1, \ldots, x^{n-1}, 0) \d{x}^1 \wedge \cdots \wedge \widehat{\d{x}^i} \wedge \cdots \wedge \d{x}^n \\
            &= \int_{A \cap \partial \mathbb{H}^n} \omega_n(x^1, \ldots, x^{n-1}, 0) \d{x}^1 \wedge \cdots \wedge \d{x}^{n-1},
        \end{split}
    \end{align}
    since $x^n$ vanishes on $\partial \mathbb{H}^n$ and the pullback of $\d{x}^n$ to the boundary is thus zero.
    From the fact that the coordinates $(x^1, \ldots, x^{n-1})$ are positively oriented for $\partial \mathbb{H}^n$ when $n$ is even and negatively oriented when $n$ is odd, we see that this is equal to \eqref{eq:Stokes'_theorem_half-space_LHS}.
    For $n = 1$, the iterated integrals over $x^1, \ldots, x^{n-1}$ are absent, and both sides equal $-\omega_1(0)$: the right-hand side is the signed evaluation at the boundary point $0$, which is negatively oriented because the outward vector $-\partial_1$ is negatively oriented.
    If instead the chart onto $\mathbb{H}^1$ reverses orientation, then the left-hand side changes sign, so it equals $\omega_1(0)$.
    The outward vector of $M$ at the boundary point maps to $-\partial_1$, which is negatively oriented in $\mathbb{H}^1$, so the outward vector is positively oriented in $M$.
    Hence the boundary point is positively oriented, and the right-hand side also equals $\omega_1(0)$.

    Finally, we need to consider $\R^n$.
    The support of $\omega$ is contained within the interior of the cube $A = [-R, R]^{\times n}$.
    The computations are the same as those for $\mathbb{H}^n$, except the $i = n$ term also vanishes, so the LHS of \eqref{eq:stokes'_theorem} is zero.
    The boundary of $\R^n$ is empty so the RHS is zero as well.
    If $\partial M = \varnothing$, every chart in the atlas is of this kind, and the argument above shows that both sides vanish.
    This completes our proof.
\end{proof}
\begin{corollary}\label{cor:integral_of_exact_form}
    Let $M$ be a compact oriented smooth manifold without boundary.
    Then, the integral of any exact form over $M$ is zero.
\end{corollary}
\begin{proof}
    Let $\omega = \d{\eta}$ with $\eta \in \Omega^{n-1}(M)$, which has compact support because $M$ is compact.
    By \cref{thm:stokes'_theorem}, $\int_M \d{\eta} = \int_{\partial M} \eta = 0$, since $\partial M = \varnothing$.
\end{proof}
\begin{corollary}\label{integral_of_closed_form_over_boundary}
    Let $M$ be a compact oriented smooth manifold with boundary.
    The integral of any closed form over $\partial M$ is zero.
\end{corollary}
\begin{proof}
    Let $\omega \in \Omega^{n-1}(M)$ be closed, which is the only degree for which $\int_{\partial M} \omega$ is defined.
    Since $M$ is compact, \cref{thm:stokes'_theorem} gives $\int_{\partial M} \omega = \int_M \d{\omega} = 0$.
\end{proof}
We can use Stokes' theorem to obtain some of the more elementary integral theorems of calculus.
Note that their hypotheses are not in their weakest forms due to the assumptions we chose for Stokes' theorem; it is by choice, and we could have proved the theorem with a weaker set of conditions.
We provide three integral theorems below that are immediately apparent.
The curl and divergence theorems, on the other hand, come more naturally after the discussion of Riemannian manifolds.
\begin{theorem}[Fundamental Theorem of Calculus]\label{thm:fundamental_theorem_of_calculus}
    Let $[a, b]$ be a closed interval of $\R$, and let $f \colon \R \to \R$ be a function defined by $f(x) = F'(x)$ for some smooth function $F \colon \R \to \R$.
    Then,
    \begin{equation}
        \int_a^b f(x) \d{x} = F(b) - F(a).
    \end{equation}
\end{theorem}
\begin{proof}
    Give $[a, b]$ the standard orientation and apply \cref{thm:stokes'_theorem} to the zero-form $F$, which has compact support.
    Here $\d{F} = f \d{x}$, and the boundary point $b$ is positively oriented and $a$ is negatively oriented, since the outward vectors are $\partial_x$ and $-\partial_x$ respectively.
    Hence $\int_a^b f(x) \d{x} = F(b) - F(a)$.
\end{proof}
\begin{theorem}[Gradient Theorem]\label{thm:gradient_theorem}
    Let $C$ be a smooth curve in $\R^3$ parametrized by $\mathbf{r}(t) = (x(t), y(t), z(t))$ with $a \leq t \leq b$, and let $\mathbf{f} \colon \R^3 \to \R^3$ be a function defined by $\mathbf{f} = \grad \varphi$ for some smooth function $\varphi \colon \R^3 \to \R$.
    Then,
    \begin{equation}
        \int_C \mathbf{f} \cdot d\mathbf{r} = \varphi(\mathbf{r}(b)) - \varphi(\mathbf{r}(a)).
    \end{equation}
\end{theorem}
\begin{proof}
    Since $\mathbf{f} = \grad \varphi$, the one-form $\mathbf{f} \cdot d\mathbf{r}$ is $\d{\varphi}$, so $\int_C \mathbf{f} \cdot d\mathbf{r} = \int_{[a, b]} \mathbf{r}^*\d{\varphi} = \int_{[a, b]} \d(\varphi \circ \mathbf{r})$ by \cref{prop:naturality_of_exterior_derivative}.
    Applying \cref{thm:fundamental_theorem_of_calculus} to $\varphi \circ \mathbf{r}$ gives $\varphi(\mathbf{r}(b)) - \varphi(\mathbf{r}(a))$.
\end{proof}
\begin{theorem}[Green's Theorem]\label{thm:green's_theorem}
    Let $D$ be a compact regular domain in $\R^2$, and let $P$ and $Q$ be a smooth real-valued functions on $D$.
    Then,
    \begin{equation}
        \int_D \left(\frac{\partial Q}{\partial x} - \frac{\partial P}{\partial y}\right) \d{x}\d{y} = \int_{\partial D} P \d{x} + Q \d{y}.
    \end{equation}
\end{theorem}
\begin{proof}
    Let $\omega = P \d{x} + Q \d{y}$.
    Then $\d{\omega} = \left(\frac{\partial Q}{\partial x} - \frac{\partial P}{\partial y}\right) \d{x} \wedge \d{y}$, and \cref{thm:stokes'_theorem} applied to $D$ with the standard orientation of $\R^2$ gives the claim.
\end{proof}

\section{De Rham Theory}\label{sec:de_rham_theory}

\subsection{De Rham Cohomology}\label{subsec:de_rham_theory:de_rham_cohomology}

\subsubsection{Motivation}\label{subsubsec:de_rham_theory:de_rham_cohomology:motivation}

De Rham cohomology measures how far closed forms are from being exact \cite{tao2008differential}.
On $\R$, a function with zero derivative is constant, by the fundamental theorem of calculus: $f(b) - f(a) = \int_a^b f'(x) \d{x} = 0$.
This fails on $\R \setminus \{0\}$, where the function equal to $+1$ on positive numbers and $-1$ on negative numbers has zero derivative but is not constant.
On $R_n = \R \setminus \{p_1, \ldots, p_n\}$, the locally constant functions form a space $A \cong \R^{n+1}$, with one constant for each of the $n+1$ pieces, and the globally constant functions form a subspace $B \cong \R$.
The quotient $A/B \cong \R^n$ counts the holes, and it is the zeroth reduced de Rham cohomology group $\widetilde{H}^0_{\mathrm{dR}}(R_n)$.

The same phenomenon occurs for forms of higher degree.
The one-form $\d{\theta}$ of \eqref{eq:angle-form} on $\R^2 \setminus \{0\}$ is closed but not exact, whereas on the slit plane $\R^2 \setminus \{(x, 0) : x \leq 0\}$ it is the exterior derivative of the argument function $\theta$ with values in $(-\pi, \pi)$.
Its integral over the unit circle $C$, parametrized by $t \mapsto (\cos t, \sin t)$, is
\begin{equation*}
    \int_C \d{\theta} = \int_0^{2\pi} (\sin^2 t + \cos^2 t) \d{t} = 2\pi \neq 0.
\end{equation*}
Hence Stokes' theorem cannot be applied to write this integral as the integral of a form over the empty boundary of the circle.
De Rham cohomology makes this failure precise.

\subsubsection{De Rham Complex}\label{subsubsec:de_rham_theory:de_rham_cohomology:de_rham_complex}

The space of differential $k$-forms $\Omega^k(M)$ for successive integer values of $k$ are connected via the exterior derivative.
We can thus construct the following sequence of vector spaces.
Here, note that 0 refers to the zero vector space $\{0\}$.
\[\begin{tikzcd}
	0 & {\Omega^0(M)} & {\Omega^1(M)} & {\Omega^2(M)} & {\Omega^3(M)} & \cdots
	\arrow[from=1-1, to=1-2]
	\arrow["{\d^0}", from=1-2, to=1-3]
	\arrow["{\d^1}", from=1-3, to=1-4]
	\arrow["{\d^2}", from=1-4, to=1-5]
	\arrow["{\d^3}", from=1-5, to=1-6]
\end{tikzcd}\]
Whilst all $\d^k \colon \Omega^k(M) \to \Omega^{k+1}(M)$ are specific instances of the exterior derivative $\d$, we have labelled them separately to show that the above sequence is an example of a more general concept.
The concept is called a cochain complex, which is the dual version of the chain complex.
What makes it a complex is the fact that $\d^2 = 0$, meaning that the image of $\d^k$, the space of exact $(k+1)$-forms, is contained in the kernel of $\d^{k+1}$, the space of closed $(k+1)$-forms, i.e. exact forms are closed.
\begin{definition}[(Co)Chain Complex]\label{def:(co)chain_complex}
    A chain complex $(A_\bullet, \partial_\bullet)$ is a $\Z$-indexed sequence $\ldots, A_0, A_1, A_2, \ldots$ of modules over a commutative ring $R$ connected by homomorphisms $\partial_k \colon A_k \to A_{k-1}$, such that $\partial_k \circ \partial_{k+1} = 0$.
    \[\begin{tikzcd}
	\cdots & {A_0} & {A_1} & {A_2} & {A_3} & \cdots
	\arrow["{\partial_0}"', from=1-2, to=1-1]
	\arrow["{\partial_1}"', from=1-3, to=1-2]
	\arrow["{\partial_2}"', from=1-4, to=1-3]
	\arrow["{\partial_3}"', from=1-5, to=1-4]
	\arrow["{\partial_4}"', from=1-6, to=1-5]
    \end{tikzcd}\]
    A cochain complex $(A^\bullet, \partial^\bullet)$ is a $\Z$-indexed sequence $\ldots, A^0, A^1, A^2, \ldots$ of modules over a commutative ring $R$ connected by homomorphisms $\partial^k \colon A^k \to A^{k+1}$, such that $\partial^{k+1} \circ \partial^k = 0$.
    \[\begin{tikzcd}
	\cdots & {A^0} & {A^1} & {A^2} & {A^3} & \cdots
	\arrow["{\partial^{-1}}", from=1-1, to=1-2]
	\arrow["{\partial^0}", from=1-2, to=1-3]
	\arrow["{\partial^1}", from=1-3, to=1-4]
	\arrow["{\partial^2}", from=1-4, to=1-5]
	\arrow["{\partial^3}", from=1-5, to=1-6]
    \end{tikzcd}\]
    The module homomorphisms are called boundary operators or differentials.
    The individual elements of a (co)chain complex are called (co)chains.
    The elements in the kernels of the boundary operators are called (co)cycles, and the elements of the images are called (co)boundaries.
    The suffix -$k$ attached to \enquote*{(co)chain}, \enquote*{(co)cycle}, or \enquote*{(co)boundary} indicates their associated element $A_k$.
    The space of $k$-(co)cycles may be denoted by $Z_k(A_\bullet)$ ($Z^k(A^\bullet)$), and the space of $k$-(co)boundaries is denoted by $B_k(A_\bullet)$ ($B^k(A^\bullet)$).

    The object $A_\bullet$ may be considered to be the collection of $R$-modules $A_i$ or the (co)chain complex itself; $\partial_\bullet$ may be considered to be the collection of boundary operators or a global boundary operator.
    The subscript indexing boundary operators are often omitted, i.e. $\partial \circ \partial = 0$ is written instead of $\partial_k \circ \partial_{k+1} = 0$.
\end{definition}
\begin{remark}
    We could have instead defined (co)chain complexes with Abelian groups, as is reasonably standard.
    However, defining them over $R$-modules is more general, simplifying to Abelian groups, vector spaces, or other desired structures.
    Note that at least the structure of an Abelian group is usually required as a minimum for associated theory to work well.
    Also, note that we use the notation for chain complexes -- primarily for basic concepts -- to refer to both their chain and cochain versions for convenience.
\end{remark}
\begin{definition}[Bounded (Co)Chain Complex]\label{def:bounded_(co)chain_complex}
    A (co)chain complex is bounded if all but a finite number of the (co)chains $A_i$ are zero.
    A (co)chain complex is bounded above if all $A_i$ are zero for all $i > n$, where $n \in \N$ is fixed; it is bounded below if $A_i = 0$ for all $i < n$.
    The $R$-module $A_n$ in the prior two cases may be called the end and start of the complex respectively.
    It may be called an endpoint in either case.
\end{definition}
\begin{remark}
    We often omit the trailing zeroes of a bounded (co)chain complex, writing just one zero at the start and end.
\end{remark}
\begin{definition}[De Rham Complex]\label{def:de_rham_complex}
    Let $M$ be a smooth manifold.
    The de Rham complex of $M$ is the cochain complex $(\Omega^\bullet(M), \d)$ consisting of the differential forms on $M$ and the exterior derivative.
    In this context, the space $\Omega^\bullet(M)$ may be called the de Rham algebra.
\end{definition}
\begin{remark}
    In the de Rham complex, cochains are spaces of differential forms, cocycles are closed forms, and coboundaries are exact forms.
\end{remark}
Note that the cochain complex structure on the space of differential forms, as well as its grading, makes the space a differential graded $\R$-algebra.
\begin{definition}[Differential Graded Algebra]\label{def:differential_graded_algebra}
    Let $A_\bullet$ be a $\Z$-graded algebra with a product $\cdot \colon A_\bullet \times A_\bullet \to A_\bullet$ and a graded derivation $D \colon A_\bullet \to A_\bullet$ of degree $\pm 1$.
    The space $(A_\bullet, D, \cdot)$ is a differential graded algebra if $(A_\bullet, D)$ is a chain complex for $\abs{D} = -1$, or a cochain complex for $\abs{D} = +1$.
\end{definition}
With the de Rham complex, we can consider the difference between the space of closed forms and exact forms in terms of the complex.
If for some positive integer $k$, the two spaces are equal, then $\d^{k-1}$ maps onto the space of closed $k$-forms.
Thus, a difference between the spaces is some degree of failure of $\d^{k-1}$ to map onto the closed $k$-forms.
We formalize this with exact sequences and measure it with cohomology.
\begin{definition}[Exact Sequence]\label{def:exact_sequence}
    A finite sequence of groups $G_i$ with group homomorphisms $f_i \colon G_i \to G_{i+1}$ is exact at $G_i$ if $\ker f_i = \im f_{i-1}$.
    A sequence is exact if it is exact at each $G_i$ except the initial and final groups of the sequence.
    An infinite sequence of groups is exact if it is exact at each $G_i$ in the sequence.
    A short exact sequence is a bounded exact sequence of the form
    \[\begin{tikzcd}
	0 & A & B & C & 0.
	\arrow[from=1-1, to=1-2]
	\arrow["f", from=1-2, to=1-3]
	\arrow["g", from=1-3, to=1-4]
	\arrow[from=1-4, to=1-5]
    \end{tikzcd}\]
    An exact sequence that is not short is called a long exact sequence.
\end{definition}
\begin{remark}
    If $A = 0$, the sequence
    \[\begin{tikzcd}
	0 & B & C
	\arrow["f", from=1-1, to=1-2]
	\arrow["g", from=1-2, to=1-3]
    \end{tikzcd}\]
    is exact if and only if $\ker g = \im f = 0$, so $g$ is injective.
    Similarly, if $C = 0$, the sequence
    \[\begin{tikzcd}
	A & B & 0
	\arrow["f", from=1-1, to=1-2]
	\arrow["g", from=1-2, to=1-3]
    \end{tikzcd}\]
    is exact if and only if $\im f = \ker g = B$, so $f$ is surjective.
\end{remark}
If a (co)chain complex is exact, then $\ker \partial_k = \im \partial_{k+1}$ ($\ker \partial^k = \im \partial^{k-1}$).
In general, due to $\partial_k \circ \partial_{k+1} = 0$ ($\partial^{k+1} \circ \partial^k = 0$), we at least have $\im \partial_{k+1} \subseteq \ker \partial_k$ ($\im \partial^{k-1} \subseteq \ker \partial^k$).
This allows us to construct quotient spaces out of the images and kernels of the boundary maps, giving us the (co)homology groups; they are quotient $R$-modules.
\begin{definition}[(Co)Homology]\label{def:(co)homology}
    Let $(A_\bullet, \partial_\bullet)$ be a chain complex.
    The $k$-th homology group $H_k$ is defined to be
    \begin{equation}
        H_k = \frac{\ker \partial_k}{\im \partial_{k+1}}.
    \end{equation}
    Let $(A^\bullet, \partial^\bullet)$ be a cochain complex.
    The $k$-th cohomology group $H^k$ is defined to be
    \begin{equation}
        H^k = \frac{\ker \partial^k}{\im \partial^{k-1}}.
    \end{equation}
    The (co)homology groups may also be denoted by $H_k(A_\bullet)$ and $H^k(A^\bullet)$, respectively.
    Elements in $A_\bullet$ that are equivalent by the induced equivalence relation are called homologous, and equivalent elements of $A^\bullet$ are called cohomologous.
    The induced equivalence classes are called (co)homology classes.
\end{definition}
\begin{remark}
    In the direction of the boundary operators, the (co)homology at the first term of a bounded (co)chain complex is the kernel of the boundary operator out of it, and at the last term it is the cokernel of the boundary operator into it.
    Also, there is the concept of reduced (co)homology \cite{wikipedia_reduced_homology}, briefly mentioned in \Cref{subsubsec:de_rham_theory:de_rham_cohomology:motivation}; it differs only in its definition of the zeroth (co)homology groups.
    We will not discuss it any further.
\end{remark}
\begin{definition}[De Rham Cohomology]\label{def:de_rham_cohomology}
    Let $M$ be a smooth manifold.
    The $k$-th de Rham cohomology group of $M$, denoted by $H^k_{\mathrm{dR}}(M)$ or $H^k(M)$ when unambiguous, is defined to be the $k$-th cohomology group of $(\Omega^\bullet(M), \d)$.
    \begin{equation}
        H^k_{\mathrm{dR}}(M) = \frac{\ker \d^k}{\im \d^{k-1}}
    \end{equation}
    We may denote $\ker \d^k$ by $Z^k(M)$ and $\im \d^{k-1}$ by $B^k(M)$, so $H^k_{\mathrm{dR}}(M) = Z^k(M)/B^k(M)$.
    The de Rham cohomology of $M$, denoted by $H^\bullet_{\mathrm{dR}}(M)$ or $H^\bullet(M)$, is the anticommutative graded $\R$-algebra defined by
    \begin{equation}
        H^\bullet_{\mathrm{dR}}(M) = \bigoplus_{k=0}^\infty H^k_{\mathrm{dR}}(M).
    \end{equation}
    The algebra $H^\bullet_{\mathrm{dR}}(M)$ may also be called the de Rham cohomology ring or the de Rham cohomology algebra.
\end{definition}
With the goal of studying differential forms, we have invented much more general mathematics.
(Co)Homology generally measures the failure of the boundary operators of a (co)chain complex to map onto the (co)cycles.
Equivalently, it measures the failure of the (co)chain complex itself to be exact; exact sequences are thus sometimes defined as those with zero homology.
Specifically, the de Rham cohomology measures the exactness of the de Rham complex or the failure of the exterior derivative to map onto the closed forms, which translates to the difference between closed and exact differential forms.

As with any (co)homology, we have an induced equivalence relation between elements of the kernel of the boundary operators, which happen to be closed forms.
Two closed forms $\omega$ and $\eta$ are cohomologous if their difference is an exact form.
\begin{equation}
    \omega \sim \eta \quad \iff \quad \omega - \eta \in B^k(M)
\end{equation}
Thus, $\omega$ and $\eta$ determine the same cohomology class, $[\omega] = [\eta]$, if
\begin{equation}
    \eta = \omega + \d{\tau}
\end{equation}
for some differential form $\tau$ of one fewer degree.
We can also consider algebra on cohomology classes via the wedge product, defined by
\begin{equation}
    [\omega] \wedge [\eta] = [\omega \wedge \eta].
\end{equation}
This is well-defined.
If $\omega \in Z^k(M)$ and $\eta$ are closed, then $\d(\omega \wedge \eta) = \d{\omega} \wedge \eta + (-1)^k \omega \wedge \d{\eta} = 0$.
Moreover, for differential forms $\tau$ and $\sigma$, we have $(\omega + \d{\tau}) \wedge \eta = \omega \wedge \eta + \d(\tau \wedge \eta)$ and $\omega \wedge (\eta + \d{\sigma}) = \omega \wedge \eta + (-1)^k \d(\omega \wedge \sigma)$, since $\d{\omega} = \d{\eta} = 0$, so $[\omega \wedge \eta]$ does not depend on the choice of representatives.
Elements of the de Rham cohomology algebra $H^\bullet(M)$ are finite sums of cohomology classes in $H^k(M)$ for different values of $k$, and they may be combined via addition and (wedge) multiplication.
The algebra being an anticommutative graded algebra follows from the properties of the wedge product.
Note that it is vacuously a differential graded algebra as well, with $\d[\omega] = 0$.

\begin{example}
    Consider $\R^n$.
    For $n = 0$, where $\R^n$ is the single point $\{0\}$, we have
    \begin{equation}
        H^k(\R^0) \cong \begin{cases}
            \R & \text{if } k = 0, \\
            0 & \text{if } k > 0.
        \end{cases}
    \end{equation}
    The closed zero-forms on $\{0\}$ are constant functions characterized by some real number, so $H^0(\R^0) \cong \R$.
    For $k > 0$, alternating $k$-forms on the zero vector space vanish, so $\Omega^k(\R^0) = 0$ and $H^k(\R^0) = 0$.
    Note that we may often make statements such as $H^1(\R^0) = 0$, using the symbol for equality and the symbol for isomorphism interchangeably.

    For $n = 1$, we have $\ker \d^0 \cong \R$, since the closed zero-forms on $\R$ are constant functions.
    Consider a one-form $\omega = f(x) \d{x}$ on $\R$; we wish to show that this is exact.
    Define $g \colon \R \to \R$ by
    \begin{equation}
        g(x) = \int_0^x f(t) \d{t},
    \end{equation}
    so
    \begin{equation}
        \d{g} = f(x) \d{x}.
    \end{equation}
    Hence, every one-form on $\R$ is exact, meaning that $H^1(\R) = 0$.
    All differential forms on $\R$ of degree greater than one will be zero, because the only basis covector available is $\d{x}$, and $\d{x} \wedge \d{x} = 0$.
    This gives us $H^k(\R) = 0$ for $k > 1$.
    Therefore, we have
    \begin{equation}
        H^k(\R) \cong \begin{cases}
            \R & \text{if } k = 0, \\
            0 & \text{if } k > 0.
        \end{cases}
    \end{equation}
    From our prior knowledge and intuition, and the fact that $\R^n$ contains no holes, we expect its de Rham cohomology groups to be zero for $k > 0$.
    The zeroth de Rham cohomology group should be $\R$ by the reasoning of constant functions.
    We will prove the case of $\R^n$ later using the Poincar\'{e} lemma.
\end{example}
The computations above are special cases of the following results, which are often useful when computing the de Rham cohomology of other manifolds.
\begin{proposition}\label{prop:de_rham_cohomology_connected_components}
    Let $M$ be a smooth manifold with $r$ connected components.
    Then, its zeroth de Rham cohomology is $H^0(M) \cong \R^r$, and each element of the ordered $r$-tuple representing an element of $H^0(M)$ represents a constant function on a connected component of $M$.
\end{proposition}
\begin{proof}
    We know that $H^0(M) = Z^0(M)$, the space of closed zero-forms.
    Let $f \in Z^0(M)$, so $f$ is a smooth function such that $d{f} = 0$.
    On any chart $(U, \varphi)$ with $\varphi = (x^1, \ldots, x^n)$, we have
    \begin{equation*}
        \d{f} = \frac{\partial f}{\partial x^i} \d{x}^i = 0,
    \end{equation*}
    so the partial derivatives $\partial f / \partial x^i$ must all be zero on $U$.
    Thus, $f$ is locally constant on $U$, so all closed zero-forms on $M$ are locally constant functions on $M$.
    It is constant on each connected component of $M$, since the sets on which it takes a given value are open and disjoint, so if $M$ has $r$ connected components, then a locally constant function on $M$ is specified by an ordered $r$-tuple of real numbers.
    Therefore, $H^0(M) = Z^0(M) \cong \R^r$.
\end{proof}
\begin{proposition}\label{prop:de_rham_cohomology_max_dimension}
    Let $M$ be a smooth $n$-dimensional manifold.
    Then, its $k$-th de Rham cohomology is zero for $k > n$.
\end{proposition}
\begin{proof}
    At a point $p \in M$, the tangent space $T_p M$ has dimension $n$.
    A differential $k$-form $\omega \in \Omega^k(M)$ assigns to $p$ a $k$-form $\omega_p \in \Alt^k(T_p M, \R)$.
    If $k > n$, then $\Alt^k(T_p M, \R) = 0$, so the only $k$-form on $M$ is the zero form.
    Therefore, $H^k(M) = \{0\} / \{0\} \cong \{0\}$.
\end{proof}
We can already carry out some further cohomology computations.
However, thus far, we do not have the tools and results to do much efficiently.
\begin{example}
    Consider the unit circle $\mathbb{S}^1$ embedded in $\R^2$.
    Since $\mathbb{S}^1$ is connected, we have $H^0(\mathbb{S}^1) \cong \R$ by \Cref{prop:de_rham_cohomology_connected_components}.
    As it is one dimensional, we also have $H^k(\mathbb{S}^1) = 0$ for $k > 1$.
    That just leaves us with $H^1(\mathbb{S}^1)$.

    The restriction of the map $h \colon \R \to \mathbb{S}^1$, $h(t) = (\cos t, \sin t)$ to the interval $[0, 2\pi]$ via the inclusion map $i \colon [0, 2\pi] \hookrightarrow \R$ gives us a parametrization $\gamma = h \circ i \colon [0, 2\pi] \to \mathbb{S}^1$ of the circle.
    We orient $\mathbb{S}^1$ counterclockwise, so that $\gamma$ is orientation preserving.
    The one-form $\omega = -y \d{x} + x \d{y}$ is nowhere zero on $\mathbb{S}^1$, and $\gamma^*\omega = i^*h^*\omega = i^*\d{t} = \d{t}$.
    Thus,
    \begin{equation*}
        \int_{\mathbb{S}^1} \omega = \int_{\gamma([0, 2\pi])} \omega = \int_{[0, 2\pi]} \gamma^*\omega = \int_0^{2\pi} \d{t} = 2\pi.
    \end{equation*}
    Since $\mathbb{S}^1$ is one dimensional, all one-forms on $\mathbb{S}^1$ are closed, so $\Omega^1(\mathbb{S}^1) = Z^1(\mathbb{S}^1)$.
    Define the linear map $\varphi \colon Z^1(\mathbb{S}^1) \to \R$ by
    \begin{equation}
        \varphi(\alpha) = \int_{\mathbb{S}^1} \alpha.
    \end{equation}
    We will prove the desired result using the first isomorphism theorem.
    First, note that $\varphi$ is surjective, since $\varphi(\omega) = 2\pi \neq 0$.
    We want to show that $\ker \varphi = B^1(\mathbb{S}^1)$.
    By \hyperref[thm:stokes'_theorem]{Stokes' theorem}, the exact one-forms on $\mathbb{S}^1$ are in $\ker \varphi$, so it is left to show that all one-forms in $\ker \varphi$ are exact.
    Let $\alpha = f\omega$ be a one-form on $\mathbb{S}^1$, where $f$ is a smooth function on $\mathbb{S}^1$, such that $\varphi(\alpha) = 0$.
    Let $\bar{f} = h^*f = f \circ h \in \Omega^0(\R)$, so $\bar{f}$ is periodic with period $2\pi$.
    Also,
    \begin{equation*}
        0 = \int_{\mathbb{S}^1} \alpha = \int_{\gamma([0, 2\pi])} \alpha = \int_{[0, 2\pi]} \gamma^*\alpha = \int_{[0, 2\pi]} (i^*h^*f)(t) \cdot \gamma^*\omega = \int_0^{2\pi} \bar{f}(t) \d{t}.
    \end{equation*}
    Define $\bar{g} \in \Omega^0(\R)$ by
    \begin{equation}
        \bar{g}(t) = \int_0^t \bar{f}(u) \d{u}.
    \end{equation}
    Since $\bar{f}$ has period $2\pi$ and $\int_0^{2\pi} \bar{f}(u) \d{u} = 0$, we get
    \begin{equation*}
        \bar{g}(t + 2\pi) = \int_0^{2\pi} \bar{f}(u) \d{u} + \int_{2\pi}^{t+2\pi} \bar{f}(u) \d{u} = \int_{2\pi}^{t+2\pi} \bar{f}(u) \d{u} = \int_0^t \bar{f}(u) \d{u} = \bar{g}(t).
    \end{equation*}
    Thus, $\bar{g}$ is periodic with period $2\pi$.
    Furthermore, we have
    \begin{equation}
        \d{\bar{g}} = \bar{g}'(t) \d{t} = \bar{f}(t) \d{t}.
    \end{equation}
    There exists a smooth function $g$ on $\mathbb{S}^1$ such that $\bar{g} = h^*g$; the pullback map $h^*$ identifies smooth functions on $\mathbb{S}^1$ with smooth periodic functions on $\R$ with period $2\pi$.
    Hence,
    \begin{equation}\label{eq:circle_de_rham_cohomology_d_bar_g_formula}
        \d{\bar{g}} = \d{h}^*g = h^*(\d{g}).
    \end{equation}
    Using the fact that $\bar{f} = h^*f$ and $\d{t} = h^*\omega$,
    \begin{equation}\label{eq:circle_de_rham_cohomology_bar_f_t_formula}
        \bar{f}(t) \d{t} = (h^*f)(h^*\omega) = h^*(f\omega) = h^*\alpha.
    \end{equation}
    Since $h^* \colon \Omega^1(\mathbb{S}^1) \to \Omega^1(\R)$ is injective, we get $\alpha = \d{g}$ from \eqref{eq:circle_de_rham_cohomology_d_bar_g_formula} and \eqref{eq:circle_de_rham_cohomology_bar_f_t_formula}.
    This means that the kernel of $\varphi$ consists of exact forms, and it is precisely the exact forms: $\ker \varphi = B^1(\mathbb{S}^1)$.
    We also know that $\varphi$ is a surjective linear map (homomorphism).
    Therefore, by the first isomorphism theorem,
    \begin{equation}
        H^1(\mathbb{S}^1) = \frac{Z^1(\mathbb{S}^1)}{B^1(\mathbb{S}^1)} = \frac{Z^1(\mathbb{S}^1)}{\ker \varphi} \cong \im \varphi = \R.
    \end{equation}
\end{example}
A very important property of de Rham cohomology is that it is a diffeomorphism invariant.
Given a smooth map $f \colon M \to N$, there is an induced pullback map $f^* \colon \Omega^\bullet(N) \to \Omega^\bullet(M)$. 
From this, there is an induced map between $H^k(M)$ and $H^k(N)$ that is an isomorphism of vector spaces when $f$ is a diffeomorphism.
\begin{lemma}\label{lem:pullback_closed_and_exact_forms}
    Let $M$ and $N$ be smooth manifolds, and let $f \colon M \to N$ be a smooth function.
    The pullback map $f^* \colon \Omega^\bullet(N) \to \Omega^\bullet(M)$ sends closed forms to closed forms and exact forms to exact forms.
\end{lemma}
\begin{proof}
    Let $\omega \in \Omega^k(N)$ be a closed differential $k$-form, so $\d{\omega} = 0$.
    By \cref{prop:naturality_of_exterior_derivative}, we have $\d(f^*\omega) = f^*\d{\omega} = f^*0 = 0$, so $f^*\omega$ is closed.
    Now, instead suppose that $\omega$ is an exact differential $k$-form, so $\omega = \d{\eta}$ for some differential $(k-1)$-form $\eta \in \Omega^{k-1}(N)$.
    Then, $f^*\omega = f^*\d{\eta} = \d(f^*\eta)$, so $f^*\omega$ is exact.
\end{proof}
\begin{proposition}[Induced De Rham Cohomology Map]\label{prop:induced_de_rham_cohomology_maps}
    Let $M$ and $N$ be smooth manifolds, and let $f \colon M \to N$ be a smooth function.
    The pullback map $f^* \colon \Omega^\bullet(N) \to \Omega^\bullet(M)$ induces a linear map $f^* \colon H^k(N) \to H^k(M)$, called the induced cohomology map, with the following properties.
    \begin{enumerate}
        \item If $P$ is another smooth manifold and $g \colon N \to P$ is a smooth map, then
        \begin{equation}
            (g \circ f)^* = f^* \circ g^* \colon H^k(P) \to H^k(M).
        \end{equation}
        \item The identity map of $H^k(M)$ is given by $\id_M^*$.
    \end{enumerate}
\end{proposition}
\begin{proof}
    The induced cohomology map $f^* \colon H^k(N) \to H^k(M)$ is defined by
    \begin{equation}
        f^*[\omega] = [f^*\omega].
    \end{equation}
    For the well-definedness of this map, first note \Cref{lem:pullback_closed_and_exact_forms}.
    Also, if $\eta = \omega + \d{\tau}$, i.e. $[\eta] = [\omega]$, then $[f^*\eta] = [f^*\omega + \d(f^*\tau)] = [f^*\omega]$.
    The two properties of the induced map characterizing its functoriality follow from the functoriality of the pullback map.
\end{proof}
\begin{remark}
    If there is ambiguity in the notation for the induced cohomology map, then we may use a different symbol such as $f^\#$ to distinguish it from the pullback map.
\end{remark}
\begin{corollary}\label{cor:de_rham_cohomology_diffeomorphism_invariance}
    Let $M$ and $N$ be smooth manifolds, and let $f \colon M \to N$ be a diffeomorphism.
    Then, the induced cohomology map $f^* \colon H^k(N) \to H^k(M)$ is a vector-space isomorphism.
\end{corollary}
\begin{proof}
    By \cref{prop:induced_de_rham_cohomology_maps}, $(f^{-1})^* \circ f^* = (f \circ f^{-1})^* = \id_N^*$ and $f^* \circ (f^{-1})^* = (f^{-1} \circ f)^* = \id_M^*$, and these are the identity maps on cohomology.
    Hence $f^*$ is invertible with inverse $(f^{-1})^*$.
\end{proof}
Since, for any smooth map $f \colon M \to N$, we have $f^*(\omega \wedge \eta) = f^*\omega \wedge f^*\eta$, the induced linear map $f^* \colon H^\bullet(N) \to H^\bullet(M)$ is an algebra homomorphism.
If $f$ is a diffeomorphism, then $f^*$ is an algebra isomorphism, making the de Rham cohomology algebra a diffeomorphism invariant as well.

\subsubsection{Mayer--Vietoris Sequence}\label{subsubsec:de_rham_theory:de_rham_cohomology:mayer--vietoris_sequence}

We wish to introduce the Mayer--Vietoris sequence, which is one of the most useful tools for the computation of the (co)homology of spaces, including de Rham cohomology.
Before we are able to construct it, we require further results of (co)chain complexes.
When we provide generic results, if the objects of our complexes are not mentioned, then they are assumed to be $R$-modules.
\begin{proposition}\label{prop:three-term_exact_sequence}
    Let
    \[\begin{tikzcd}
	A & B & C
	\arrow["f", from=1-1, to=1-2]
	\arrow["g", from=1-2, to=1-3]
    \end{tikzcd}\]
    be an exact sequence.
    Then,
    \begin{enumerate}
        \item the map $f$ is surjective if and only if $g$ is the zero map; and
        \item the map $g$ is injective if and only if $f$ is the zero map.
    \end{enumerate}
\end{proposition}
\begin{proof} \leavevmode
    \begin{enumerate}
        \item First, suppose that $f$ is surjective, so $\im f = B$.
              Since the sequence is exact, $\im f = \ker g$, so $\ker g = B$.
              Hence, $g$ is the zero map.
              The converse is similarly shown.
        \item Suppose that $g$ is injective, so $\ker g = 0$.
              Since the sequence is exact, $\im f = \ker g$, so $\im f = 0$.
              Hence, $f$ is the zero map.
              The converse is similarly shown. \qedhere
    \end{enumerate}
\end{proof}
\begin{proposition}\label{prop:four-term_exact_sequence} \leavevmode
    \begin{enumerate}
        \item The four-term sequence $0 \to A \xrightarrow{f} B \to 0$ is exact if and only if $f \colon A \to B$ is an isomorphism.
        \item If $A \xrightarrow{f} B \to C \to 0$ is an exact sequence, then there is an isomorphism $C \cong \coker f \coloneq B / \im f$.
    \end{enumerate}
\end{proposition}
\begin{proof} \leavevmode
    \begin{enumerate}
        \item Suppose that $0 \to A \xrightarrow{f} B \to 0$ is exact, so $\im f = \ker 0_B = B$ and $0 = \im 0^A = \ker f$.
              This shows that $f$ is surjective and injective, respectively, so $f$ is bijective.
              Hence, $f$ is an isomorphism.
              Apply this reasoning in reverse to prove the converse.
        \item Suppose that $A \xrightarrow{f} B \to C \to 0$ is an exact sequence.
              Denote the map from $B$ to $C$ by $g$.
              Exactness implies $\im f = \ker g$ and $\im g = \ker 0_C = C$.
              The latter fact means that $g$ is surjective.
              We are thus able to use the first isomorphism theorem to conclude that $\coker f \coloneq B / \im f = B / \ker g \cong C$. \qedhere
    \end{enumerate}
\end{proof}
We can also consider maps between (co)chain complexes, which are sequences of maps between the elements of each complex such that those maps commute with the boundary operators.
\begin{definition}[(Co)Chain Map]\label{def:(co)chain_map}
    Let $(A_\bullet, \partial_{A,\bullet})$ and $(B_\bullet, \partial_{B,\bullet})$ be chain complexes.
    A chain map $F$ between the two complexes is a sequence $F_\bullet$ of homomorphisms $F_k \colon A_k \to B_k$ for each $k$ that each commute with the boundary operators of the chain complexes, i.e. $\partial_{B,k} \circ F_k = F_{k-1} \circ \partial_{A,k}$, so that the following diagram commutes.
    \[\begin{tikzcd}[sep=large]
	\cdots & {A_{k-1}} & {A_k} & {A_{k+1}} & \cdots \\
	\cdots & {B_{k-1}} & {B_k} & {B_{k+1}} & \cdots
	\arrow["{\partial_{A,k-1}}"', from=1-2, to=1-1]
	\arrow["{F_{k-1}}"', from=1-2, to=2-2]
	\arrow["{\partial_{A,k}}"', from=1-3, to=1-2]
	\arrow["{F_k}"', from=1-3, to=2-3]
	\arrow["{\partial_{A,k+1}}"', from=1-4, to=1-3]
	\arrow["{F_{k+1}}"', from=1-4, to=2-4]
	\arrow["{\partial_{A,k+2}}"', from=1-5, to=1-4]
	\arrow["{\partial_{B,k-1}}"', from=2-2, to=2-1]
	\arrow["{\partial_{B,k}}"', from=2-3, to=2-2]
	\arrow["{\partial_{B,k+1}}"', from=2-4, to=2-3]
	\arrow["{\partial_{B,k+2}}"', from=2-5, to=2-4]
    \end{tikzcd}\]
    Let $(A^\bullet, \partial^{A,\bullet})$ and $(B^\bullet, \partial^{B,\bullet})$ be cochain complexes.
    A cochain map $F$ between the two complexes is a sequence $F^\bullet$ of homomorphisms $F^k \colon A^k \to B^k$ for each $k$ that each commute with the boundary operators of the cochain complexes, i.e. $\partial^{B,k} \circ F^k = F^{k+1} \circ \partial^{A,k}$, so that the following diagram commutes.
    \[\begin{tikzcd}[sep=large]
	\cdots & {A^{k-1}} & {A^k} & {A^{k+1}} & \cdots \\
	\cdots & {B^{k-1}} & {B^k} & {B^{k+1}} & \cdots
	\arrow["{{\partial^{A,k-2}}}", from=1-1, to=1-2]
	\arrow["{{\partial^{A,k-1}}}", from=1-2, to=1-3]
	\arrow["{{F^{k-1}}}", from=1-2, to=2-2]
	\arrow["{{\partial^{A,k}}}", from=1-3, to=1-4]
	\arrow["{{F^k}}", from=1-3, to=2-3]
	\arrow["{{\partial^{A,k+1}}}", from=1-4, to=1-5]
	\arrow["{{F^{k+1}}}", from=1-4, to=2-4]
	\arrow["{{\partial^{B,k-2}}}", from=2-1, to=2-2]
	\arrow["{{\partial^{B,k-1}}}", from=2-2, to=2-3]
	\arrow["{{\partial^{B,k}}}", from=2-3, to=2-4]
	\arrow["{{\partial^{B,k+1}}}", from=2-4, to=2-5]
\end{tikzcd}\]
\end{definition}
\begin{remark}
    We often omit the index $k$ when referencing the $k$-th homomorphism of a (co)chain map.
    This is also true for derived concepts such as maps induced from the (co)chain maps.
\end{remark}
\begin{proposition}[Induced (Co)Homology Map]\label{prop:induced_(co)chain_(co)homology_map}
    A chain map $F \colon A_\bullet \to B_\bullet$ between two chain complexes induces a homomorphism of homology groups $F_* \colon H_k(A_\bullet) \to H_k(B_\bullet)$ given by
    \begin{equation}
        F_*[a] = [F(a)].
    \end{equation}
    A cochain map $F \colon A^\bullet \to B^\bullet$ between two cochain complexes induces a homomorphism of cohomology groups $F^* \colon H^k(A^\bullet) \to H^k(B^\bullet)$ given by
    \begin{equation}
        F^*[a] = [F(a)].
    \end{equation}
\end{proposition}
\begin{proof}
    We show that $F$ sends (co)cycles to (co)cycles and (co)boundaries to (co)boundaries, so that the induced map is well-defined.
    If $\partial a = 0$, then $\partial F(a) = F(\partial a) = 0$.
    If $a = \partial b$, then $F(a) = F(\partial b) = \partial F(b)$.
    The same computation applies to cochains and cochain maps.
\end{proof}
\begin{example}
    The induced pullback map $f^* \colon H_{\mathrm{dR}}^k(N) \to H_{\mathrm{dR}}^k(M)$ between de Rham cohomology groups is an example of an induced cochain map.
    For some smooth vector field on $M$, the Lie derivative $\mathcal{L}_X \colon \Omega^\bullet(M) \to \Omega^\bullet(M)$ commutes with the exterior derivative, so it induces a linear map $\mathcal{L}^*_X \colon H_{\mathrm{dR}}^\bullet(M) \to H_{\mathrm{dR}}^\bullet(M)$ in cohomology.
    Note here that the induced map between cohomology algebras comes first from the induced map between cohomology groups.
    We have not yet discussed such algebra structure for (co)homology besides de Rham cohomology.
\end{example}
With (co)chain maps, we can now consider sequences of (co)chain complexes.
Such a sequence, e.g.
\[\begin{tikzcd}
	{A_\bullet} & {B_\bullet} & {C_\bullet},
	\arrow["F", from=1-1, to=1-2]
	\arrow["G", from=1-2, to=1-3]
\end{tikzcd}\]
is really a compact version of a much more complicated diagram.
We can consider properties of such diagrams and prove results about them.
\begin{definition}[Short Exact Sequence of (Co)Chain Complexes]\label{def:short_exact_sequence_of_(co)chain_complexes}
    A short exact sequence of chain complexes is a sequence of three chain complexes $A_\bullet$, $B_\bullet$, and $C_\bullet$ with chain maps $F$ and $G$,
    \[\begin{tikzcd}
	0 & {A_\bullet} & {B_\bullet} & {C_\bullet} & 0,
	\arrow[from=1-1, to=1-2]
	\arrow["F", from=1-2, to=1-3]
	\arrow["G", from=1-3, to=1-4]
	\arrow[from=1-4, to=1-5]
    \end{tikzcd}\]
    such that each sequence
    \[\begin{tikzcd}
	0 & {A_k} & {B_k} & {C_k} & 0
	\arrow[from=1-1, to=1-2]
	\arrow["{F_k}", from=1-2, to=1-3]
	\arrow["{G_k}", from=1-3, to=1-4]
	\arrow[from=1-4, to=1-5]
    \end{tikzcd}\]
    is exact.
    A short exact sequence of cochain complexes $A^\bullet$, $B^\bullet$, and $C^\bullet$ is defined analogously.
\end{definition}
\begin{remark}
    In such a sequence, $F$ is injective, $G$ is surjective, and $\im F = \ker G$.
    Also, $G \circ F = 0$.
\end{remark}
We will state and prove the zig-zag lemma, which happens to be an important result required for the proof of the (existence of the) Mayer--Vietoris sequence.
The process of our proof involves checking the exactness at different parts of a process; the technique used is commonly called diagram chasing.
\begin{lemma}[Zig-Zag Lemma]\label{lem:zig-zag_lemma}
    Given a short exact sequence of chain complexes,
    \[\begin{tikzcd}
	0 & {A_\bullet} & {B_\bullet} & {C_\bullet} & 0,
	\arrow[from=1-1, to=1-2]
	\arrow["F", from=1-2, to=1-3]
	\arrow["G", from=1-3, to=1-4]
	\arrow[from=1-4, to=1-5]
    \end{tikzcd}\]
    there exists a collection of homomorphisms $\delta_k \colon H_k(C_\bullet) \to H_{k-1}(A_\bullet)$, called connecting homomorphisms, such that the following sequence is exact.
    \[\begin{tikzcd}
    & \cdots \ar[r, "G_*"] & H_{k+1}(C_\bullet) \ar[out=0, in=180, looseness=2, overlay, lld, "\delta_{k+1}"{description}]\\
    H_k(A_\bullet) \ar[r, "F_*"] & H_k(B_\bullet) \ar[r, "G_*"] & H_k(C_\bullet) \ar[out=0, in=180, looseness=2, overlay, lld, "\delta_k"{description}]\\
    H_{k-1}(A_\bullet) \ar[r, "F_*"] & \cdots &
    \end{tikzcd}\]
    Given a short exact sequence of cochain complexes,
    \[\begin{tikzcd}
	0 & {A^\bullet} & {B^\bullet} & {C^\bullet} & 0,
	\arrow[from=1-1, to=1-2]
	\arrow["F", from=1-2, to=1-3]
	\arrow["G", from=1-3, to=1-4]
	\arrow[from=1-4, to=1-5]
    \end{tikzcd}\]
    there exists a collection of homomorphisms $\delta^k \colon H^k(C^\bullet) \to H^{k+1}(A^\bullet)$, called connecting homomorphisms, such that the following sequence is exact.
    \[\begin{tikzcd}
    & \cdots \ar[r, "G^*"] & H^{k-1}(C^\bullet) \ar[out=0, in=180, looseness=2, overlay, lld, "\delta^{k-1}"{description}]\\
    H^k(A^\bullet) \ar[r, "F^*"] & H^k(B^\bullet) \ar[r, "G^*"] & H^k(C^\bullet) \ar[out=0, in=180, looseness=2, overlay, lld, "\delta^k"{description}]\\
    H^{k+1}(A^\bullet) \ar[r, "F^*"] & \cdots &
    \end{tikzcd}\]
\end{lemma}
\begin{proof}
    We will prove the result for cochain complexes, which can be very slightly modified to work for chain complexes.
    The short exact sequence of cochain complexes is contained in the following commutative diagram.
    \[\begin{tikzcd}[sep=large]
	& \vdots & \vdots & \vdots & \\
	0 & {A^{k-1}} & {B^{k-1}} & {C^{k-1}} & 0 \\
	0 & {A^k} & {B^k} & {C^k} & 0 \\
	0 & {A^{k+1}} & {B^{k+1}} & {C^{k+1}} & 0 \\
	& \vdots & \vdots & \vdots
	\arrow[from=1-2, to=2-2]
	\arrow[from=1-3, to=2-3]
	\arrow[from=1-4, to=2-4]
	\arrow[from=2-1, to=2-2]
	\arrow["{F^{k-1}}", from=2-2, to=2-3]
	\arrow["{\partial^{A, k-1}}", from=2-2, to=3-2]
	\arrow["{G^{k-1}}", from=2-3, to=2-4]
	\arrow["{\partial^{B, k-1}}", from=2-3, to=3-3]
	\arrow[from=2-4, to=2-5]
	\arrow["{\partial^{C, k-1}}", from=2-4, to=3-4]
	\arrow[from=3-1, to=3-2]
	\arrow["{F^k}", from=3-2, to=3-3]
	\arrow["{\partial^{A, k}}", from=3-2, to=4-2]
	\arrow["{G^k}", from=3-3, to=3-4]
	\arrow["{\partial^{B, k}}", from=3-3, to=4-3]
	\arrow[from=3-4, to=3-5]
	\arrow["{\partial^{C, k}}", from=3-4, to=4-4]
	\arrow[from=4-1, to=4-2]
	\arrow["{F^{k+1}}", from=4-2, to=4-3]
	\arrow[from=4-2, to=5-2]
	\arrow["{G^{k+1}}", from=4-3, to=4-4]
	\arrow[from=4-3, to=5-3]
	\arrow[from=4-4, to=4-5]
	\arrow[from=4-4, to=5-4]
    \end{tikzcd}\]
    We need to show exactness at $H^k(A^\bullet)$, $H^k(B^\bullet)$, and $H^k(C^\bullet)$ in the zig-zag sequence.
    We may reference the above commutative diagram as the commutative grid, and we will omit the indexes of chain maps and boundary operators.

    First, consider $H^k(B^\bullet)$, so we need to show that $\im F^* = \ker G^*$.
    Let $[a] \in H^k(A^\bullet)$ be a cohomology class with $a \in A^k$ being an arbitrary representative of $[a]$.
    Then, $G^*(F^*[a]) = G^*[F(a)] = [G(F(a))] = [0]$ by the exactness of the commutative grid at $B^k$, so $\im F^* \subseteq \ker G^*$.
    Now, we instead start with the cohomology class $[b] \in H^k(B^\bullet)$ such that $G^*[b] = [0]$, meaning that $G(b)$ is a coboundary in $C^k$: $G(b) = \partial^C(c')$ for some $c'  \in C^{k-1}$.
    By exactness, $G \colon B^{k-1} \to C^{k-1}$ is surjective, so there exists $b' \in B^{k-1}$ such that $G(b') = c'$, so $\partial^C(G(b')) = G(b)$.

    Let $\bar{b} \in B^k$ be defined by $\bar{b} = b - \partial^B(b')$.
    We have $[\bar{b}] = [b]$ and
    \begin{equation}
        G(\bar{b}) = G(b) - G(\partial^B(b')) = G(b) - \partial^C(G(b')) = \partial^C(c') - \partial^C(c') = 0.
    \end{equation}
    By exactness, this means that there exists $a \in A^k$ such that $F(a) = \bar{b}$.
    Since $b$ is a cocycle, and $\partial^B \circ \partial^B = 0$, it follows that
    \begin{equation}
        \partial^B(F(a)) = F(\partial^A(a)) = \partial^B(\bar{b}) = \partial^B(b) - \partial^B(\partial^B(b')) = 0.
    \end{equation}
    With this, and the fact that $F$ is injective, we have $\partial^A(a) = 0$, i.e. $a$ is a cocycle, and $[a]$ is well-defined.
    Hence, $[b] = [\bar{b}] = [F(a)] = F^*[a]$, so $\ker G^* \subseteq \im F^*$.
    Therefore, $\im F^* = \ker G^*$.

    To show the exactness of the zig-zag sequence at $H^k(A^\bullet)$ and $H^k(C^\bullet)$, we need to actually define the connecting homomorphisms.
    Let $c \in C^k$ be a cocycle.
    Since $G \colon B^k \to C^k$ is surjective, there exists $b \in B^k$ such that $c = G(b)$,
    Also,
    \begin{equation}
        G(\partial^B(b)) = \partial^C(G(b)) = \partial^C(c) = 0
    \end{equation}
    so $\partial^B(b) \in \ker G = \im F$.
    Thus, there exists an element $a \in A^{k+1}$ such that $F(a) = \partial^B(b)$.
    By the injectivity of $F$, this choice of element is unique.
    Since
    \begin{equation}
        F(\partial^A(a)) = \partial^B(F(a)) = \partial^B(\partial^B(b)) = 0,
    \end{equation}
    and $F$ is injective, we have $\partial^A(a) = 0$, i.e. $a$ is a cocycle.
    We now define the connecting homomorphism $\delta^k \colon H^k(C^\bullet) \to H^{k+1}(A^\bullet)$ by
    \begin{equation}
        \delta^k[c] \coloneq [a].
    \end{equation}
    This definition involves a choice of elements $b$ and $c$, so we will show that the cohomology class $[a]$ is independent of these two choices.
    Suppose we instead choose $\bar{c} = c + \partial^C(c')$, where $c' \in C^{k-1}$, to represent $[c]$.
    By the surjectivity of $G$, there exists an element $b' \in B^{k-1}$ such that $c' = G(b')$.
    We define $\bar{b} \in B^k$ by $\bar{b} = b + \partial^B(b')$, giving us
    \begin{equation}
        G(\bar{b}) = G(b + \partial^B(b')) = G(b) + \partial^C(G(b')) = G(b) + \partial^C(c') = \bar{c}.
    \end{equation}
    By the reasoning we used previously, there exists a unique cocycle $\bar{a} \in A^{k+1}$ such that $F(\bar{a}) = \partial^B(\bar{b})$.
    This means that $F(\bar{a}) = \partial^B(b + \partial^B(b')) = \partial^B(b) = F(a)$.
    As $F$ is injective, we have $a = \bar{a}$, so $[a] = [\bar{a}]$.
    Thus, the cohomology class $[a]$ is independent of our choice of representative of $[c]$.

    Now, we wish to show that it is independent of our choice of element $b$.
    Suppose that we choose a different element $\tilde{b} \in B^k$ such that $c = G(\tilde{b})$.
    This means we have $G(\tilde{b} - b) = G(\tilde{b}) - G(b) = c - c = 0$, so $\tilde{b} - b \in \ker G = \im F$.
    There thus exists an element $\tilde{a} \in A^k$ such that $F(\tilde{a}) = \tilde{b} - b$, so $\tilde{b} = b + F(\tilde{a})$.
    \begin{equation}
        \partial^B(\tilde{b}) = \partial^B(b + F(\tilde{a})) = \partial^B(b) + \partial^B(F(\tilde{a})) = \partial^B(b) + F(\partial^A(\tilde{a}))
    \end{equation}
    Again, using established reasoning, there is a unique cocycle $\hat{a} \in A^{k+1}$ such that $F(\hat{a}) = \partial^B(\tilde{b})$, so we have
    \begin{equation}
        F(\hat{a}) = \partial^B(b) + F(\partial^A(\tilde{a})) = F(a) + F(\partial^A(\tilde{a})) = F(a + \partial^A(\tilde{a})).
    \end{equation}
    Since $F$ is injective, we get $\hat{a} = a + \partial^A(\tilde{a})$.
    Then, $[\hat{a}] = [a]$.
    Hence, $[a]$ is independent of our choice of the element $b$, and the connecting homomorphism is therefore well-defined.

    The connecting homomorphisms are characterized by the following zig-zag diagram of mappings between elements.
    \[\begin{tikzcd}[sep=large]
	a & {\partial^B(b)} & \\
	& b & c
	\arrow["F", tail, from=1-1, to=1-2]
	\arrow["{\partial^B}"', maps to, from=2-2, to=1-2]
	\arrow["G", maps to, two heads, from=2-2, to=2-3]
	\arrow["{\delta^k}", curve={height=-50pt}, maps to, from=2-3, to=1-1]
    \end{tikzcd}\]
    We now wish to show that the zig-zag sequence is exact at $H^k(C^\bullet)$, so we need to show that $\im G^* = \ker \delta^k$.
    Let $[b] \in H^k(B^\bullet)$ be a cohomology class, with $b \in B^k$ is an arbitrary representative of $[b]$.
    Then, $\delta^k[G^*[b]] = \delta^k[G(b)]$.
    Since $b$ is a cocycle, $\partial^B(b) = 0$.
    Following the zig-zag diagram above, $F(a) = 0$, where $a \in A^{k+1}$ is an arbitrary representative of the image of $G^*[b]$.
    By the injectivity of $F$, $a = 0$.
    Hence, $\delta^k[G^*[b]] = [0]$, so $\im G^* \subseteq \ker \delta^k$.

    For the converse, let $[c] \in H^k(C^\bullet)$ be a cohomology class such that $\delta^k[c] = [0]$.
    Then, an arbitrary representative $a \in A^{k+1}$ of $[a] \in H^{k+1}(A^\bullet)$ is a coboundary: $a = \partial^A(a')$ for some $a' \in A^k$.
    Also, there is some $b \in B^k$ such that $F(a) = \partial^B(b)$, so $\partial^B(b) = F(\partial^A(a')) = \partial^B(F(a'))$.
    Thus, $\partial^B(b - F(a')) = 0$, so $b - F(a')$ is a cocycle in $B^k$.
    Since, by the exactness of the commutative grid at $B^k$,
    \begin{equation}
        G(b - F(a')) = G(b) - G(F(a')) = c - 0 = c,
    \end{equation}
    the class $[c]$ is the image of $[b - F(a')]$.
    Therefore, $\ker \delta^k \subseteq \im G^*$, meaning that $\im G^* = \ker \delta^k$.

    Finally, we need to show that the zig-zag sequence is exact at $H^k(A^\bullet)$; we have to show that $\im \delta^{k-1} = \ker F^*$.
    Let $[c] \in H^{k-1}(C^\bullet)$ be a cohomology class.
    Then, $\delta^{k-1}[c] = [a]$, where $F(a) = \partial^B(b)$ for some $b \in B^{k-1}$ by construction.
    Hence, $F^*[a] = [F(a)] = [\partial^B(b)] = [0]$, so $\im \delta^{k-1} \subseteq \ker F^*$.
    Now, suppose that $[a] \in H^k(A^\bullet)$ is a cohomology class such that $F^*[a] = [0]$.
    This means that $[F(a)] = [0]$, so $F(a) = \partial^B(b)$ for some $b \in B^{k-1}$.
    Let $c = G(b)$, so 
    \begin{equation}
        \partial^C(c) = \partial^C(G(b)) = G(\partial^B(b)) = G(F(a)) = 0,
    \end{equation}
    meaning that $c$ is a cocycle.
    By construction, we have $\delta^{k-1}[c] = [a]$, so $\ker F^* \subseteq \im \delta^{k-1}$.
    Therefore, $\im \delta^{k-1} = \ker F^*$, completing our proof.
\end{proof}
A similar result to the zig-zag lemma is the snake lemma.
Proving it involves diagram chasing to show exactness at the various parts of a sequence.
However, we are now able to use the zig-zag lemma to avoid a fair portion of this tedious process.
\begin{lemma}[Snake Lemma]\label{lem:snake_lemma}
    Consider the following commutative diagram of $R$-modules over a commutative ring $R$ with $R$-module homomorphisms such that the rows are exact.
    \[\begin{tikzcd}[sep=large]
	& A & B & C & 0 \\
	0 & {A'} & {B'} & {C'}
	\arrow["f", from=1-2, to=1-3]
	\arrow["\alpha", from=1-2, to=2-2]
	\arrow["g", from=1-3, to=1-4]
	\arrow["\beta", from=1-3, to=2-3]
	\arrow[from=1-4, to=1-5]
	\arrow["\gamma", from=1-4, to=2-4]
	\arrow[from=2-1, to=2-2]
	\arrow["{f'}", from=2-2, to=2-3]
	\arrow["{g'}", from=2-3, to=2-4]
    \end{tikzcd}\]
    Then, there exists an $R$-module homomorphism $\delta \colon \ker \gamma \to \coker \alpha$, called the connecting homomorphism, making the following sequence exact.
    \[\begin{tikzcd}
        & 0 \ar[r] & \ker f \ar[r] & \ker \alpha \ar[r] & \ker \beta \ar[r] & \ker \gamma \arrow[out=0, in=180, looseness=2, overlay, lllld, "\delta"{description}]\\
        & \coker \alpha \ar[r] & \coker \beta \ar[r] & \coker \gamma \ar[r] & \coker g' \ar[r] & 0
    \end{tikzcd}\]
\end{lemma}
\begin{proof}
    We first consider the case where the two rows of the commutative diagram are short exact and not just exact.
    That is, the diagram is the following.
    \[\begin{tikzcd}
	0 & A & B & C & 0 \\
	0 & {A'} & {B'} & {C'} & 0
	\arrow[from=1-1, to=1-2]
	\arrow["f", from=1-2, to=1-3]
	\arrow["\alpha", from=1-2, to=2-2]
	\arrow["g", from=1-3, to=1-4]
	\arrow["\beta", from=1-3, to=2-3]
	\arrow[from=1-4, to=1-5]
	\arrow["\gamma", from=1-4, to=2-4]
	\arrow[from=2-1, to=2-2]
	\arrow["{{f'}}", from=2-2, to=2-3]
	\arrow["{{g'}}", from=2-3, to=2-4]
	\arrow[from=2-4, to=2-5]
    \end{tikzcd}\]
    We define the following three cochain complexes $A^\bullet$, $B^\bullet$, and $C^\bullet$.
    \[\begin{tikzcd}
	0 & A & {A'} & 0 \\
	0 & B & {B'} & 0 \\
	0 & C & {C'} & 0
	\arrow[from=1-1, to=1-2]
	\arrow["\alpha", from=1-2, to=1-3]
	\arrow[from=1-3, to=1-4]
	\arrow[from=2-1, to=2-2]
	\arrow["\beta", from=2-2, to=2-3]
	\arrow[from=2-3, to=2-4]
	\arrow[from=3-1, to=3-2]
	\arrow["\gamma", from=3-2, to=3-3]
	\arrow[from=3-3, to=3-4]
    \end{tikzcd}\]
    By the exactness of the rows of the commutative diagram, we can form a short exact sequence of these cochain complexes.
    \[\begin{tikzcd}
	0 & {A^\bullet} & {B^\bullet} & {C^\bullet} & 0
	\arrow[from=1-1, to=1-2]
	\arrow["F", from=1-2, to=1-3]
	\arrow["G", from=1-3, to=1-4]
	\arrow[from=1-4, to=1-5]
    \end{tikzcd}\]
    Here, commutativity of the original diagram means that the map $F$, representing the sequence containing $f$ and $f'$, is indeed a cochain map -- and likewise for $G$.
    Note that our complexes only have zeroth and first nontrivial cohomology groups.
    We apply the \hyperref[lem:zig-zag_lemma]{zig-zag lemma} to our sequence of cochain complexes to obtain our desired long exact sequence.
    This is because $H^0(A^\bullet) = \ker \alpha$ and $H^1(A^\bullet) = \ker 0_{A'} / \im \alpha = A' / \im \alpha \eqcolon \coker \alpha$.
    Similar results hold for $B^\bullet$ and $C^\bullet$, giving us the following long exact sequence.
    \[\begin{tikzcd}
        & 0 \ar[r] & \ker \alpha \ar[r] & \ker \beta \ar[r] & \ker \gamma \arrow[out=0, in=180, looseness=2, overlay, llld, "\delta"{description}]\\
        & \coker \alpha \ar[r] & \coker \beta \ar[r] & \coker \gamma \ar[r] & 0
    \end{tikzcd}\]
    The map $0^{\ker \alpha}$ from zero and the map $0_{\coker \gamma}$ to zero are used to show the exactness of the sequence at $\ker \alpha$ and $\coker \gamma$.
    Note that the connecting homomorphism of the snake lemma is the same as that of the zig-zag lemma.
    This completes the proof of the snake lemma in our simplified case.
    The commutative diagram and the long exact sequence can be combined into a single diagram, yielding an explanation for the name of the snake lemma.
    \[\begin{tikzcd}[row sep=large]
	0 & {\ker \alpha} & {\ker \beta} & {\ker \gamma} & \\
	0 & A & B & C & 0 \\
	0 & {A'} & {B'} & {C'} & 0 \\
	& {\coker \alpha} & {\coker \beta} & {\coker  \gamma} & 0
	\arrow[from=1-1, to=1-2, red]
	\arrow[from=1-2, to=1-3, red]
	\arrow[from=1-2, to=2-2]
	\arrow[from=1-3, to=1-4, red]
	\arrow[from=1-3, to=2-3]
	\arrow[from=1-4, to=2-4]
    \arrow[from=1-4, to=4-2, overlay, red, rounded corners, to path={-- ++(3.0,0) -- ++(0,-2.4) node[midway,right]{\(\delta\)} -- ++(-11,0) -- ++(0,-2.85) -- (\tikztotarget)}]
	\arrow[from=2-1, to=2-2, blue]
    \arrow["f", from=2-2, to=2-3, blue]
	\arrow["\alpha", from=2-2, to=3-2, blue]
	\arrow["g", from=2-3, to=2-4, blue]
	\arrow["\beta", from=2-3, to=3-3, blue]
	\arrow[from=2-4, to=2-5, blue]
	\arrow["\gamma", from=2-4, to=3-4, blue]
	\arrow[from=3-1, to=3-2, blue]
	\arrow["{{f'}}", from=3-2, to=3-3, blue]
	\arrow[from=3-2, to=4-2]
	\arrow["{{g'}}", from=3-3, to=3-4, blue]
	\arrow[from=3-3, to=4-3]
    \arrow[from=3-4, to=3-5, blue]
	\arrow[from=3-4, to=4-4]
	\arrow[from=4-2, to=4-3, red]
	\arrow[from=4-3, to=4-4, red]
	\arrow[from=4-4, to=4-5, red]
    \end{tikzcd}\]
    Now, we will consider the general case, where our combined diagram is instead the following.
    \[\begin{tikzcd}[row sep=large]
	0 & {\ker f} & {\ker \alpha} & {\ker \beta} & {\ker \gamma} && \\
	&& A & B & C & 0 \\
	& 0 & {A'} & {B'} & {C'} \\
	&& {\coker \alpha} & {\coker \beta} & {\coker  \gamma} & {\coker g'} & 0
	\arrow[from=1-1, to=1-2, red]
	\arrow[from=1-2, to=1-3, red]
	\arrow[from=1-3, to=1-4, red]
	\arrow[from=1-3, to=2-3]
	\arrow[from=1-4, to=1-5, red]
	\arrow[from=1-4, to=2-4]
	\arrow[from=1-5, to=2-5]
    \arrow[from=1-5, to=4-3, overlay, red, rounded corners, to path={-- ++(3.5,0) -- ++(0,-2.4) node[midway,right]{\(\delta\)} -- ++(-12,0) -- ++(0,-2.85) -- (\tikztotarget)}]
	\arrow["f", from=2-3, to=2-4, blue]
	\arrow["\alpha", from=2-3, to=3-3, blue]
	\arrow["g", from=2-4, to=2-5, blue]
	\arrow["\beta", from=2-4, to=3-4, blue]
	\arrow[from=2-5, to=2-6, blue]
	\arrow["\gamma", from=2-5, to=3-5, blue]
	\arrow[from=3-2, to=3-3, blue]
	\arrow["{f'}", from=3-3, to=3-4, blue]
	\arrow[from=3-3, to=4-3]
	\arrow["{g'}", from=3-4, to=3-5, blue]
	\arrow[from=3-4, to=4-4]
	\arrow[from=3-5, to=4-5]
	\arrow[from=4-3, to=4-4, red]
	\arrow[from=4-4, to=4-5, red]
	\arrow[from=4-5, to=4-6, red]
	\arrow[from=4-6, to=4-7, red]
    \end{tikzcd}\]
    Instead of our prior cochain complexes, we define $\bar{A} = A / \ker f$ and $\overbar{C}' = \im g'$.
    The map $A \to B$ factors into $A \surject \bar{A} \inject B$, and the map $B' \to C'$ factors into $B' \surject \overbar{C}' \inject C'$.
    The injectivity of $\bar{A} \inject B$ and the surjectivity of $B' \surject \overbar{C}'$ allow for us to create a short exact sequence -- in order to apply the zig-zag lemma.
    Now, we consider the following cochain complexes $\bar{A}^\bullet$, $B^\bullet$, and $\overbar{C}^\bullet$.
    \[\begin{tikzcd}
	0 & \bar{A} & {A'} & 0 \\
	0 & B & {B'} & 0 \\
	0 & C & {\overbar{C}'} & 0
	\arrow[from=1-1, to=1-2]
	\arrow["\bar{\alpha}", from=1-2, to=1-3]
	\arrow[from=1-3, to=1-4]
	\arrow[from=2-1, to=2-2]
	\arrow["\beta", from=2-2, to=2-3]
	\arrow[from=2-3, to=2-4]
	\arrow[from=3-1, to=3-2]
	\arrow["\bar{\gamma}", from=3-2, to=3-3]
	\arrow[from=3-3, to=3-4]
    \end{tikzcd}\]
    The map $\bar{\alpha} \colon \bar{A} \to A'$, defined by $\bar{\alpha}[a]_f = \alpha(a)$, is well-defined by the following reasoning.
    Note the commutativity of the left square of the commutative diagram.
    Given an element $x \in \ker f \subseteq A$, we have $f'(\alpha(x)) = \beta(f(x))$.
    Since $f(x) = 0$, $\beta(f(x)) = 0$.
    The injectivity of $f'$, which may be deduced from the diagram, means that $\alpha(x) = 0$.
    Thus, $\ker f \subseteq \ker \alpha$, so $f(a) = f(\tilde{a}) \implies \alpha(a) = \alpha(\tilde{a})$.
    Hence, $\bar{\alpha}[a]_f$ does not depend on the representative of $[a]_f$.

    Now, consider the map $\bar{\gamma} \colon C \to \overbar{C}'$ given by $\bar{\gamma}(c) = \gamma(c)$.
    This is well-defined by the commutativity of the right square of the commutative diagram: $g' \circ \beta = \gamma \circ g$.
    From the diagram, it can be deduced that $g$ is surjective, so any element in $\gamma(c) \in \im \gamma$ can be written as $\gamma(g(b))$ for some element $b \in B$.
    This equals $g'(\beta(b))$, so $\im \gamma \subseteq \im g' = \overbar{C}'$.
    As $\bar{\alpha}$ and $\bar{\gamma}$ are well-defined maps, we can now consider the following sequence of cochain complexes; it is short exact by construction.
    \[\begin{tikzcd}
	0 & {\bar{A}^\bullet} & {B^\bullet} & {\overbar{C}^\bullet} & 0
	\arrow[from=1-1, to=1-2]
	\arrow["F", from=1-2, to=1-3]
	\arrow["G", from=1-3, to=1-4]
	\arrow[from=1-4, to=1-5]
    \end{tikzcd}\]
    Applying the \hyperref[lem:zig-zag_lemma]{zig-zag lemma} to this short exact sequence yields the following long exact sequence; we may call this the augmented snake sequence.
    The only different computations here are $\ker \bar{\alpha} = \ker \alpha / \ker f$, $\ker \bar{\gamma} = \ker \gamma$, $\coker \bar{\alpha} = A' / \im \bar{\alpha} = A' / \im \alpha = \coker \alpha$, and $\coker \bar{\gamma} = \im g' / \im \gamma$.
    \[\begin{tikzcd}
        & 0 \ar[r] & \frac{\ker \alpha}{\ker f} \ar[r] & \ker \beta \ar[r] & \ker \gamma \arrow[out=0, in=180, looseness=2, overlay, llld, "\delta"{description}]\\
        & \coker \alpha \ar[r] & \coker \beta \ar[r] & \frac{\im g'}{\im \gamma} \ar[r] & 0
    \end{tikzcd}\]
    The first isomorphism theorem applied to the projection map $\pi \colon \ker \alpha \to \ker \alpha / \ker f$ gives us the following short exact sequence involving the inclusion and projection maps. 
    \[\begin{tikzcd}
	0 & {\ker f} & {\ker \alpha} & {\frac{\ker \alpha}{\ker f}} & 0
	\arrow[from=1-1, to=1-2]
	\arrow["\iota", hook, from=1-2, to=1-3]
	\arrow["\pi", two heads, from=1-3, to=1-4]
	\arrow[from=1-4, to=1-5]
    \end{tikzcd}\]
    We are able to splice this with the augmented snake sequence to obtain the following diagram.
    \[\begin{tikzcd}[row sep=small]
	0 & {\ker f} & {\ker \alpha} && {\ker \beta} & \cdots \\
	&&& {\frac{\ker \alpha}{\ker f}} \\
	&& 0 && 0
	\arrow[from=1-1, to=1-2]
	\arrow["\iota", hook, from=1-2, to=1-3]
	\arrow["{\bar{f} \circ \pi}", dashed, from=1-3, to=1-5]
	\arrow["\pi"', two heads, from=1-3, to=2-4]
	\arrow[from=1-5, to=1-6]
	\arrow["{\bar{f}}"', from=2-4, to=1-5]
	\arrow[from=2-4, to=3-5]
	\arrow[from=3-3, to=2-4]
    \end{tikzcd}\]
    Since the map $\bar{f}$ is injective, $\ker (\bar{f} \circ \pi) = \ker \pi$.
    As the projection map $\pi$ is surjective, $\im (\bar{f} \circ \pi) = \im \bar{f}$.
    Hence, our new sequence -- the top row of the above diagram -- is exact $\ker \alpha$ and $\ker \beta$.
    
    We apply a similar approach to the other end of our sequence.
    Applying the third isomorphism theorem gives us a short exact sequence,
    \[\begin{tikzcd}
	0 & {\frac{\im g'}{\im \gamma}} & {\frac{C'}{\im \gamma}} & {\frac{C'}{\im g'}} & 0,
	\arrow[from=1-1, to=1-2]
	\arrow["\kappa", hook, from=1-2, to=1-3]
	\arrow["\tau", two heads, from=1-3, to=1-4]
	\arrow[from=1-4, to=1-5]
    \end{tikzcd}\]
    which can be rewritten in terms of cokernels.
    \[\begin{tikzcd}
	0 & {\frac{\im g'}{\im \gamma}} & {\coker \gamma} & {\coker g'} & 0,
	\arrow[from=1-1, to=1-2]
	\arrow["\kappa", hook, from=1-2, to=1-3]
	\arrow["\tau", two heads, from=1-3, to=1-4]
	\arrow[from=1-4, to=1-5]
    \end{tikzcd}\]
    We splice this with the augmented snake sequence to obtain the following diagram.
    \[\begin{tikzcd}[row sep=small]
	\cdots & {\coker \beta} && {\coker \gamma} & {\coker g'} & 0 \\
	&& {\frac{\im g'}{\im \gamma}} \\
	& 0 && 0
	\arrow[from=1-1, to=1-2]
	\arrow["{\kappa \circ \bar{g}'}", dashed, from=1-2, to=1-4]
	\arrow["{\bar{g}'}"', from=1-2, to=2-3]
	\arrow["\tau", two heads, from=1-4, to=1-5]
	\arrow[from=1-5, to=1-6]
	\arrow["\kappa"', hook, from=2-3, to=1-4]
	\arrow[from=2-3, to=3-4]
	\arrow[from=3-2, to=2-3]
    \end{tikzcd}\]
    Since the inclusion map $\kappa$ is injective, $\ker (\kappa \circ \bar{g}') = \ker \bar{g}'$.
    As the map $\bar{g}'$ is surjective, $\im (\kappa \circ \bar{g}') = \im \kappa$.
    Hence, our new sequence -- the top row of the above diagram -- is exact at $\coker \gamma$ and $\coker g'$.
    Applying both of our splicings to the augmented snake sequence yields the desired long exact sequence, the full snake sequence.
\end{proof}

%! TODO MAYER--VIETORIS SEQUENCE

\begin{theorem}[Mayer--Vietoris Sequence]\label{thm:mayer--vietoris_sequence}
    
\end{theorem}
\begin{proof}
    
\end{proof}
\[
\begin{tikzcd}
    & \cdots \ar[r, "j^*"] & H_{\mathrm{dR}}^{k-1}(U \cap V) \ar[out=0, in=180, looseness=2, overlay, lld, "\d^*"{description}]\\
    H^k_{\mathrm{dR}}(M) \ar[r, "i^*"] & H^k_{\mathrm{dR}}(U) \oplus H^k_{\mathrm{dR}}(V) \ar[r, "j^*"] & H^k_{\mathrm{dR}}(U \cap V) \ar[out=0, in=180, looseness=2, overlay, lld, "\d^*"{description}]\\
    H_{\mathrm{dR}}^{k+1}(M) \ar[r, "i^*"] & \cdots &
\end{tikzcd}
\]

\subsubsection{Homotopy Invariance}\label{subsubsec:de_rham_theory:de_rham_cohomology:homotopy_invariance}

\subsection{\v{C}ech-De Rham Cohomology}

\printbibliography

\begin{appendices}
\crefalias{section}{appsec}
\crefalias{subsection}{appsec}

\section{Topological Spaces}\label{sec:topological_spaces}

This appendix fixes the conventions of point-set topology that are used in the text.
The standard theory is not reproduced here, and it may be found in \cite{Lee2011}.

Topological spaces allow for the concept of closeness to be defined without necessarily having any notion of distance.
They are fundamental for the definition of manifolds, and thus for differential forms.
The typical definition involves the use of open sets.
\begin{definition}[Topological Space]\label{def:topological_space}
    A topological space is a set $X$ equipped with a collection $\mathcal{T}$ of subsets of $X$, called \emph{open sets}, which satisfy the following axioms.
    \begin{enumerate}
        \item The empty set $\varnothing$ and $X$ are elements of $\mathcal{T}$.
        \item Any union of sets in $\mathcal{T}$ is an element of $\mathcal{T}$.
        \item Any finite intersection of sets in $\mathcal{T}$ is an element of $\mathcal{T}$.
    \end{enumerate}
    The space may be denoted by $(X, \mathcal{T})$, and the set $\mathcal{T}$ is called a \emph{topology} on $X$.
\end{definition}
\begin{remark}
    For convenience, we may refer to the topological space as $X$, and we may do this when dealing with any other space that is a set with additional structure.
\end{remark}
\begin{definition}[Neighbourhood]\label{def:neighbourhoods_via_open_sets}
    Given a topological space $(X, \mathcal{T})$, a neighbourhood of $x \in X$ is a set $N \subseteq X$ that includes an open set $U \in \mathcal{T}$ such that $x \in U$.
    If $N$ is an open set, then it is an open neighbourhood.
\end{definition}
\begin{definition}[Continuity]\label{def:continuity}
    A function $f \colon X \to Y$ between topological spaces is continuous if $f^{-1}(V) \coloneq \{x \in X : f(x) \in V\}$ is an open set in $X$ for every open set $V$ in $Y$.
\end{definition}
\begin{definition}[Homeomorphism]\label{def:homeomorphism}
    A function $f \colon X \to Y$ between topological spaces is a \emph{homeomorphism} if it is a continuous bijection with a continuous inverse.
    If such a function exists, then the spaces $X$ and $Y$ are said to be \emph{homeomorphic}, denoted by $X \cong Y$.
\end{definition}
\begin{definition}[Hausdorff]\label{def:hausdorff}
    A topological space $X$ is \emph{Hausdorff} if for any distinct points $x, y \in X$, there exist open neighbourhoods $U$ of $x$ and $V$ of $y$ in $X$ that are disjoint, i.e. $U \cap V = \varnothing$.
\end{definition}
\begin{definition}[Second Countable]\label{def:second_countable}
    A topological space $(X, \mathcal{T})$ is \emph{second countable} if $\mathcal{T}$ has a countable basis, i.e. a countable collection $\mathcal{B} \subseteq \mathcal{T}$ such that every open set is a union of elements of $\mathcal{B}$.
\end{definition}
\begin{definition}[Cover]\label{def:cover}
    A \emph{cover} of a topological space $X$ is a collection of subsets of $X$ whose union equals $X$.
    If every subset is open, it is called an \emph{open cover}.
    If there are a finite number of subsets in the cover, it is called a \emph{finite cover}.
    A \emph{subcover} of a cover is a subcollection that is still a cover of $X$.
\end{definition}
\begin{definition}[Locally Finite Collection]\label{def:locally_finite_cover}
    A collection $\{A_i \subseteq X\}_{i \in I}$ of subsets of a topological space $X$, such as a cover, is \emph{locally finite} if for every point $x \in X$, there exists a neighbourhood $U_x$ of $x$ that intersects only finitely many elements of the collection, i.e. $U_x \cap A_i \neq \varnothing$ for only a finite number of $i \in I$.
\end{definition}
\begin{definition}[Refinement of Open Covers]\label{def:refinement_of_open_covers}
    A \emph{refinement} of an open cover $\{U_i \subseteq X\}_{i \in I}$ is an open cover $\{V_j \subseteq X\}_{j \in J}$ such that for each $j \in J$, there exists an $i \in I$ with $V_j \subseteq U_i$.
\end{definition}
\begin{definition}[Compact]\label{def:compact}
    A topological space $X$ is \emph{compact} if every open cover has a finite subcover.
\end{definition}
\begin{definition}[Paracompact]\label{def:paracompact}
    A topological space $X$ is \emph{paracompact} if every open cover of $X$ has a refinement by a locally finite open cover.
\end{definition}
\begin{definition}[Connected]\label{def:connected}
    A topological space $X$ is \emph{connected} if it is not the union of two disjoint nonempty open subsets.
    The connected components of $X$ are its maximal connected subsets.
\end{definition}
\begin{definition}[Partition of Unity]\label{def:partition_of_unity}
    Let $M$ be a smooth manifold, and let $\{U_i\}_{i \in I}$ be an open cover of $M$.
    A smooth partition of unity subordinate to $\{U_i\}_{i \in I}$ is a collection of smooth functions $\psi_i \colon M \to [0, 1]$ such that $\supp \psi_i \subseteq U_i$ for each $i$, the collection $\{\supp \psi_i\}_{i \in I}$ is locally finite, and $\sum_{i \in I} \psi_i = 1$.
\end{definition}
The following results are used in the text without proof.
\begin{itemize}
    \item Every second-countable, Hausdorff, locally Euclidean space is paracompact \cite[Thm.~1.15]{Lee2012}.
    \item Every open cover of a smooth manifold admits a subordinate smooth partition of unity \cite[Thm.~2.23]{Lee2012}.
    \item The interior and boundary of a manifold with boundary are disjoint, which rests on invariance of domain \cite[Thm.~1.37]{Lee2012}.
\end{itemize}

\end{appendices}

\end{document}